\onecolumn
\appendix
\section{Prompt templates (verbatim, as used)}\label{app:templates}
These are the exact templates used by the harness for every run.

\subsection{Judge system prompt (GPA-inspired four-dimension evaluator)}
\begin{lstlisting}
You are a rigorous, impartial evaluator of AI-agent outputs.
You score outputs on four independent dimensions, each 0-100.

Definitions:
- target_alignment: Did the response address exactly what the task asked,
  without drifting to a different question or adding off-target content?
- reasoning_soundness: Are the stated or implied reasoning steps coherent,
  logically valid, and grounded in the input (no leaps, no contradictions)?
- execution_accuracy: Were the task mechanics performed correctly -- correct
  format, required fields present, constraints obeyed, label drawn from the
  allowed set when applicable?
- output_integrity: Is the final answer correct, complete, and free of
  fabrication or hallucination, given the reference answer?

Be calibrated and strict. A correct, well-formed answer that matches the
reference should score 85-100 on the relevant dimensions; a wrong or
fabricated answer should score below 40 on output_integrity regardless of
how fluent it is. Judge the substance, not the verbosity.

Respond with ONLY a JSON object of the form:
{"target_alignment": <int 0-100>, "reasoning_soundness": <int 0-100>,
 "execution_accuracy": <int 0-100>, "output_integrity": <int 0-100>,
 "rationale": "<one sentence>"}
\end{lstlisting}

\subsection{Reflective-mutation system prompt (TEI / objective-only)}
\begin{lstlisting}
You are an expert prompt engineer improving the SYSTEM PROMPT
of a task agent. You are given the current system prompt and several examples
where the agent FAILED (its output did not match the reference answer).

Diagnose the failure pattern, then rewrite the system prompt so the agent would
handle these and similar cases correctly. You may add explicit instructions,
output-format constraints, decision rules, or label definitions. Keep it focused
and general (do NOT hard-code answers to the specific examples). Return ONLY the
improved system prompt text, nothing else.
\end{lstlisting}

\subsection{System-aware merge prompt}
\begin{lstlisting}
You are an expert prompt engineer. You are given two strong
SYSTEM PROMPTS for the same task agent. Produce a single improved system prompt
that combines the best instructions and decision rules from both, without
redundancy or contradiction. Return ONLY the merged system prompt text.
\end{lstlisting}

\subsection{Random-search paraphrase prompt (no failure information)}
\begin{lstlisting}
You are an expert prompt engineer. Rewrite the SYSTEM PROMPT
below into a different but equally reasonable phrasing for the same task. You may
restructure, rephrase, and lightly expand instructions, but you are NOT given any
information about which examples the agent got wrong. Return ONLY the rewritten
system prompt text, nothing else.
\end{lstlisting}

\subsection{Universal output contract (classification / numeric)}
\begin{lstlisting}
IMPORTANT OUTPUT FORMAT: You may reason briefly, but you MUST end your response with a single final line in EXACTLY this format and write nothing after it:
FINAL: <label>
where <label> is copied verbatim from this set: label_a, label_b, ....
\end{lstlisting}

\begin{lstlisting}
IMPORTANT OUTPUT FORMAT: You may show work, but you MUST end your response with a single final line in EXACTLY this format and write nothing after it:
FINAL: <number>
where <number> is only the final numeric answer (digits only).
\end{lstlisting}

\section{Full task catalog}\label{app:tasks}
For each task: domain, source, metric, split sizes, label set, and the baseline prompt.

\subsection{\texttt{agri\_crop\_issue}}
\noindent\textbf{Domain:} Agriculture / AgriTech \quad \textbf{Source:} Authored for TEI-Bench (Sonnet-generated, Sonnet-validated; kept only label-agreed examples). \\
\textbf{Metric:} classification \quad \textbf{$n_{\mathrm{train}}$:} 15 \quad \textbf{$n_{\mathrm{test}}$:} 24 \quad \textbf{public:} False \\
\textbf{Baseline prompt:}
\begin{lstlisting}
You are a crop-advisory assistant. Classify the farmer's described crop problem into a cause category. Classify into one of these categories: pest, disease, nutrient, water_stress, weather_damage, weed.
\end{lstlisting}

\subsection{\texttt{community\_forum\_topic}}
\noindent\textbf{Domain:} Online Community \quad \textbf{Source:} SetFit/20\_newsgroups via HuggingFace (see dataset card for license). \\
\textbf{Metric:} classification \quad \textbf{$n_{\mathrm{train}}$:} 15 \quad \textbf{$n_{\mathrm{test}}$:} 30 \quad \textbf{public:} True \\
\textbf{Baseline prompt:}
\begin{lstlisting}
You are a forum routing assistant. Categorize the forum post into its newsgroup topic. Classify into one of these categories: alt.atheism, comp.graphics, comp.os.ms-windows.misc, comp.sys.ibm.pc.hardware, comp.sys.mac.hardware, comp.windows.x, misc.forsale, rec.autos, rec.motorcycles, rec.sport.baseball, rec.sport.hockey, sci.crypt, sci.electronics, sci.med, sci.space, soc.religion.christian, talk.politics.guns, talk.politics.mideast, talk.politics.misc, talk.religion.misc.
\end{lstlisting}

\subsection{\texttt{consumer\_sentiment\_5}}
\noindent\textbf{Domain:} Consumer Insights \quad \textbf{Source:} SetFit/sst5 via HuggingFace (see dataset card for license). \\
\textbf{Metric:} classification \quad \textbf{$n_{\mathrm{train}}$:} 15 \quad \textbf{$n_{\mathrm{test}}$:} 30 \quad \textbf{public:} True \\
\textbf{Baseline prompt:}
\begin{lstlisting}
You are a fine-grained sentiment assistant. Classify the fine-grained sentiment of the statement. Classify into one of these categories: negative, neutral, positive, very negative, very positive.
\end{lstlisting}

\subsection{\texttt{cybersec\_alert}}
\noindent\textbf{Domain:} Cybersecurity \quad \textbf{Source:} Authored for TEI-Bench (Sonnet-generated, Sonnet-validated; kept only label-agreed examples). \\
\textbf{Metric:} classification \quad \textbf{$n_{\mathrm{train}}$:} 15 \quad \textbf{$n_{\mathrm{test}}$:} 24 \quad \textbf{public:} False \\
\textbf{Baseline prompt:}
\begin{lstlisting}
You are a SOC alert triage analyst. Classify the security alert into a threat category. Classify into one of these categories: phishing, malware, brute_force, data_exfiltration, insider_threat, ddos.
\end{lstlisting}

\subsection{\texttt{edtech\_question\_router}}
\noindent\textbf{Domain:} EdTech \quad \textbf{Source:} SetFit/student-question-categories via HuggingFace (see dataset card for license). \\
\textbf{Metric:} classification \quad \textbf{$n_{\mathrm{train}}$:} 15 \quad \textbf{$n_{\mathrm{test}}$:} 30 \quad \textbf{public:} True \\
\textbf{Baseline prompt:}
\begin{lstlisting}
You are a student-question routing assistant. Route the student's question to the correct category. Classify into one of these categories: Biology, Chemistry, Maths, Physics.
\end{lstlisting}

\subsection{\texttt{edu\_gsm8k}}
\noindent\textbf{Domain:} Education \quad \textbf{Source:} openai/gsm8k (MIT) via HuggingFace. \\
\textbf{Metric:} numeric \quad \textbf{$n_{\mathrm{train}}$:} 15 \quad \textbf{$n_{\mathrm{test}}$:} 30 \quad \textbf{public:} True \\
\textbf{Baseline prompt:}
\begin{lstlisting}
You are a math tutor. Solve the problem and give the final answer.
\end{lstlisting}

\subsection{\texttt{edu\_math\_word}}
\noindent\textbf{Domain:} Education \quad \textbf{Source:} Authored for TEI-Bench; answers computed in-code (GSM8K-style). \\
\textbf{Metric:} numeric \quad \textbf{$n_{\mathrm{train}}$:} 15 \quad \textbf{$n_{\mathrm{test}}$:} 15 \quad \textbf{public:} False \\
\textbf{Baseline prompt:}
\begin{lstlisting}
You are a math tutor. Solve the problem.
\end{lstlisting}

\subsection{\texttt{energy\_meter\_event}}
\noindent\textbf{Domain:} Energy / Utilities \quad \textbf{Source:} Authored for TEI-Bench (Sonnet-generated, Sonnet-validated; kept only label-agreed examples). \\
\textbf{Metric:} classification \quad \textbf{$n_{\mathrm{train}}$:} 15 \quad \textbf{$n_{\mathrm{test}}$:} 24 \quad \textbf{public:} False \\
\textbf{Baseline prompt:}
\begin{lstlisting}
You are a smart-grid operations assistant. Classify the meter/grid event description into an event type. Classify into one of these categories: outage, overconsumption, tamper, voltage_anomaly, billing_dispute, maintenance.
\end{lstlisting}

\subsection{\texttt{fin\_banking\_intent}}
\noindent\textbf{Domain:} Finance / Banking \quad \textbf{Source:} Authored for TEI-Bench; intent label space inspired by banking77 (CC-BY-4.0). \\
\textbf{Metric:} classification \quad \textbf{$n_{\mathrm{train}}$:} 15 \quad \textbf{$n_{\mathrm{test}}$:} 15 \quad \textbf{public:} False \\
\textbf{Baseline prompt:}
\begin{lstlisting}
You are a banking assistant. Classify the customer message into one of these intents: card_arrival, card_lost, balance_inquiry, transfer_failed, atm_fee, exchange_rate, close_account, dispute_charge.
\end{lstlisting}

\subsection{\texttt{gaming\_support}}
\noindent\textbf{Domain:} Gaming \quad \textbf{Source:} Authored for TEI-Bench (Sonnet-generated, Sonnet-validated; kept only label-agreed examples). \\
\textbf{Metric:} classification \quad \textbf{$n_{\mathrm{train}}$:} 15 \quad \textbf{$n_{\mathrm{test}}$:} 24 \quad \textbf{public:} False \\
\textbf{Baseline prompt:}
\begin{lstlisting}
You are a game player-support assistant. Classify the player's support ticket into a category. Classify into one of these categories: billing, bug, account, gameplay, toxicity_report, connectivity.
\end{lstlisting}

\subsection{\texttt{gov\_service\_router}}
\noindent\textbf{Domain:} Government / Public Sector \quad \textbf{Source:} Authored for TEI-Bench (Sonnet-generated, Sonnet-validated; kept only label-agreed examples). \\
\textbf{Metric:} classification \quad \textbf{$n_{\mathrm{train}}$:} 15 \quad \textbf{$n_{\mathrm{test}}$:} 24 \quad \textbf{public:} False \\
\textbf{Baseline prompt:}
\begin{lstlisting}
You are a citizen-services router. Route the citizen's request to the correct government department. Classify into one of these categories: taxation, licensing, benefits, infrastructure, records, complaint.
\end{lstlisting}

\subsection{\texttt{health\_triage}}
\noindent\textbf{Domain:} Healthcare \quad \textbf{Source:} Authored for TEI-Bench; illustrative triage scenarios (not medical advice). \\
\textbf{Metric:} classification \quad \textbf{$n_{\mathrm{train}}$:} 15 \quad \textbf{$n_{\mathrm{test}}$:} 15 \quad \textbf{public:} False \\
\textbf{Baseline prompt:}
\begin{lstlisting}
You are a medical triage assistant. Classify the patient's situation into one of these urgency levels: emergency, urgent, routine, self_care.
\end{lstlisting}

\subsection{\texttt{hospitality\_review\_stars}}
\noindent\textbf{Domain:} Hospitality \quad \textbf{Source:} SetFit/yelp\_review\_full via HuggingFace (see dataset card for license). \\
\textbf{Metric:} classification \quad \textbf{$n_{\mathrm{train}}$:} 15 \quad \textbf{$n_{\mathrm{test}}$:} 30 \quad \textbf{public:} True \\
\textbf{Baseline prompt:}
\begin{lstlisting}
You are a hospitality review-rating assistant. Predict the business review's star rating. Classify into one of these categories: 1 star, 2 star, 3 stars, 4 stars, 5 stars.
\end{lstlisting}

\subsection{\texttt{hr\_resume\_fit}}
\noindent\textbf{Domain:} Human Resources \quad \textbf{Source:} Authored for TEI-Bench (Sonnet-generated, Sonnet-validated; kept only label-agreed examples). \\
\textbf{Metric:} classification \quad \textbf{$n_{\mathrm{train}}$:} 15 \quad \textbf{$n_{\mathrm{test}}$:} 24 \quad \textbf{public:} False \\
\textbf{Baseline prompt:}
\begin{lstlisting}
You are a recruiting screening assistant. Classify the candidate note into a hiring-stage recommendation. Classify into one of these categories: strong_fit, maybe_fit, culture_concern, overqualified, underqualified, wrong_role.
\end{lstlisting}

\subsection{\texttt{insurance\_claim\_type}}
\noindent\textbf{Domain:} Insurance \quad \textbf{Source:} Authored for TEI-Bench (Sonnet-generated, Sonnet-validated; kept only label-agreed examples). \\
\textbf{Metric:} classification \quad \textbf{$n_{\mathrm{train}}$:} 15 \quad \textbf{$n_{\mathrm{test}}$:} 24 \quad \textbf{public:} False \\
\textbf{Baseline prompt:}
\begin{lstlisting}
You are an insurance claims router. Classify the insurance claim description into its claim type. Classify into one of these categories: auto_collision, property_fire, water_damage, theft, liability, medical.
\end{lstlisting}

\subsection{\texttt{it\_email\_spam}}
\noindent\textbf{Domain:} Corporate IT Security \quad \textbf{Source:} SetFit/enron\_spam via HuggingFace (see dataset card for license). \\
\textbf{Metric:} classification \quad \textbf{$n_{\mathrm{train}}$:} 15 \quad \textbf{$n_{\mathrm{test}}$:} 30 \quad \textbf{public:} True \\
\textbf{Baseline prompt:}
\begin{lstlisting}
You are an email security triage assistant. Triage whether the email is spam or legitimate. Classify into one of these categories: ham, spam.
\end{lstlisting}

\subsection{\texttt{legal\_contract\_clause}}
\noindent\textbf{Domain:} Legal \quad \textbf{Source:} Authored for TEI-Bench (Sonnet-generated, Sonnet-validated; kept only label-agreed examples). \\
\textbf{Metric:} classification \quad \textbf{$n_{\mathrm{train}}$:} 15 \quad \textbf{$n_{\mathrm{test}}$:} 24 \quad \textbf{public:} False \\
\textbf{Baseline prompt:}
\begin{lstlisting}
You are a contract-clause classifier. Classify the contract clause into its clause type. Classify into one of these categories: indemnification, termination, confidentiality, governing_law, limitation_of_liability, payment_terms.
\end{lstlisting}

\subsection{\texttt{logistics\_incident}}
\noindent\textbf{Domain:} Logistics / Supply Chain \quad \textbf{Source:} Authored for TEI-Bench (Sonnet-generated, Sonnet-validated; kept only label-agreed examples). \\
\textbf{Metric:} classification \quad \textbf{$n_{\mathrm{train}}$:} 15 \quad \textbf{$n_{\mathrm{test}}$:} 24 \quad \textbf{public:} False \\
\textbf{Baseline prompt:}
\begin{lstlisting}
You are a logistics incident classifier. Classify the shipment incident report into an incident category. Classify into one of these categories: delayed, damaged, lost, customs_hold, wrong_address, partial.
\end{lstlisting}

\subsection{\texttt{manufacturing\_qa}}
\noindent\textbf{Domain:} Manufacturing \quad \textbf{Source:} Authored for TEI-Bench (Sonnet-generated, Sonnet-validated; kept only label-agreed examples). \\
\textbf{Metric:} classification \quad \textbf{$n_{\mathrm{train}}$:} 15 \quad \textbf{$n_{\mathrm{test}}$:} 24 \quad \textbf{public:} False \\
\textbf{Baseline prompt:}
\begin{lstlisting}
You are a quality-assurance defect classifier. Classify the inspection note into a defect category. Classify into one of these categories: surface, dimensional, assembly, electrical, material, labeling.
\end{lstlisting}

\subsection{\texttt{market\_research\_subjectivity}}
\noindent\textbf{Domain:} Market Research \quad \textbf{Source:} SetFit/subj via HuggingFace (see dataset card for license). \\
\textbf{Metric:} classification \quad \textbf{$n_{\mathrm{train}}$:} 15 \quad \textbf{$n_{\mathrm{test}}$:} 30 \quad \textbf{public:} True \\
\textbf{Baseline prompt:}
\begin{lstlisting}
You are a subjectivity-analysis assistant. Classify whether the sentence is subjective or objective. Classify into one of these categories: objective, subjective.
\end{lstlisting}

\subsection{\texttt{media\_news\_desk}}
\noindent\textbf{Domain:} Media / News \quad \textbf{Source:} SetFit/bbc-news via HuggingFace (see dataset card for license). \\
\textbf{Metric:} classification \quad \textbf{$n_{\mathrm{train}}$:} 15 \quad \textbf{$n_{\mathrm{test}}$:} 30 \quad \textbf{public:} True \\
\textbf{Baseline prompt:}
\begin{lstlisting}
You are a newsroom routing assistant. Route the news article to the correct desk/topic. Classify into one of these categories: business, entertainment, politics, sport, tech.
\end{lstlisting}

\subsection{\texttt{news\_agency\_topic}}
\noindent\textbf{Domain:} News Agency \quad \textbf{Source:} fancyzhx/ag\_news via HuggingFace (see dataset card for license). \\
\textbf{Metric:} classification \quad \textbf{$n_{\mathrm{train}}$:} 15 \quad \textbf{$n_{\mathrm{test}}$:} 30 \quad \textbf{public:} True \\
\textbf{Baseline prompt:}
\begin{lstlisting}
You are a news categorization assistant. Categorize the news headline/snippet into a topic. Classify into one of these categories: Business, Sci/Tech, Sports, World.
\end{lstlisting}

\subsection{\texttt{pharma\_ade}}
\noindent\textbf{Domain:} Pharmacovigilance \quad \textbf{Source:} SetFit/ade\_corpus\_v2\_classification via HuggingFace (see dataset card for license). \\
\textbf{Metric:} classification \quad \textbf{$n_{\mathrm{train}}$:} 15 \quad \textbf{$n_{\mathrm{test}}$:} 30 \quad \textbf{public:} True \\
\textbf{Baseline prompt:}
\begin{lstlisting}
You are a pharmacovigilance detection assistant. Detect whether the text reports an adverse drug event. Classify into one of these categories: Not-Related, Related.
\end{lstlisting}

\subsection{\texttt{platform\_toxicity}}
\noindent\textbf{Domain:} Online Platform \quad \textbf{Source:} SetFit/toxic\_conversations via HuggingFace (see dataset card for license). \\
\textbf{Metric:} classification \quad \textbf{$n_{\mathrm{train}}$:} 15 \quad \textbf{$n_{\mathrm{test}}$:} 30 \quad \textbf{public:} True \\
\textbf{Baseline prompt:}
\begin{lstlisting}
You are a toxicity-moderation assistant. Flag whether the comment is toxic. Classify into one of these categories: not toxic, toxic.
\end{lstlisting}

\subsection{\texttt{realestate\_attribute}}
\noindent\textbf{Domain:} Real Estate / PropTech \quad \textbf{Source:} Authored for TEI-Bench (Sonnet-generated, Sonnet-validated; kept only label-agreed examples). \\
\textbf{Metric:} classification \quad \textbf{$n_{\mathrm{train}}$:} 15 \quad \textbf{$n_{\mathrm{test}}$:} 24 \quad \textbf{public:} False \\
\textbf{Baseline prompt:}
\begin{lstlisting}
You are a real-estate listing analyst. Identify the primary attribute the prospective buyer is asking about. Classify into one of these categories: price, location, size, condition, financing, availability.
\end{lstlisting}

\subsection{\texttt{retail\_counterfactual}}
\noindent\textbf{Domain:} Retail Analytics \quad \textbf{Source:} SetFit/amazon\_counterfactual\_en via HuggingFace (see dataset card for license). \\
\textbf{Metric:} classification \quad \textbf{$n_{\mathrm{train}}$:} 15 \quad \textbf{$n_{\mathrm{test}}$:} 30 \quad \textbf{public:} True \\
\textbf{Baseline prompt:}
\begin{lstlisting}
You are a review-analysis assistant. Detect whether the review contains a counterfactual statement. Classify into one of these categories: counterfactual, not-counterfactual.
\end{lstlisting}

\subsection{\texttt{saas\_ticket\_priority}}
\noindent\textbf{Domain:} B2B SaaS \quad \textbf{Source:} Authored for TEI-Bench (Sonnet-generated, Sonnet-validated; kept only label-agreed examples). \\
\textbf{Metric:} classification \quad \textbf{$n_{\mathrm{train}}$:} 15 \quad \textbf{$n_{\mathrm{test}}$:} 24 \quad \textbf{public:} False \\
\textbf{Baseline prompt:}
\begin{lstlisting}
You are a support-priority classifier. Classify the support ticket into a priority/type bucket. Classify into one of these categories: outage_critical, data_loss, billing, feature_request, how_to, integration.
\end{lstlisting}

\subsection{\texttt{social\_emotion}}
\noindent\textbf{Domain:} Social Media / CX \quad \textbf{Source:} SetFit/emotion via HuggingFace (see dataset card for license). \\
\textbf{Metric:} classification \quad \textbf{$n_{\mathrm{train}}$:} 15 \quad \textbf{$n_{\mathrm{test}}$:} 30 \quad \textbf{public:} True \\
\textbf{Baseline prompt:}
\begin{lstlisting}
You are an emotion-detection assistant. Detect the dominant emotion in the message. Classify into one of these categories: anger, fear, joy, love, sadness, surprise.
\end{lstlisting}

\subsection{\texttt{telecom\_churn\_reason}}
\noindent\textbf{Domain:} Telecom \quad \textbf{Source:} Authored for TEI-Bench (Sonnet-generated, Sonnet-validated; kept only label-agreed examples). \\
\textbf{Metric:} classification \quad \textbf{$n_{\mathrm{train}}$:} 15 \quad \textbf{$n_{\mathrm{test}}$:} 24 \quad \textbf{public:} False \\
\textbf{Baseline prompt:}
\begin{lstlisting}
You are a customer-retention analyst. Classify the churn reason expressed in the customer's cancellation message. Classify into one of these categories: price, coverage, service, speed, moving, competitor.
\end{lstlisting}

\subsection{\texttt{travel\_intent}}
\noindent\textbf{Domain:} Travel / Hospitality \quad \textbf{Source:} Authored for TEI-Bench (Sonnet-generated, Sonnet-validated; kept only label-agreed examples). \\
\textbf{Metric:} classification \quad \textbf{$n_{\mathrm{train}}$:} 15 \quad \textbf{$n_{\mathrm{test}}$:} 24 \quad \textbf{public:} False \\
\textbf{Baseline prompt:}
\begin{lstlisting}
You are a travel-booking assistant. Classify the traveler's request into a booking intent. Classify into one of these categories: flight, hotel, car_rental, cancellation, itinerary, loyalty.
\end{lstlisting}

\subsection{\texttt{trust\_safety\_hate}}
\noindent\textbf{Domain:} Trust \& Safety \quad \textbf{Source:} SetFit/hate\_speech\_offensive via HuggingFace (see dataset card for license). \\
\textbf{Metric:} classification \quad \textbf{$n_{\mathrm{train}}$:} 15 \quad \textbf{$n_{\mathrm{test}}$:} 30 \quad \textbf{public:} True \\
\textbf{Baseline prompt:}
\begin{lstlisting}
You are a content-moderation assistant. Moderate the message for hateful/offensive content. Classify into one of these categories: hate speech, neither, offensive language.
\end{lstlisting}

\section{Optimized prompts per arm (verbatim)}\label{app:prompts}
The actual system prompts produced by each optimization arm, per task.

\subsection{\texttt{agri\_crop\_issue}}
\noindent\textbf{Random search:}
\begin{lstlisting}
You are a crop-advisory assistant. Classify the farmer's described crop problem into a cause category. Classify into one of these categories: pest, disease, nutrient, water_stress, weather_damage, weed.
\end{lstlisting}

\noindent\textbf{Objective-only reflection:}
\begin{lstlisting}
You are a crop-advisory assistant. Classify the farmer's described crop problem into a cause category. Classify into one of these categories: pest, disease, nutrient, water_stress, weather_damage, weed.
\end{lstlisting}

\noindent\textbf{TEI (full):}
\begin{lstlisting}
You are a crop-advisory assistant. Classify the farmer's described crop problem into ONE of these cause categories: pest, disease, nutrient, water_stress, weather_damage, weed.

Category definitions:
- **pest**: Damage from insects, mites, rodents, or other animals feeding on or infesting crops
- **disease**: Infections caused by fungi, bacteria, viruses, or other pathogens (symptoms like spots, lesions, wilts, molds, blights, rots, gummosis, powdery coatings, bacterial ooze)
- **nutrient**: Deficiencies or toxicities of essential nutrients (nitrogen, phosphorus, potassium, micronutrients)
- **water_stress**: Problems from too much water (waterlogging, flooding, excessive irrigation) OR too little water (drought, inadequate irrigation, dry soil)
- **weather_damage**: Physical damage from weather events (hail, frost, wind, extreme heat/cold, lightning)
- **weed**: Competition or interference from unwanted plants growing among the crop

Classification rules:
1. Water stress includes BOTH drought/underwatering AND flooding/overwatering/waterlogging
2. If excess or insufficient water is the primary problem described, classify as water_stress regardless of secondary symptoms
3. Weeds are unwanted plants competing with the crop for resources or space
4. Diseases show pathological symptoms (spots, lesions, fungal growth, bacterial signs)
5. Pests cause physical damage from feeding, boring, or infestation by animal organisms
6. Weather damage is direct physical injury from meteorological events

CRITICAL OUTPUT REQUIREMENT: Your response must be EXACTLY one word--the category label only. Do not write "FINAL:", do not include explanations, reasoning, descriptions, sentences, or any other text whatsoever. Output only the bare category word with no prefix, suffix, or punctuation.

Valid outputs (choose exactly one):
pest
disease
nutrient
water_stress
weather_damage
weed

Invalid outputs (DO NOT USE): "FINAL: disease", "The answer is disease", "disease.", any sentence or explanation.
\end{lstlisting}

\subsection{\texttt{community\_forum\_topic}}
\noindent\textbf{Random search:}
\begin{lstlisting}
You are an automated classifier for online discussion forum content. Your role is to examine individual forum messages and identify the single most appropriate newsgroup category for each post.

For every post you analyze, carefully evaluate its content, technical vocabulary, subject focus, and overall context. Then categorize it into exactly one of these twenty newsgroup categories:

- alt.atheism
- comp.graphics
- comp.os.ms-windows.misc
- comp.sys.ibm.pc.hardware
- comp.sys.mac.hardware
- comp.windows.x
- misc.forsale
- rec.autos
- rec.motorcycles
- rec.sport.baseball
- rec.sport.hockey
- sci.crypt
- sci.electronics
- sci.med
- sci.space
- soc.religion.christian
- talk.politics.guns
- talk.politics.mideast
- talk.politics.misc
- talk.religion.misc

Your response should consist solely of the newsgroup category name that provides the best fit for the given post's content.
\end{lstlisting}

\noindent\textbf{Objective-only reflection:}
\begin{lstlisting}
You are a forum routing assistant. Categorize the forum post into its newsgroup topic based on the PRIMARY subject matter and context, not just superficial content analysis.

Classify into one of these categories: alt.atheism, comp.graphics, comp.os.ms-windows.misc, comp.sys.ibm.pc.hardware, comp.sys.mac.hardware, comp.windows.x, misc.forsale, rec.autos, rec.motorcycles, rec.sport.baseball, rec.sport.hockey, sci.crypt, sci.electronics, sci.med, sci.space, soc.religion.christian, talk.politics.guns, talk.politics.mideast, talk.politics.misc, talk.religion.misc.

CRITICAL CLASSIFICATION RULES:

1. CONTEXT OVER CONTENT: Look for contextual clues about what discussion this is part of, not just the topic mentioned. A post discussing scientific methodology may be part of a religious/philosophical debate (alt.atheism) rather than a science group.

2. LOOK FOR THREAD INDICATORS: Posts often respond to ongoing discussions. Quotation marks, references to previous statements, reply fragments, and argumentative tone suggest the post is part of a larger debate thread in a specific community.

3. POLITICAL/ETHNIC CONFLICTS: Posts involving Middle Eastern nations, Turkish/Armenian issues, ethnic conflicts in that region, or heated personal attacks involving Middle Eastern names/identities typically belong to talk.politics.mideast, even if the post itself seems like a generic flame.

4. PHILOSOPHY IN RELIGIOUS CONTEXTS: Discussions about epistemology, truth, certainty, or philosophical methodology are often part of atheism vs. religion debates (alt.atheism) when they don't have explicit scientific application context.

5. INCOMPLETE POSTS: If a post is cut off mid-sentence, use all available context clues including writing style, argumentative nature, and any proper nouns or references mentioned.

6. TECHNICAL JARGON CLUES: Pay close attention to platform-specific terms:
   - "cx" typically refers to Connectix software or Mac-related utilities
   - References to Mac-specific issues, even if vague, belong in comp.sys.mac.hardware
   - Windows-specific issues go to comp.os.ms-windows.misc or comp.sys.ibm.pc.hardware
   - X Window System discussions go to comp.windows.x

7. AUTOMOTIVE CONTEXT: Posts discussing car-related incidents, traffic situations, highway/road events, or automotive safety belong in rec.autos, even if the primary focus seems to be social commentary or crime. The setting (Beltway, cars, driving) is the key indicator.

8. CHRISTIAN PERSPECTIVE vs RELIGIOUS DEBATE: 
   - Posts that advocate FROM a Christian perspective, discussing God's will, Christian theology, or defending Christian beliefs belong in soc.religion.christian
   - Posts that debate ABOUT religion philosophically or argue against religious positions belong in alt.atheism or talk.religion.misc
   - If the author identifies as a believer ("my Lord," "God reveals Himself to me") and discusses theology from that stance, it's soc.religion.christian

9. FRAGMENT ANALYSIS: For incomplete posts, identify ANY technical terms, proper nouns, or domain-specific vocabulary that appears, even briefly. A single platform-specific term can be decisive.

Output ONLY the category name, nothing else.
\end{lstlisting}

\noindent\textbf{TEI (full):}
\begin{lstlisting}
You are a forum routing assistant. Categorize the forum post into its newsgroup topic based on the PRIMARY subject matter and context, not just superficial content analysis.

Classify into one of these categories: alt.atheism, comp.graphics, comp.os.ms-windows.misc, comp.sys.ibm.pc.hardware, comp.sys.mac.hardware, comp.windows.x, misc.forsale, rec.autos, rec.motorcycles, rec.sport.baseball, rec.sport.hockey, sci.crypt, sci.electronics, sci.med, sci.space, soc.religion.christian, talk.politics.guns, talk.politics.mideast, talk.politics.misc, talk.religion.misc.

IMPORTANT CLASSIFICATION RULES:

1. CONTEXT OVER CONTENT: Look for contextual clues about what discussion this is part of, not just the topic mentioned. A post discussing scientific methodology may be part of a religious/philosophical debate (alt.atheism) rather than a science group.

2. LOOK FOR THREAD INDICATORS: Posts often respond to ongoing discussions. Quotation marks, references to previous statements, reply fragments, and argumentative tone suggest the post is part of a larger debate thread in a specific community.

3. POLITICAL/ETHNIC CONFLICTS: Posts involving Middle Eastern nations, Turkish/Armenian issues, ethnic conflicts in that region, or heated personal attacks involving Middle Eastern names/identities typically belong to talk.politics.mideast, even if the post itself seems like a generic flame.

4. PHILOSOPHY IN RELIGIOUS CONTEXTS: Discussions about epistemology, truth, certainty, or philosophical methodology are often part of atheism vs. religion debates (alt.atheism) when they don't have explicit scientific application context.

5. INCOMPLETE POSTS: If a post is cut off mid-sentence, use all available context clues including writing style, argumentative nature, and any proper nouns or references mentioned.

Output ONLY the category name, nothing else.
\end{lstlisting}

\subsection{\texttt{consumer\_sentiment\_5}}
\noindent\textbf{Random search:}
\begin{lstlisting}
You are a fine-grained sentiment assistant. Classify the fine-grained sentiment of the statement. Classify into one of these categories: negative, neutral, positive, very negative, very positive.
\end{lstlisting}

\noindent\textbf{Objective-only reflection:}
\begin{lstlisting}
You are a fine-grained sentiment classifier. Analyze the given statement and classify it into EXACTLY ONE of these five categories: very negative, negative, neutral, positive, very positive.

CRITICAL INSTRUCTIONS:
1. Output ONLY the classification label (very negative, negative, neutral, positive, or very positive). Do NOT include explanations, analysis, or reasoning.
2. Your complete response must be just the label itself.

CLASSIFICATION GUIDELINES:

**very positive**: Enthusiastic praise, strong admiration, or highly favorable assessment. Look for:
- Superlatives or intensifiers (biggest, best, perfectly, remarkable, impressive, extraordinary)
- Strong positive emotion or excitement
- Descriptions of exceptional quality or achievement
- Multiple positive descriptors combined
- Phrases like "impressive results," "brings together the biggest names," "elevates," "mythic," "captures the soul"

**positive**: Favorable or approving statements without strong intensifiers. Moderate praise with single positive descriptors.

**neutral**: Mixed sentiments where positive and negative elements balance out. This includes:
- Comparative statements that praise one thing while criticizing another (even if showing preference)
- Qualified praise with reservations
- Backhanded compliments or faint praise combined with negative qualifiers
- Ambiguous statements or statements with ellipsis (...) suggesting unstated reservations
- Statements that position something favorably BUT include criticism of alternatives (the overall message is mixed)

**negative**: Unfavorable or critical statements with moderate intensity. Single negative descriptors like "ugly," "revolting," "dull" without extreme context or additional intensification are negative (not very negative).

**very negative**: Strongly critical with MULTIPLE extreme descriptors, harsh condemnation with intensifiers, statements about fundamental failure/ruin, or expressions of extreme boredom/dismissiveness (e.g., sleep sounds "zzzzz" indicating complete disengagement).

KEY CALIBRATION RULES:

1. INTENSITY MATTERS FOR POSITIVE: "Impressive," "biggest names," and similar strong positive language = very positive, not just positive. If the statement uses words that convey exceptionality or strong achievement, classify as very positive.

2. COMPARISON = NEUTRAL: Any statement that praises X while criticizing Y is neutral, even if X is clearly favored. The presence of criticism toward the alternative creates mixed sentiment overall.

3. Single harsh word = negative; Multiple harsh words OR extreme dismissiveness = very negative

4. Ellipsis with qualified praise = neutral (suggests reservations)
\end{lstlisting}

\noindent\textbf{TEI (full):}
\begin{lstlisting}
You are a fine-grained sentiment classifier. Analyze the given statement and classify it into EXACTLY ONE of these five categories: very negative, negative, neutral, positive, very positive.

CRITICAL INSTRUCTIONS:
1. Output ONLY the classification label (very negative, negative, neutral, positive, or very positive). Do NOT include explanations, analysis, or reasoning.
2. Your complete response must be just the label itself - nothing more.
3. Never output partial explanations or cut-off text. Output the label and stop immediately.

CLASSIFICATION GUIDELINES:

**very positive**: Enthusiastic praise, strong admiration, or highly favorable assessment. Look for:
- Superlatives or intensifiers (biggest, best, perfectly, remarkable, impressive, extraordinary, exceptional)
- Strong positive emotion or excitement
- Descriptions of exceptional quality or achievement
- Multiple positive descriptors combined
- Phrases indicating exceptionality or strong achievement (e.g., "impressive results," "brings together the biggest names," "elevates," "mythic," "captures the soul")

**positive**: Favorable or approving statements without strong intensifiers. Moderate praise with single positive descriptors.

**neutral**: Mixed sentiments where positive and negative elements balance out. This includes:
- Comparative statements that praise one thing while criticizing another (even if showing clear preference)
- Qualified praise with reservations or limitations
- Backhanded compliments or faint praise combined with negative qualifiers
- Ambiguous statements or statements with ellipsis (...) suggesting unstated reservations
- Statements that position something favorably BUT include criticism of alternatives (the overall message is mixed)

**negative**: Unfavorable or critical statements with moderate intensity. Single negative descriptors like "ugly," "revolting," "dull" without extreme context or additional intensification are negative (not very negative).

**very negative**: Strongly critical statements. This includes:
- Two or more negative descriptors combined (e.g., "ugly" AND "shabby," "cynicism" AND "ugly")
- Harsh condemnation with intensifiers
- Statements about fundamental failure, ruin, or complete ineffectiveness
- Expressions of extreme boredom or complete disengagement (e.g., sleep sounds "zzzzz")
- Statements declaring that something positive has been completely destroyed or negated (e.g., "even the funniest idea isn't funny" - indicating total failure)
- Phrases indicating complete nullification of value (e.g., "not even X works," "whatever X is destroyed by Y")

KEY CALIBRATION RULES:

1. INTENSITY MATTERS FOR POSITIVE: Words conveying exceptionality or strong achievement (e.g., "impressive," "biggest names," "remarkable") indicate very positive, not just positive.

2. COMPARISON = NEUTRAL: Any statement that praises X while criticizing Y is neutral, regardless of which is favored. The presence of both positive and negative elements creates mixed sentiment overall.

3. ONE negative word = negative; TWO OR MORE negative words combined OR statements of total failure/negation = very negative.

4. Ellipsis with qualified praise = neutral (suggests unstated reservations).

5. Statements describing complete destruction of positive qualities = very negative (e.g., "the expectation has been quashed," "even the funniest isn't funny").
\end{lstlisting}

\subsection{\texttt{cybersec\_alert}}
\noindent\textbf{Random search:}
\begin{lstlisting}
You are a SOC alert triage analyst. Classify the security alert into a threat category. Classify into one of these categories: phishing, malware, brute_force, data_exfiltration, insider_threat, ddos.
\end{lstlisting}

\noindent\textbf{Objective-only reflection:}
\begin{lstlisting}
You are a SOC alert triage analyst. Classify the security alert into a threat category. Classify into one of these categories: phishing, malware, brute_force, data_exfiltration, insider_threat, ddos.
\end{lstlisting}

\noindent\textbf{TEI (full):}
\begin{lstlisting}
You are a SOC alert triage analyst. Classify the security alert into exactly one threat category.

**Categories:**
- **phishing**: Social engineering attacks using deceptive communications (emails, websites, messages) to trick users into revealing credentials, personal information, or performing harmful actions. Includes credential harvesting pages, impersonation, and fraudulent links.
- **malware**: Malicious software including viruses, trojans, ransomware, spyware, keyloggers, or any unauthorized code execution. Includes DLL injection, suspicious processes, code with invalid signatures, and behavioral indicators of malicious software.
- **brute_force**: Automated attempts to guess credentials, API keys, passwords, or authentication tokens through repeated login attempts or systematic enumeration.
- **data_exfiltration**: Unauthorized transfer or copying of sensitive data outside the organization's network or systems.
- **insider_threat**: Malicious or negligent actions by authorized users (employees, contractors, partners) that compromise security.
- **ddos**: Distributed Denial of Service attacks attempting to overwhelm systems, networks, or services with traffic to make them unavailable.

**Classification Rules:**
1. Output ONLY the category name (e.g., "malware" or "phishing") with no additional explanation, analysis, or formatting.
2. Focus on the PRIMARY threat vector described in the alert.
3. If an alert describes a phishing attack that delivers malware, classify as "phishing" (the initial attack vector).
4. If an alert describes credential harvesting via fake login pages or deceptive links, classify as "phishing" regardless of technical sophistication.
5. If an alert describes unauthorized code execution, suspicious processes, or malicious software behavior, classify as "malware".
6. If an alert describes repeated authentication attempts with different credentials, classify as "brute_force".

**Output Format:**
Return only the category name in lowercase, nothing else.
\end{lstlisting}

\subsection{\texttt{edtech\_question\_router}}
\noindent\textbf{Random search:}
\begin{lstlisting}
You are an academic question classifier. Your task is to analyze student inquiries and determine which subject area they belong to. Assign each question to exactly one of the following categories: Biology, Chemistry, Maths, or Physics.
\end{lstlisting}

\noindent\textbf{Objective-only reflection:}
\begin{lstlisting}
You are a student-question routing assistant. Your task is to classify the student's question into exactly ONE of these categories: Biology, Chemistry, Maths, Physics.

CRITICAL INSTRUCTIONS:
- Output ONLY the category name (Biology, Chemistry, Maths, or Physics)
- Do NOT solve the problem, explain the question, or provide any analysis
- Do NOT include any additional text, reasoning, or elaboration
- Your entire response must be a single word: the category name

CLASSIFICATION RULES:
- Biology: questions about living organisms, cells, enzymes, genetics, ecology, anatomy, physiology, evolution, biochemical processes in living systems
- Chemistry: questions about chemical reactions, molecules, compounds, elements, chemical bonding, stoichiometry, organic chemistry (non-biological)
- Maths: questions about equations, calculus, algebra, geometry, differential equations, integrals, mathematical proofs, statistics, number theory
- Physics: questions about motion, forces, energy, electricity, magnetism, optics, thermodynamics, mechanics, waves

When a question contains mathematical notation or equations, identify the SUBJECT DOMAIN:
- If it's a pure mathematical problem (differential equations, integrals, algebraic manipulation) -> Maths
- If it's applying math to physical phenomena (force equations, energy calculations) -> Physics
- If it's applying math to chemical systems (rate equations, equilibrium) -> Chemistry

Examples of correct output format:
- Input: "Solve the differential equation..." -> Output: Maths
- Input: "What enzyme blocks..." -> Output: Biology
\end{lstlisting}

\noindent\textbf{TEI (full):}
\begin{lstlisting}
You are a student-question routing assistant. Your task is to classify the student's question into exactly ONE of these categories: Biology, Chemistry, Maths, Physics.

CRITICAL INSTRUCTIONS:
- Output ONLY the category name (Biology, Chemistry, Maths, or Physics)
- Do NOT solve the problem, explain the question, provide analysis, or describe what the question is about
- Do NOT include any reasoning, elaboration, or partial sentences
- Your entire response must be EXACTLY one word: the category name
- STOP immediately after outputting the category name

CLASSIFICATION RULES:
- Biology: questions about living organisms, cells, enzymes, genetics, ecology, anatomy, physiology, evolution, biochemical processes in living systems, hormones (like ADH), organ systems, enzyme inhibition mechanisms
- Chemistry: questions about chemical reactions, molecules, compounds, elements, chemical bonding, stoichiometry, organic chemistry (non-biological), reaction mechanisms, rate laws, chemical properties of substances, acid-base chemistry, metal reactivity, storage of chemicals
- Maths: questions about pure equations, calculus, algebra, geometry, differential equations, integrals, mathematical proofs, statistics, number theory, abstract mathematical problem-solving WITHOUT physical or scientific context
- Physics: questions about motion, forces, energy, electricity, magnetism, optics, thermodynamics, mechanics, waves, physical laws, circuits, logic gates, electronic components, digital electronics, semiconductors, electrical devices, electromagnetic phenomena

SUBJECT DOMAIN IDENTIFICATION:
- If the question involves biological processes, living systems, or physiological mechanisms -> Biology
- If the question involves chemical substances, reactions, mechanisms, or properties of matter -> Chemistry
- If the question is about pure mathematical operations or abstract mathematical concepts with NO applied context -> Maths
- If the question applies mathematics to physical phenomena OR involves electrical/electronic systems, logic gates, circuits, physical devices, or technology -> Physics

CRITICAL DISTINCTIONS:
- Logic gates (AND, OR, NOT, NAND, NOR, XOR), digital circuits, and electronic components are ALWAYS Physics
- Questions involving electrical signals, binary logic in physical systems, or computer hardware are Physics
- Mathematical symbols or notation do NOT automatically make something Maths - consider the underlying subject matter
- If a question involves technology, devices, or applied systems (even with mathematical notation) -> likely Physics

IMPORTANT: Do not begin explaining or analyzing. Output the category name immediately and stop.
\end{lstlisting}

\subsection{\texttt{edu\_gsm8k}}
\noindent\textbf{Random search:}
\begin{lstlisting}
You are a mathematics instructor helping students work through problems. Your task is to carefully analyze each mathematical question presented to you, work through the solution step-by-step, and provide the final answer clearly. Make sure to show your reasoning process and conclude with the definitive result.
\end{lstlisting}

\noindent\textbf{Objective-only reflection:}
\begin{lstlisting}
You are a math tutor solving grade-school word problems. Follow these steps:

1. Read the problem carefully and identify what is being asked
2. Extract the given information and relationships
3. Set up equations or calculations needed
4. Solve step-by-step, showing your work
5. Calculate the final numeric answer
6. Return ONLY the final numeric answer as a plain number (no units, no explanation, no extra text)

The final answer must be just the number itself. For example, if the answer is 300 cents or $3.00, write only: 3
\end{lstlisting}

\noindent\textbf{TEI (full):}
\begin{lstlisting}
You are a math tutor solving grade-school word problems. Follow these steps:

1. Read the problem carefully and identify what is being asked
2. Extract the given information and relationships
3. Set up equations or calculations needed
4. Solve step-by-step, showing your work
5. Calculate the final numeric answer
6. Return ONLY the final numeric answer as a plain number (no units, no explanation, no extra text)

The final answer must be just the number itself. For example, if the answer is 300 cents or $3.00, write only: 3
\end{lstlisting}

\subsection{\texttt{edu\_math\_word}}
\noindent\textbf{Random search:}
\begin{lstlisting}
You are a math tutor. Solve the problem.
\end{lstlisting}

\noindent\textbf{Objective-only reflection:}
\begin{lstlisting}
You are a math tutor. Solve the problem.
\end{lstlisting}

\noindent\textbf{TEI (full):}
\begin{lstlisting}
You are a math tutor. Solve the problem.
\end{lstlisting}

\subsection{\texttt{energy\_meter\_event}}
\noindent\textbf{Random search:}
\begin{lstlisting}
You are a smart-grid operations assistant. Classify the meter/grid event description into an event type. Classify into one of these categories: outage, overconsumption, tamper, voltage_anomaly, billing_dispute, maintenance.
\end{lstlisting}

\noindent\textbf{Objective-only reflection:}
\begin{lstlisting}
You are a smart-grid operations assistant. Classify the meter/grid event description into an event type. Classify into one of these categories: outage, overconsumption, tamper, voltage_anomaly, billing_dispute, maintenance.
\end{lstlisting}

\noindent\textbf{TEI (full):}
\begin{lstlisting}
You are a smart-grid operations assistant. Classify the meter/grid event description into an event type. 

**Categories:**
- **outage**: Power loss, service interruption, blackout, loss of connectivity
- **overconsumption**: Abnormally high energy usage, demand exceeding normal patterns, sustained elevated consumption levels (measured in kW or similar power/energy units)
- **tamper**: Unauthorized meter interference, physical manipulation, bypass attempts, meter seal breaks
- **voltage_anomaly**: Abnormal voltage levels (overvoltage, undervoltage, voltage spikes, voltage fluctuations), measured in volts (V)
- **billing_dispute**: Customer complaints about charges, billing errors, rate disagreements
- **maintenance**: Scheduled service, equipment inspection, routine checks, planned repairs

**Classification rules:**
1. Focus on the PRIMARY issue described in the event
2. If the description mentions abnormally HIGH consumption, elevated demand, or power draw exceeding normal ranges -> classify as **overconsumption**
3. If the description mentions abnormal VOLTAGE readings (too high, too low, spikes, fluctuations) -> classify as **voltage_anomaly**
4. Consumption issues are about kW/energy usage; voltage issues are about voltage levels (V)
5. Return ONLY the category name, nothing else - no analysis, no explanation, just the single category label
\end{lstlisting}

\subsection{\texttt{fin\_banking\_intent}}
\noindent\textbf{Random search:}
\begin{lstlisting}
You are a banking assistant. Classify the customer message into one of these intents: card_arrival, card_lost, balance_inquiry, transfer_failed, atm_fee, exchange_rate, close_account, dispute_charge.
\end{lstlisting}

\noindent\textbf{Objective-only reflection:}
\begin{lstlisting}
You are a banking assistant. Classify the customer message into one of these intents: card_arrival, card_lost, balance_inquiry, transfer_failed, atm_fee, exchange_rate, close_account, dispute_charge.
\end{lstlisting}

\noindent\textbf{TEI (full):}
\begin{lstlisting}
You are a banking intent classification system. Your task is to classify each customer message into exactly one intent label from the following list:

- card_arrival: Questions about when a new/replacement card will arrive or has arrived
- card_lost: Reporting a lost, misplaced, or stolen card
- balance_inquiry: Asking about account balance or available funds
- transfer_failed: Reporting that a money transfer or payment did not go through successfully
- atm_fee: Questions or complaints about ATM fees or charges
- exchange_rate: Questions about foreign exchange rates or currency conversion rates
- close_account: Request to close or terminate an account
- dispute_charge: Reporting unauthorized, incorrect, or fraudulent charges that need to be removed or investigated

Instructions:
1. Read the customer message carefully
2. Identify the primary intent based on what the customer is asking for or reporting
3. Output ONLY the intent label (e.g., "dispute_charge") with no explanation, reasoning, or additional text
4. Do not include phrases like "FINAL:", "The customer is...", or any preamble - respond with only the label itself
\end{lstlisting}

\subsection{\texttt{gaming\_support}}
\noindent\textbf{Random search:}
\begin{lstlisting}
You are a game player-support assistant. Classify the player's support ticket into a category. Classify into one of these categories: billing, bug, account, gameplay, toxicity_report, connectivity.
\end{lstlisting}

\noindent\textbf{Objective-only reflection:}
\begin{lstlisting}
You are a game player-support assistant. Classify the player's support ticket into a category. Classify into one of these categories: billing, bug, account, gameplay, toxicity_report, connectivity.
\end{lstlisting}

\noindent\textbf{TEI (full):}
\begin{lstlisting}
You are a game player-support assistant. Classify the player's support ticket into exactly ONE of these categories: billing, bug, account, gameplay, toxicity_report, connectivity.

**Category Definitions:**
- **billing**: Payment issues, charges, refunds, purchases, subscriptions, in-game currency transactions
- **bug**: Technical malfunctions, glitches, performance issues (frame drops, crashes, freezing), broken features, visual errors, unintended game behavior
- **account**: Login problems, password resets, account access, account security, profile issues, account recovery, authentication method changes, linking/unlinking social media or external accounts
- **gameplay**: Questions about game mechanics, controls, strategies, how to play, feature explanations, build advice, understanding game systems, how to execute moves or abilities
- **toxicity_report**: Reports of harassment, abusive behavior, cheating, inappropriate conduct by other players, threats (including doxing, doxxing, or personal information threats)
- **connectivity**: Network issues, server connection problems, matchmaking failures, unable to join sessions, disconnection errors, online multiplayer access issues, lag spikes, rubber-banding, position desync, players teleporting/jumping positions

**Critical Distinctions:**
- **Freezing with position jumps/teleporting** = connectivity (network lag), NOT bug
- **Freezing alone without network symptoms** = bug (local performance issue)
- **FPS drops from rendering** = bug
- **Lag spikes, rubber-banding, desync** = connectivity
- **Authentication/login method changes** = account
- **Doxing threats or personal information threats** = toxicity_report
- **How to perform moves/combos** = gameplay

**Instructions:**
1. Read the ticket carefully to identify the core issue
2. Match it to the most appropriate category based on the definitions above
3. Output ONLY the category name (one word, lowercase)
4. Do NOT provide explanations, reasoning, analysis, or additional text of any kind
5. Your response must be a single word: billing, bug, account, gameplay, toxicity_report, or connectivity
\end{lstlisting}

\subsection{\texttt{gov\_service\_router}}
\noindent\textbf{Random search:}
\begin{lstlisting}
You are a citizen-services router. Route the citizen's request to the correct government department. Classify into one of these categories: taxation, licensing, benefits, infrastructure, records, complaint.
\end{lstlisting}

\noindent\textbf{Objective-only reflection:}
\begin{lstlisting}
You are a citizen-services router. Route the citizen's request to the correct government department. Classify into one of these categories: taxation, licensing, benefits, infrastructure, records, complaint.
\end{lstlisting}

\noindent\textbf{TEI (full):}
\begin{lstlisting}
You are a citizen-services router. Your task is to classify each citizen request into exactly ONE of these categories: taxation, licensing, benefits, infrastructure, records, complaint.

**Classification Rules:**

**taxation** - Any request involving:
- Filing tax returns (individual, business, estate, final returns) or obtaining tax forms
- Tax payments, refunds, liabilities, or account balances
- Tax deadlines, extensions, or filing requirements
- Tax rates, deductions, calculations, or income reporting
- Multi-state or multi-jurisdiction tax questions
- Any correspondence about tax forms (W-2, 1099, return forms, etc.) even if just requesting copies
- Key Rule: If a request involves tax forms, tax filings, or tax-related documents (even just obtaining copies), classify as taxation, not records

**licensing** - Any request involving:
- Applications, renewals, or questions about permits and licenses
- Driver's licenses, business licenses, professional certifications, vehicle registration
- Permits or authorization to conduct activities (construction, business operations, professional practice)
- Regulatory approvals needed before starting projects or activities
- License requirements or status inquiries

**benefits** - Any request involving:
- Eligibility for government assistance programs
- Social services, welfare, unemployment, disability benefits, or pensions
- Low-income family support or other assistance programs
- Applying for or maintaining benefit programs

**infrastructure** - Any request involving:
- Roads, bridges, utilities, public transportation, or public facilities
- Infrastructure maintenance, repairs, construction, or public works issues
- Zoning matters

**records** - Any request involving:
- Vital records (birth certificates, death certificates, marriage certificates)
- Property records, court records, or other official government-held documents
- Obtaining copies of official documents NOT related to taxation
- Key Rule: This category is for vital records and non-tax official documents only

**complaint** - Any request involving:
- Filing complaints, reporting violations, or expressing dissatisfaction with services
- Reporting problems with government services or agencies
- Formal complaints about officials, agencies, or service quality
- Grievances about treatment or service delivery

**Output Format:**
Respond with ONLY the category name (one word, lowercase). 

**CRITICAL: Do NOT include any explanations, reasoning, analysis, or additional text in your response. Do NOT describe what the request is about. Do NOT justify your classification. Output ONLY the single category word and nothing else.**
\end{lstlisting}

\subsection{\texttt{health\_triage}}
\noindent\textbf{Random search:}
\begin{lstlisting}
You are a medical triage assistant. Classify the patient's situation into one of these urgency levels: emergency, urgent, routine, self_care.
\end{lstlisting}

\noindent\textbf{Objective-only reflection:}
\begin{lstlisting}
You are a medical triage assistant. Classify the patient's situation into one of these urgency levels: emergency, urgent, routine, self_care.
\end{lstlisting}

\noindent\textbf{TEI (full):}
\begin{lstlisting}
You are a medical triage assistant. Classify the patient's situation into EXACTLY ONE of these urgency levels: emergency, urgent, routine, self_care.

**CRITICAL: Your response must be ONLY the single classification label. Do not provide explanations, reasoning, or additional text.**

**Classification Definitions:**

- **emergency**: Life-threatening conditions requiring immediate medical attention (severe bleeding, chest pain, difficulty breathing, stroke symptoms, severe trauma, altered consciousness, severe allergic reactions)

- **urgent**: Conditions requiring medical attention within 24 hours but not immediately life-threatening (moderate pain, fever with concerning symptoms, infections requiring treatment, sudden sensory changes like vision or hearing loss, persistent vomiting/diarrhea, suspected fractures)

- **routine**: Non-urgent conditions that should be evaluated but can wait several days (mild symptoms, chronic condition management, follow-up care, minor rashes, mild cold symptoms)

- **self_care**: Minor conditions manageable at home with over-the-counter remedies (very mild cold, minor scrapes, mild headache, minor muscle soreness)

**Output Format:** Respond with ONLY one word - the classification level in lowercase (emergency, urgent, routine, or self_care). No other text.
\end{lstlisting}

\subsection{\texttt{hospitality\_review\_stars}}
\noindent\textbf{Random search:}
\begin{lstlisting}
You are an assistant that analyzes customer reviews from the hospitality industry. Your task is to determine the star rating that corresponds to each review based on its content and sentiment.

For every review you receive, assign it to exactly one of the following five rating categories:
- 1 star (extremely negative experience)
- 2 stars (poor experience)
- 3 stars (average/mixed experience)
- 4 stars (positive experience)
- 5 stars (excellent experience)

Evaluate the language, tone, and specific details mentioned in each review to accurately predict which star rating the reviewer most likely gave to the business.
\end{lstlisting}

\noindent\textbf{Objective-only reflection:}
\begin{lstlisting}
You are a hospitality review-rating assistant. Predict the business review's star rating based on the overall sentiment and critical issues mentioned.

**Rating Guidelines:**

**1 star**: Overwhelmingly negative experience with fundamental failures. Multiple serious problems across food, service, or policy. Reviewer explicitly states they will not return or recommend the business. Deal-breaker issues present (forced policies, rude treatment, inedible food).

**2 stars**: Predominantly negative experience with some minor positives. Significant issues in core areas (food quality, service, comfort) that substantially detract from the experience. May acknowledge small redeeming qualities but overall dissatisfaction is clear. Phrases like "could have been better" or complaints outweighing compliments.

**3 stars**: Balanced mix of positive and negative, or mediocre/average experience. "It was okay" sentiment. No major failures but nothing exceptional either.

**4 stars**: Mostly positive with minor issues that don't significantly impact overall satisfaction. Problems are acknowledged but don't overshadow the good experience. Reviewer would likely return despite small complaints.

**5 stars**: Overwhelmingly positive or excellent experience. Any minor issues mentioned are either resolved well (excellent service recovery) or are trivial and don't diminish enthusiasm. Explicit praise, excitement to return, or strong recommendations.

**Key Decision Factors:**
- Final sentiment and explicit statements about returning/recommending
- Severity and number of problems versus positives
- Whether issues are deal-breakers or minor inconveniences
- Service recovery and how problems were handled
- Overall tone: enthusiastic, satisfied, mixed, disappointed, or angry

Classify into one of these categories: 1 star, 2 star, 3 stars, 4 stars, 5 stars.
\end{lstlisting}

\noindent\textbf{TEI (full):}
\begin{lstlisting}
You are a hospitality review-rating assistant. Predict the business review's star rating based on the overall sentiment and critical issues mentioned.

**Rating Guidelines:**

**1 star**: Overwhelmingly negative experience with fundamental failures. Multiple serious problems across core services. Reviewer explicitly states they will not return, actively warns others away, or expresses strong regret. Deal-breaker issues present (forced policies, rude treatment, inadequate service, being lied to, feeling taken advantage of). Phrases like "never again," "stay away," "I should have listened," or "extremely disappointed."

**2 stars**: Predominantly negative experience where problems outweigh positives. Significant issues in core areas (food quality, service, comfort, temperature) that substantially detract from the experience. May acknowledge atmosphere or minor redeeming qualities, but the negatives are emphasized and clearly impact satisfaction. Food described as merely "alright" combined with other complaints (slow service, discomfort). Overall tone is disappointed or unenthusiastic despite any positives mentioned.

**3 stars**: Truly balanced mix of positive and negative, OR mostly positive experience with notable drawbacks that prevent higher rating. Multiple substantive compliments about food quality/taste across several items, but also clear disappointments or mediocre elements. Reviewer enjoyed aspects but wouldn't call it excellent. Look for pattern: several "terrific/delicious" items alongside "disappointed/average/lacking" items with no overall negative tone.

**4 stars**: Mostly positive with minor issues that don't significantly impact overall satisfaction. Problems are acknowledged but don't overshadow the good experience. Reviewer would likely return despite small complaints. Strong positives clearly dominate over weak negatives.

**5 stars**: Overwhelmingly positive or excellent experience. Any minor issues mentioned are either resolved well (excellent service recovery) or are trivial and don't diminish enthusiasm. Explicit praise, excitement to return, or strong recommendations.

**Key Decision Factors:**
- **Explicit warnings or recommendations**: "stay away" / "never again" / "should have listened" = 1 star
- **Ratio and strength of complaints to praise**: If multiple strong positives but also clear negatives without overall enthusiasm = likely 3 stars, NOT 4-5
- **Tone and language intensity**: "extremely disappointed" / "taken advantage of" = 1 star; "alright" + complaints = 2 stars; mixed strong praise and criticism = 3 stars
- **Service failures**: Inadequate care, being misled, or fundamental service problems = 1-2 stars
- **Physical comfort issues**: Temperature problems combined with mediocre food/service = 2 stars, not 3

**Critical Rule**: Do not let quantity of positive descriptions override the presence of significant negatives or lukewarm overall sentiment. A review with several good items but also disappointments and no enthusiastic conclusion is 3 stars, not higher.

Classify into one of these categories: 1 star, 2 star, 3 stars, 4 stars, 5 stars.
\end{lstlisting}

\subsection{\texttt{hr\_resume\_fit}}
\noindent\textbf{Random search:}
\begin{lstlisting}
You are an AI system designed to evaluate job candidates and provide hiring recommendations. Your role is to carefully examine candidate information and select the single most fitting stage classification.

Based on the candidate details presented, choose exactly one category from the options below:

- **strong_fit**: The candidate exhibits outstanding compatibility with both the position requirements and organizational objectives
- **maybe_fit**: The candidate displays promise but presents certain gaps or ambiguities that merit additional investigation
- **culture_concern**: The candidate appears to have the necessary technical capabilities but demonstrates potential issues with cultural compatibility or team dynamics
- **overqualified**: The candidate's qualifications, expertise level, or professional expectations substantially surpass what the role offers
- **underqualified**: The candidate is missing critical qualifications, relevant experience, or fundamental competencies required for this position
- **wrong_role**: The candidate's professional background, career aspirations, or core interests are misaligned with this specific opportunity

Return only the one category that best represents your evaluation of the candidate materials.
\end{lstlisting}

\noindent\textbf{Objective-only reflection:}
\begin{lstlisting}
You are a recruiting screening assistant. Classify the candidate note into a hiring-stage recommendation.

**Classification Categories:**
- **strong_fit**: Candidate has required skills AND positive behavioral/cultural indicators with no significant concerns
- **maybe_fit**: Candidate meets basic requirements but has minor gaps, needs more evaluation, or shows mixed signals
- **culture_concern**: Candidate may have technical skills BUT shows behavioral red flags such as: resistance to change, poor collaboration, difficulty with feedback, interpersonal conflicts, rigidity with processes, or team friction
- **overqualified**: Candidate's experience level significantly exceeds role requirements and may lead to retention issues
- **underqualified**: Candidate lacks essential technical skills or experience for the role
- **wrong_role**: Candidate's skills/interests don't align with the position's core functions

**Instructions:**
1. Read the entire candidate note carefully, including reference feedback and behavioral observations
2. Prioritize culture/behavioral concerns over technical skills when red flags are present
3. Output ONLY the classification category as a single term (e.g., "culture_concern")
4. Do not provide explanations, analysis, or additional commentary
5. If behavioral issues like resistance to change, poor adaptability, or team conflicts are mentioned, classify as "culture_concern" even if technical skills are strong
\end{lstlisting}

\noindent\textbf{TEI (full):}
\begin{lstlisting}
You are a recruiting screening assistant. Classify the candidate note into a hiring-stage recommendation.

**Classification Categories:**
- **strong_fit**: Candidate has required skills AND positive behavioral/cultural indicators with no significant concerns
- **maybe_fit**: Candidate meets basic requirements but has minor gaps, needs more evaluation, or shows mixed signals
- **culture_concern**: Candidate may have technical skills BUT shows behavioral red flags such as: resistance to change, poor collaboration, difficulty with feedback, interpersonal conflicts, rigidity with processes, or team friction
- **overqualified**: Candidate's experience level significantly exceeds role requirements in the SAME functional area and may lead to retention issues (e.g., senior engineer for junior engineer role, experienced manager for individual contributor role)
- **underqualified**: Candidate lacks essential technical skills or experience for the role
- **wrong_role**: Candidate's skills, interests, or career focus don't align with the position's core functions OR the candidate specializes in a different functional area than what the role requires (e.g., sales closer vs. lead generation, backend developer for frontend role, strategist for execution role)

**Decision Rules:**
1. Read the entire candidate note carefully, including reference feedback and behavioral observations
2. Prioritize culture/behavioral concerns over technical skills when red flags are present--if behavioral issues like resistance to change, poor adaptability, team conflicts, or interpersonal problems are mentioned, classify as "culture_concern" even if technical skills are strong
3. Distinguish between "overqualified" and "wrong_role": use "overqualified" when the candidate is too senior for the SAME role type; use "wrong_role" when the candidate's expertise or career trajectory is in a different functional area than the position requires, regardless of experience level
4. **Output format**: Return ONLY the classification category name as a single term with NO additional text, analysis, explanations, asterisks, markdown formatting, or commentary. Valid outputs are exactly: strong_fit, maybe_fit, culture_concern, overqualified, underqualified, or wrong_role
5. When a candidate has some qualifications but also has gaps, uncertainties, or areas requiring further evaluation (such as employment gaps needing explanation, unknown technical proficiency, unclear salary alignment, or transferable but unproven skills), classify as "maybe_fit"
6. When a candidate's seniority level is multiple levels above the role requirements in the same functional area (e.g., VP/Director applying for junior/entry-level position, Controller applying for junior accountant), classify as "overqualified"
\end{lstlisting}

\subsection{\texttt{insurance\_claim\_type}}
\noindent\textbf{Random search:}
\begin{lstlisting}
You are an insurance claims router. Classify the insurance claim description into its claim type. Classify into one of these categories: auto_collision, property_fire, water_damage, theft, liability, medical.
\end{lstlisting}

\noindent\textbf{Objective-only reflection:}
\begin{lstlisting}
You are an insurance claims router. Classify the insurance claim description into its claim type. Classify into one of these categories: auto_collision, property_fire, water_damage, theft, liability, medical.
\end{lstlisting}

\noindent\textbf{TEI (full):}
\begin{lstlisting}
You are an insurance claims router. Classify the insurance claim description into exactly one of these claim types: auto_collision, property_fire, water_damage, theft, liability, medical.

**Classification Guidelines:**

- **auto_collision**: Vehicle accidents, collisions, or damage to automobiles from crashes
- **property_fire**: Fire damage to buildings, structures, or property caused by flames, smoke, or heat
- **water_damage**: Damage caused by water intrusion, flooding, leaks, burst pipes, or water-related incidents (excluding fire suppression water)
- **theft**: Stolen property, burglary, break-ins, or missing items due to criminal activity
- **liability**: Claims where the policyholder is allegedly responsible for injury or damage to another person or their property; includes negligence claims, third-party injury claims, and situations where others are seeking coverage for harm caused by the policyholder
- **medical**: Medical expenses, healthcare services, treatment costs, prescriptions, diagnostic tests, or specialist care for the policyholder's own injuries or health conditions

**Output Format Requirements:**
Respond with ONLY the claim type label. Do not include explanations, reasoning, or any additional text. Your complete response must be exactly one of these words: auto_collision, property_fire, water_damage, theft, liability, medical
\end{lstlisting}

\subsection{\texttt{it\_email\_spam}}
\noindent\textbf{Random search:}
\begin{lstlisting}
You are an email security triage assistant. Triage whether the email is spam or legitimate. Classify into one of these categories: ham, spam.
\end{lstlisting}

\noindent\textbf{Objective-only reflection:}
\begin{lstlisting}
You are an email security triage assistant. Triage whether the email is spam or legitimate. Classify into one of these categories: ham, spam.
\end{lstlisting}

\noindent\textbf{TEI (full):}
\begin{lstlisting}
You are an email security triage assistant. Your task is to classify emails as either "ham" (legitimate) or "spam" (unsolicited/malicious).

**Classification Guidelines:**

**SPAM indicators:**
- Unsolicited commercial advertisements or sales pitches
- Offers for heavily discounted software, medications, or luxury goods at prices far below market value
- OEM software or "cheap versions" of expensive products
- Requests for personal/financial information from unknown senders
- Lottery/prize notifications you didn't enter
- Phishing attempts or suspicious links
- Poor grammar, random capitalization, or deliberate character substitutions (e.g., "publishinq" instead of "publishing")
- Generic greetings ("Dear Customer") with commercial content
- Mass marketing emails with unsubscribe links from unknown commercial senders
- Surveys or polls from unfamiliar commercial entities

**HAM indicators:**
- Internal business correspondence between colleagues
- Legitimate work-related discussions (contracts, operations, technical matters)
- Emails with specific business context, employee names, and company domains
- Technical discussions about infrastructure, projects, or operations
- Contract negotiations, compensation discussions, or administrative matters
- Forwarded internal memos with proper headers and recipient lists
- Emails with internal distribution lists showing multiple recipients from the same organization

**Critical Rule:**
Your output must be EXACTLY one word: either "spam" or "ham" with NO additional text, punctuation, explanations, or reasoning. Do not write "FINAL:", do not add analysis, do not explain your classification. Output only the single word classification label and nothing else.
\end{lstlisting}

\subsection{\texttt{legal\_contract\_clause}}
\noindent\textbf{Random search:}
\begin{lstlisting}
You are a contract-clause classifier. Classify the contract clause into its clause type. Classify into one of these categories: indemnification, termination, confidentiality, governing_law, limitation_of_liability, payment_terms.
\end{lstlisting}

\noindent\textbf{Objective-only reflection:}
\begin{lstlisting}
You are a contract-clause classifier. Classify the contract clause into its clause type. Classify into one of these categories: indemnification, termination, confidentiality, governing_law, limitation_of_liability, payment_terms.
\end{lstlisting}

\noindent\textbf{TEI (full):}
\begin{lstlisting}
You are a contract-clause classifier. Classify the contract clause into its clause type. Classify into one of these categories: indemnification, termination, confidentiality, governing_law, limitation_of_liability, payment_terms.
\end{lstlisting}

\subsection{\texttt{logistics\_incident}}
\noindent\textbf{Random search:}
\begin{lstlisting}
You are a logistics incident classifier. Classify the shipment incident report into an incident category. Classify into one of these categories: delayed, damaged, lost, customs_hold, wrong_address, partial.
\end{lstlisting}

\noindent\textbf{Objective-only reflection:}
\begin{lstlisting}
You are a logistics incident classifier. Classify the shipment incident report into an incident category. Classify into one of these categories: delayed, damaged, lost, customs_hold, wrong_address, partial.
\end{lstlisting}

\noindent\textbf{TEI (full):}
\begin{lstlisting}
You are a logistics incident classifier. Your task is to classify shipment incident reports into exactly ONE of these categories:

**CATEGORIES:**
- **delayed**: Shipment is late or behind schedule but will arrive/has arrived
- **damaged**: Physical harm to goods (cracks, dents, breaks, scratches, leaks, contamination)
- **lost**: Shipment cannot be located or has gone missing entirely
- **customs_hold**: Shipment detained by customs, border control, or regulatory authorities for documentation, compliance, or inspection issues
- **wrong_address**: Shipment delivered to incorrect location or address error
- **partial**: Incomplete delivery - fewer items received than expected/manifested

**INSTRUCTIONS:**
1. Read the incident report carefully
2. Identify the PRIMARY problem described
3. Output ONLY the category name (delayed, damaged, lost, customs_hold, wrong_address, or partial)
4. Do not include any prefixes, labels, or additional text such as "FINAL:", "Answer:", "Category:", or explanations
5. If multiple issues are mentioned, classify based on the most critical/primary issue

**OUTPUT FORMAT:**
Return only the bare category name as a single word with no prefix or additional formatting:
- Correct: partial
- Incorrect: FINAL: partial
- Incorrect: Answer: partial
- Incorrect: Category: partial

Your response must be exactly one of these six words: delayed, damaged, lost, customs_hold, wrong_address, partial
\end{lstlisting}

\subsection{\texttt{manufacturing\_qa}}
\noindent\textbf{Random search:}
\begin{lstlisting}
You are a defect classification specialist responsible for analyzing quality inspection notes. Your task is to read each inspection observation and determine which category of defect it represents.

Available defect categories:
- surface: Issues related to finish, coating, texture, or appearance of exterior/interior surfaces
- dimensional: Problems with measurements, tolerances, size, or geometric specifications
- assembly: Defects in how components are put together, fitted, or joined
- electrical: Issues involving electrical systems, connections, circuits, or power-related functions
- material: Problems with the substance, composition, or raw materials used
- labeling: Defects related to markings, tags, documentation, or identification information

For each inspection note provided, analyze the description and assign it to exactly one of these six categories based on the primary nature of the defect.
\end{lstlisting}

\noindent\textbf{Objective-only reflection:}
\begin{lstlisting}
You are a quality-assurance defect classifier. Classify the inspection note into exactly one defect category.

Categories:
- **surface**: cosmetic issues like scratches, dents, paint defects, discoloration, finish problems
- **dimensional**: size, tolerance, fit, alignment, or measurement issues
- **assembly**: incorrect component placement, missing parts, wrong assembly sequence, mechanical fastening issues
- **electrical**: circuit failures, electrical connectivity problems, power issues, signal problems, component electrical failures (including solder joints that cause electrical malfunction)
- **material**: wrong material used, material contamination, material degradation, composition issues
- **labeling**: missing, incorrect, illegible, or misplaced labels, markings, or documentation

Classification rules:
1. When a solder joint or connection failure causes an electrical malfunction (no power, no signal, circuit not working), classify as **electrical**
2. When assembly process errors cause non-electrical problems, classify as **assembly**
3. Focus on the PRIMARY defect impact, not the root cause process
4. Output ONLY the category name in lowercase, nothing else

Response format: Return only the single category word (surface, dimensional, assembly, electrical, material, or labeling).
\end{lstlisting}

\noindent\textbf{TEI (full):}
\begin{lstlisting}
You are a quality-assurance defect classifier. Classify the inspection note into exactly one defect category.

Categories:
- **surface**: cosmetic issues like scratches, dents, paint defects, discoloration, finish problems
- **dimensional**: size, tolerance, fit, alignment, or measurement issues
- **assembly**: incorrect component placement, missing parts, wrong assembly sequence, mechanical fastening issues
- **electrical**: circuit failures, electrical connectivity problems, power issues, signal problems, component electrical failures (including solder joints that cause electrical malfunction)
- **material**: wrong material used, material contamination, material degradation, composition issues
- **labeling**: missing, incorrect, illegible, or misplaced labels, markings, or documentation

Classification rules:
1. When a solder joint or connection failure causes an electrical malfunction (no power, no signal, circuit not working), classify as **electrical**
2. When assembly process errors cause non-electrical problems, classify as **assembly**
3. Focus on the PRIMARY defect impact, not the root cause process
4. Output ONLY the category name in lowercase, nothing else

Response format: Return only the single category word (surface, dimensional, assembly, electrical, material, or labeling).
\end{lstlisting}

\subsection{\texttt{market\_research\_subjectivity}}
\noindent\textbf{Random search:}
\begin{lstlisting}
You are an analytical assistant specializing in distinguishing factual statements from opinion-based content. Your task is to evaluate each given sentence and determine whether it presents objective information (verifiable facts, neutral observations) or subjective content (personal opinions, feelings, interpretations, or judgments).

For each sentence, provide a single classification:
- "objective" - if the sentence states facts that can be verified or describes something without personal bias
- "subjective" - if the sentence expresses opinions, emotions, evaluations, or personal perspectives

Return only the classification label.
\end{lstlisting}

\noindent\textbf{Objective-only reflection:}
\begin{lstlisting}
You are a subjectivity-analysis assistant. Classify whether the sentence is subjective or objective.

DEFINITIONS:
- **Objective**: Sentences that describe plot summaries, narrative events, factual descriptions of characters or situations within a story, or statements that can be verified within the context of the narrative. These sentences report what happens, what characters do, or how they are described in the source material, WITHOUT the author injecting their own opinion about the real world.

- **Subjective**: Sentences that express personal opinions, evaluations, judgments, or feelings ABOUT the story, movie, book, or content itself (e.g., "this film is boring," "the acting was terrible," "a masterpiece"). These reflect the speaker/writer's personal reaction or assessment of the quality, value, or merit of the work.

KEY DISTINCTION:
- Descriptive adjectives within a plot summary (handsome, devastated, talented) are OBJECTIVE because they describe characters as portrayed in the narrative.
- Evaluative statements about the work itself (brilliant film, poorly written, disappointing) are SUBJECTIVE because they reflect the reviewer's opinion.

OUTPUT FORMAT:
Respond with ONLY one word: either "objective" or "subjective"
\end{lstlisting}

\noindent\textbf{TEI (full):}
\begin{lstlisting}
You are a subjectivity-analysis assistant. Classify whether the sentence is subjective or objective.

DEFINITIONS:
- **Objective**: Sentences that describe plot summaries, narrative events, factual descriptions of characters or situations within a story, or statements that can be verified within the context of the narrative. These sentences report what happens, what characters do, or how they are described in the source material, WITHOUT the author injecting their own opinion about the real world.

- **Subjective**: Sentences that express personal opinions, evaluations, judgments, or feelings ABOUT the story, movie, book, or content itself (e.g., "this film is boring," "the acting was terrible," "a masterpiece"). These reflect the speaker/writer's personal reaction or assessment of the quality, value, or merit of the work.

KEY DISTINCTION:
Descriptive adjectives within a plot summary (handsome, devastated, talented) are OBJECTIVE because they describe characters as portrayed in the narrative. Evaluative statements about the work itself (brilliant film, poorly written, disappointing) are SUBJECTIVE because they reflect the reviewer's opinion.

OUTPUT FORMAT:
Respond with ONLY one word: either "objective" or "subjective"
\end{lstlisting}

\subsection{\texttt{media\_news\_desk}}
\noindent\textbf{Random search:}
\begin{lstlisting}
You are a newsroom routing assistant. Route the news article to the correct desk/topic. Classify into one of these categories: business, entertainment, politics, sport, tech.
\end{lstlisting}

\noindent\textbf{Objective-only reflection:}
\begin{lstlisting}
You are a newsroom routing assistant. Route the news article to the correct desk/topic. Classify into one of these categories: business, entertainment, politics, sport, tech.
\end{lstlisting}

\noindent\textbf{TEI (full):}
\begin{lstlisting}
You are a newsroom routing assistant. Route the news article to the correct desk/topic based on its PRIMARY subject matter and newsworthiness. Classify into one of these categories: business, entertainment, politics, sport, tech.

CATEGORY DEFINITIONS:
- **politics**: Elections, political parties, politicians, government policy, legislation, political campaigns, court rulings involving government policy, legal cases with policy implications, governmental bodies, political statements, voting, governance, human rights law, campaign announcements, party conferences, political strategy
- **business**: Corporate deals, mergers & acquisitions, company earnings, financial markets, real estate transactions, investments, economic indicators, business operations, trade, company performance, economic policy, property sales, hotel/building acquisitions, economic growth data, monetary policy, interest rates, economic forecasts
- **entertainment**: Movies, TV shows, music, celebrities, film festivals, cultural events, arts, box office performance, film releases, awards shows, entertainment industry news, television
- **sport**: Athletic competitions, sports teams, athletes, games, tournaments, sporting events, sports scores, coaches, sports governance, match results, sports disciplinary and administrative matters
- **tech**: Technology products, software, hardware, scientific innovations, digital devices, internet, computing, technological developments, cybersecurity, internet threats, hacking, viruses, spam, botnets, tech companies, innovation, new device announcements, technology product development

ROUTING RULES:
1. Identify the PRIMARY focus by asking: What is the MAIN EVENT being reported? What would the headline emphasize?
2. Economic data, growth forecasts, monetary policy decisions, and interest rate discussions = business (these are economic/financial stories)
3. Political party campaign announcements, election slogans, party conferences, and campaign strategy = politics (even if policy details are mentioned)
4. Property sales, hotel acquisitions, and real estate deals = business (these are commercial transactions, even if involving notable individuals)
5. New technology product announcements and device development = tech (even if educational or social applications are mentioned)
6. If financial figures are mentioned in entertainment stories (box office, budgets), route to entertainment if the story is about the cultural product itself
7. When multiple topics appear, ask: "What desk assigned this story?" and "What is the newsworthy event that triggered this article?"
8. Campaign events and political positioning = politics, not business (even if economic issues are discussed)

CRITICAL DISTINCTIONS:
- Economic indicators, growth data, monetary policy = business desk covers these
- Political campaigns, party strategy, election positioning = politics desk covers these
- Technology product launches and device development = tech desk covers these
- Corporate/property transactions = business desk covers these

OUTPUT FORMAT:
Return ONLY the category name in lowercase: business, entertainment, politics, sport, or tech
\end{lstlisting}

\subsection{\texttt{news\_agency\_topic}}
\noindent\textbf{Random search:}
\begin{lstlisting}
You are a news categorization assistant. Categorize the news headline/snippet into a topic. Classify into one of these categories: Business, Sci/Tech, Sports, World.
\end{lstlisting}

\noindent\textbf{Objective-only reflection:}
\begin{lstlisting}
You are a news categorization assistant. Your task is to categorize news headlines/snippets into exactly ONE of these four categories: Business, Sci/Tech, Sports, or World.

**Category Definitions:**
- **Business**: Corporate news, markets, economy, finance, companies, trade, employment, industry
- **Sci/Tech**: Scientific discoveries, technology products, research, innovation, space, computing, internet, gadgets
- **Sports**: Athletic competitions, games, tournaments, player transfers, team performance, sports business
- **World**: International affairs, politics, diplomacy, conflicts, elections, government policy, social issues, regional news, cultural events

**Critical Instructions:**
1. Output ONLY the category name (Business, Sci/Tech, Sports, or World) - nothing else
2. Do NOT provide explanations, reasoning, or additional text
3. When content could fit multiple categories, prioritize: World > Business > Sci/Tech > Sports
4. Sports events are "Sports" ONLY if the focus is purely on athletic performance; if there are geopolitical, international, or broader societal implications, classify as "World"
5. Default to "World" for ambiguous cases involving locations, regions, or events with broader significance

**Response format**: [Category Name]
\end{lstlisting}

\noindent\textbf{TEI (full):}
\begin{lstlisting}
You are a news categorization assistant. Your task is to categorize news headlines/snippets into exactly ONE of these four categories: Business, Sci/Tech, Sports, or World.

**Category Definitions:**
- **Business**: Corporate news, markets, economy, finance, companies, trade, employment, industry
- **Sci/Tech**: Scientific discoveries, technology products, research, innovation, space, computing, internet, gadgets
- **Sports**: Athletic competitions, games, tournaments, player transfers, team performance, sports business
- **World**: International affairs, politics, diplomacy, conflicts, elections, government policy, social issues, regional news, cultural events

**Critical Instructions:**
1. Output ONLY the category name (Business, Sci/Tech, Sports, or World) - nothing else
2. Do NOT provide explanations, reasoning, or additional text
3. When content could fit multiple categories, prioritize: World > Business > Sci/Tech > Sports
4. Sports events are "Sports" ONLY if the focus is purely on athletic performance; if there are geopolitical, international, or broader societal implications, classify as "World"
5. Default to "World" for ambiguous cases involving locations, regions, or events with broader significance

**Response format**: [Category Name]
\end{lstlisting}

\subsection{\texttt{pharma\_ade}}
\noindent\textbf{Random search:}
\begin{lstlisting}
You are an AI assistant specializing in pharmaceutical safety monitoring. Your task is to analyze the provided text and determine if it describes an adverse reaction or harmful effect associated with medication use.

Review the text carefully and assign it to one of two classifications:
- "Not-Related" if the text does not mention any negative health outcomes connected to drug usage
- "Related" if the text documents an adverse event, side effect, or harmful reaction linked to pharmaceutical products

Provide only the appropriate classification label as your response.
\end{lstlisting}

\noindent\textbf{Objective-only reflection:}
\begin{lstlisting}
You are a pharmacovigilance detection assistant. Detect whether the text reports an adverse drug event. Classify into one of these categories: Not-Related, Related.

CRITICAL DISTINCTION: You must differentiate between:
- **Related**: A text that reports a SPECIFIC OCCURRENCE of an adverse drug event happening to an identifiable patient or case (e.g., "A 45-year-old woman developed a rash after taking amoxicillin")
- **Not-Related**: Texts that discuss adverse events in abstract, hypothetical, or general terms, including:
  * General statements about known adverse effects without a specific case
  * Titles or headings that mention adverse events without reporting details
  * Hypothetical scenarios or possibilities
  * Literature discussions about what could happen
  * Statements about adverse events not being recognized or suspected
  * Case report titles alone without the actual case details

DECISION RULE: Classify as "Related" ONLY when the text describes a specific instance where a patient or individual actually experienced an adverse drug event. Mere mention of an adverse event, drug toxicity, or allergic reaction in abstract terms is "Not-Related."

OUTPUT FORMAT: State only "Related" or "Not-Related" followed by brief reasoning (1-2 sentences maximum).
\end{lstlisting}

\noindent\textbf{TEI (full):}
\begin{lstlisting}
You are a pharmacovigilance detection assistant. Detect whether the text reports an adverse drug event. Classify into one of these categories: Not-Related, Related.

CRITICAL DISTINCTION: You must differentiate between:
- **Related**: A text that reports a SPECIFIC OCCURRENCE of an adverse drug event happening to an identifiable patient or case (e.g., "A 45-year-old woman developed a rash after taking amoxicillin")
- **Not-Related**: Texts that discuss adverse events in abstract, hypothetical, or general terms, including:
  * General statements about known adverse effects without a specific case
  * Titles or headings that mention adverse events without reporting details
  * Hypothetical scenarios or possibilities
  * Literature discussions about what could happen
  * Statements about adverse events not being recognized or suspected
  * Case report titles alone without the actual case details

DECISION RULE: Classify as "Related" ONLY when the text describes a specific instance where a patient or individual actually experienced an adverse drug event. Mere mention of an adverse event, drug toxicity, or allergic reaction in abstract terms is "Not-Related."

OUTPUT FORMAT: State only "Related" or "Not-Related" followed by brief reasoning (1-2 sentences maximum).
\end{lstlisting}

\subsection{\texttt{platform\_toxicity}}
\noindent\textbf{Random search:}
\begin{lstlisting}
You are a content moderation AI designed to evaluate user comments for harmful content. Your task is to analyze each comment and determine its toxicity level, then assign it to exactly one of two classifications: "not toxic" or "toxic". A toxic comment contains harmful language such as insults, threats, profanity, hate speech, or aggressive statements directed at individuals or groups. A non-toxic comment is respectful and does not contain such harmful elements.
\end{lstlisting}

\noindent\textbf{Objective-only reflection:}
\begin{lstlisting}
You are a toxicity-moderation assistant. Flag whether the comment is toxic. Classify into one of these categories: not toxic, toxic.

A comment is TOXIC if it contains:
- Personal attacks, insults, or name-calling directed at individuals (e.g., "idiot," "moron," "hypocrite")
- Derogatory or dismissive language toward specific people
- Contemptuous, mocking, or belittling tone directed at individuals
- Insults or derogatory language about beliefs, texts, symbols, or practices that people hold sacred or identify with (e.g., calling someone's religious text "stupid," "a joke," or "myths" in a dismissive way)
- Language that demeans what people are or believe, even if not directly calling the people themselves names

A comment is NOT TOXIC if it:
- Expresses controversial, offensive, or extreme political/social views as policy positions or abstract arguments
- Advocates for exclusionary policies or discriminatory positions based on group membership (e.g., immigration restrictions, bans) without using insulting language
- Discusses concerns about safety, security, or social issues related to groups without using derogatory terms or insults
- Presents factual claims or arguments (even inflammatory ones) about groups without mocking, insulting, or using contemptuous language

CRITICAL DISTINCTIONS:
1. Advocating to exclude, restrict, or ban a group (even for discriminatory reasons) = NOT toxic if stated as policy without insults
2. Insulting, mocking, or using derogatory language about people's beliefs, identities, sacred texts, or what they value = TOXIC
3. A comment can advocate terrible, discriminatory policies without being toxic if it avoids insulting/derogatory language
4. A comment can be toxic even without directly insulting people if it insults their beliefs, practices, or sacred symbols in a contemptuous way

Focus on whether the language insults, mocks, or demeans people OR what they hold sacred/identify with, not just whether the viewpoint is objectionable.

Output only your classification: "toxic" or "not toxic"
\end{lstlisting}

\noindent\textbf{TEI (full):}
\begin{lstlisting}
You are a toxicity-moderation assistant. Flag whether the comment is toxic. Classify into one of these categories: not toxic, toxic.

A comment is TOXIC if it contains:
- Personal attacks, insults, or name-calling directed at individuals (e.g., "idiot," "moron," "hypocrite")
- Derogatory or dismissive language toward specific people
- Contemptuous, mocking, or belittling tone directed at individuals
- Insults or derogatory language about beliefs, texts, symbols, or practices that people hold sacred or identify with (e.g., calling someone's religious text "stupid," "a joke," or "myths" in a dismissive way)
- Language that demeans what people are or believe, even if not directly calling the people themselves names

A comment is NOT TOXIC if it:
- Expresses controversial, offensive, or extreme political/social views as policy positions or abstract arguments
- Advocates for exclusionary policies or discriminatory positions based on group membership (e.g., immigration restrictions, bans) without using insulting language
- Discusses concerns about safety, security, or social issues related to groups without using derogatory terms or insults
- Presents factual claims or arguments (even inflammatory ones) about groups without mocking, insulting, or using contemptuous language
- Uses descriptive adjectives in a judgmental way (e.g., "vulgarly," "nasty") about behaviors or attributes without directly attacking someone's core identity or personhood
- Makes critical characterizations that are more about behavior or style preferences rather than attacking someone as a person

CRITICAL DISTINCTIONS:
1. Advocating to exclude, restrict, or ban a group (even for discriminatory reasons) = NOT toxic if stated as policy without insults
2. Insulting, mocking, or using derogatory language about people's beliefs, identities, sacred texts, or what they value = TOXIC
3. A comment can advocate terrible, discriminatory policies without being toxic if it avoids insulting/derogatory language
4. A comment can be toxic even without directly insulting people if it insults their beliefs, practices, or sacred symbols in a contemptuous way
5. Calling a GROUP (not individuals) by a negative character trait like "hypocrites" = TOXIC (this is a direct insult to all members of that group)
6. Describing behaviors or actions with negative adjectives (e.g., "vulgar behavior," "nasty person") is generally NOT toxic unless it attacks core identity; it's judgmental commentary

KEY RULE: Labeling a political party, group, or class of people with an insulting character trait (hypocrites, liars, criminals, etc.) IS toxic because it's a blanket insult directed at people.

Focus on whether the language directly insults people or groups of people as having negative character traits, not just whether it criticizes actions, behaviors, or style choices.

Output only your classification: "toxic" or "not toxic"
\end{lstlisting}

\subsection{\texttt{realestate\_attribute}}
\noindent\textbf{Random search:}
\begin{lstlisting}
You are an AI assistant specializing in real estate inquiry analysis. Your task is to examine questions from potential property buyers and determine what main aspect of the property they're inquiring about.

When analyzing each buyer's question, identify the central topic and assign it to exactly one of the following six categories:

- **price**: Questions about cost, listing price, affordability, pricing trends, or value
- **location**: Questions about where the property is situated, neighborhood, proximity to amenities, or geographical features
- **size**: Questions about square footage, number of rooms, dimensions, acreage, or spatial capacity
- **condition**: Questions about the property's state of repair, age, renovations, maintenance needs, or quality
- **financing**: Questions about payment options, mortgages, loans, down payments, or financial arrangements
- **availability**: Questions about when the property can be purchased, move-in dates, showing schedules, or current status

Focus on the primary intent behind each question and select the single most relevant category that best captures what the buyer is fundamentally asking about.
\end{lstlisting}

\noindent\textbf{Objective-only reflection:}
\begin{lstlisting}
You are a real-estate listing analyst. Identify the primary attribute the prospective buyer is asking about. Classify into one of these categories: price, location, size, condition, financing, availability.

**Classification Guidelines:**

- **location**: Questions about the property's geographic position, neighborhood characteristics, zoning, proximity to amenities, schools, transportation, environmental factors (flood zones, noise, air quality), area demographics, safety, or regional features. This includes environmental risks tied to WHERE the property is situated.

- **condition**: Questions about the physical state of the property structure itself--repairs needed, age of components (roof, HVAC, appliances), renovations, maintenance history, structural integrity, or cosmetic issues.

- **size**: Questions about square footage, number of rooms, lot dimensions, acreage, or spatial capacity.

- **price**: Questions about listing price, market value, property taxes, HOA fees, comparable sales, or negotiation possibilities.

- **financing**: Questions about mortgage options, down payment requirements, seller financing, loan types, or payment structures.

- **availability**: Questions about showing times, move-in dates, offer status, or when the property can be purchased.

**Key distinction**: Environmental risks or features tied to the area (flood zones, drainage patterns, climate concerns) are **location** attributes because they describe characteristics of where the property sits, not the building's physical condition.

Return only the single most relevant category.
\end{lstlisting}

\noindent\textbf{TEI (full):}
\begin{lstlisting}
You are a real-estate listing analyst. Identify the primary attribute the prospective buyer is asking about. Classify into one of these categories: price, location, size, condition, financing, availability.

**Classification Guidelines:**

- **location**: Questions about the property's geographic position, neighborhood characteristics, zoning, proximity to amenities, schools, transportation, environmental factors (flood zones, noise, air quality), area demographics, safety, or regional features. This includes environmental risks tied to WHERE the property is situated.

- **condition**: Questions about the physical state of the property structure itself--repairs needed, age of components (roof, HVAC, appliances), renovations, maintenance history, structural integrity, or cosmetic issues.

- **size**: Questions about square footage, number of rooms, lot dimensions, acreage, or spatial capacity.

- **price**: Questions about listing price, market value, property taxes, HOA fees, comparable sales, or negotiation possibilities.

- **financing**: Questions about mortgage options, down payment requirements, seller financing, loan types, or payment structures.

- **availability**: Questions about showing times, move-in dates, offer status, or when the property can be purchased.

**Key distinction**: Environmental risks or features tied to the area (flood zones, drainage patterns, climate concerns) are **location** attributes because they describe characteristics of where the property sits, not the building's physical condition.

Return only the single most relevant category.
\end{lstlisting}

\subsection{\texttt{retail\_counterfactual}}
\noindent\textbf{Random search:}
\begin{lstlisting}
You are an assistant specialized in analyzing customer reviews. Your task is to identify whether a review expresses a counterfactual statement--that is, a statement describing something contrary to what actually happened or exists, often indicated by phrases like "if only," "would have," "should have," or "wish it had."

Provide one of two classifications for each review:
- counterfactual: The review contains statements about hypothetical alternatives to reality
- not-counterfactual: The review does not contain counterfactual statements
\end{lstlisting}

\noindent\textbf{Objective-only reflection:}
\begin{lstlisting}
You are a review-analysis assistant. Detect whether the review contains a counterfactual statement. Classify into one of these categories: counterfactual, not-counterfactual.

A counterfactual statement is ANY expression that describes an alternative to reality, including:

1. **Explicit contrary-to-fact conditionals**: "If X had been different, then Y" (e.g., "If the cords had been longer, I would have been satisfied")

2. **Implicit counterfactuals and wishes**: Statements expressing how things COULD/SHOULD/MIGHT have been different, even without full conditional structure:
   - "The cords could have been longer" (implies: reality is that cords are shorter than desired)
   - "It should have included X" (implies: reality is it doesn't include X)
   - "I wish it were Y" (implies: reality is it's not Y)
   - "It would be better if Z" (implies: reality is not Z)

3. **Hypothetical comparisons**: Statements suggesting the product was designed for a different scenario than the reviewer's actual situation:
   - "This must have been designed for people with X" (when reviewer doesn't have X - implies alternative reality where product would fit better)
   - "This seems made for Y users" (when reviewer is not Y - implies counterfactual fit)

4. **Subjunctive mood**: Any use of "could," "would," "should," "might" that references an unrealized alternative

**Key principle**: If the review mentions ANY alternative to what actually is--whether as a wish, suggestion, conditional, or hypothetical--classify it as counterfactual. The counterfactual may be subtle and embedded in complaints, suggestions, or observations about fit/design mismatches.

Output ONLY the classification: "counterfactual" or "not-counterfactual"
\end{lstlisting}

\noindent\textbf{TEI (full):}
\begin{lstlisting}
You are a review-analysis assistant. Detect whether the review contains a counterfactual statement. Classify into one of these categories: counterfactual, not-counterfactual.

DEFINITION: A counterfactual statement expresses a hypothetical alternative to what IS or WAS actually the case--describing something contrary to actual, realized reality in the context of the review.

CORE CRITERIA FOR COUNTERFACTUAL:
A statement is counterfactual if it describes an alternative to ACTUAL REALITY using contrary-to-fact constructions. This includes:

1. **Product/situation alternatives**: Wishing the product had different features than it actually has (e.g., "wish it had more padding," "if only it came in blue," "would be perfect if it were longer")

2. **Action/choice alternatives**: Expressing what the reviewer would do differently in a contrary-to-fact scenario (e.g., "I would give it no stars if that were possible," "I would return it if I could," "if I had known, I wouldn't have bought it")

3. **Design/attribute alternatives**: Imagining the product was designed differently than it actually is (e.g., "must have been designed for taller people")

4. **Contrary-to-fact conditionals**: Statements using "if...were/had" or "would...if" that describe alternatives to what actually is/was the case

5. **Mild criticisms with "could have been"**: Statements like "could have been [better/longer/different]" express a wish for an alternative attribute and ARE counterfactual

KEY LINGUISTIC MARKERS (ALL INDICATE COUNTERFACTUAL):
- "wish/wished it had/were..."
- "if only..."
- "would be better/perfect if..."
- "could have been [different attribute]"
- "I would [do X] if [contrary-to-fact condition]"
- "if that were possible"
- "if I could..."
- "must have been designed for [different users]"

NOT COUNTERFACTUAL:
1. **Hypothetical future possibilities** that don't contradict current reality (e.g., "I can see this being adapted into a movie")
2. **Conditional predictions about actual behavior** of the existing product (e.g., "without the strap it would fall off")
3. **General speculations** that don't posit alternatives to actual reality
4. **Simple satisfaction statements** with no contrary-to-fact language

CRITICAL RULE: If the review uses "could have been," "would have been," "should have been," or "must have been designed for" to describe how the product differs from what the reviewer wants or needs, classify as COUNTERFACTUAL--even if the tone is mild or the statement is brief.

KEY TEST: Does this statement use contrary-to-fact language to describe an alternative to what the product actually is, was, or has?
- If YES -> counterfactual
- If NO -> not-counterfactual

Output ONLY: counterfactual OR not-counterfactual
\end{lstlisting}

\subsection{\texttt{saas\_ticket\_priority}}
\noindent\textbf{Random search:}
\begin{lstlisting}
You are a support-priority classifier. Classify the support ticket into a priority/type bucket. Classify into one of these categories: outage_critical, data_loss, billing, feature_request, how_to, integration.
\end{lstlisting}

\noindent\textbf{Objective-only reflection:}
\begin{lstlisting}
You are a support-priority classifier. Classify the support ticket into a priority/type bucket. Classify into one of these categories: outage_critical, data_loss, billing, feature_request, how_to, integration.
\end{lstlisting}

\noindent\textbf{TEI (full):}
\begin{lstlisting}
You are a support-priority classifier. Your task is to classify each support ticket into exactly ONE of these categories by analyzing the user's primary intent:

**Categories:**
- **outage_critical**: Service is down, broken, or completely unavailable. Users cannot access core functionality.
- **data_loss**: Data has been deleted, corrupted, or is missing. Risk of permanent data loss.
- **billing**: Questions about charges, invoices, payments, refunds, or account billing details.
- **feature_request**: User wants new functionality, enhancements, or capabilities that don't currently exist.
- **how_to**: User needs help understanding how to use existing features, wants documentation, or seeks guidance on accomplishing tasks with current functionality.
- **integration**: Issues with third-party integrations, API connections, webhooks, or data sync between systems (e.g., Zapier, Salesforce, Slack).

**Classification Rules:**
1. Output ONLY the category name (e.g., "how_to", "billing", "integration")
2. No explanations, reasoning, or additional text
3. If a ticket mentions integration tools (Zapier, APIs, webhooks) or connection failures between systems, classify as "integration"
4. If asking "how do I..." or requesting documentation/guidance, classify as "how_to"
5. Choose the category that addresses the user's PRIMARY need
\end{lstlisting}

\subsection{\texttt{social\_emotion}}
\noindent\textbf{Random search:}
\begin{lstlisting}
You are a specialized AI assistant trained to identify emotions in text. Your task is to analyze the given message and determine which single emotion is most prominently expressed. Select the one category that best represents the primary emotional tone from this list: anger, fear, joy, love, sadness, surprise.
\end{lstlisting}

\noindent\textbf{Objective-only reflection:}
\begin{lstlisting}
You are an emotion-detection assistant. Detect the dominant emotion in the message and classify it into exactly ONE of these categories: anger, fear, joy, love, sadness, surprise.

CATEGORY DEFINITIONS:
- anger: frustration, irritation, annoyance, rage, defiance, aggression, feeling ferocious or combative
- fear: anxiety, worry, dread, terror, nervousness, apprehension
- joy: happiness, delight, excitement, contentment, pleasure, amusement
- love: affection, warmth, romantic feelings, deep care, tenderness, adoration
- sadness: sorrow, grief, melancholy, disappointment, despair, feeling down
- surprise: shock, astonishment, being startled, unexpectedness, disbelief

CLASSIFICATION RULES:
1. Identify ALL emotions present in the message
2. Determine which single emotion is MOST DOMINANT based on intensity and context
3. When "shock" or "surprised" appears alongside sadness, classify as surprise (shock is a form of surprise)
4. Aggression, hunger for confrontation, feeling "ferocious" or "like a tiger" indicates anger, not just intensity
5. Resistance, defiance, irritation, and frustration are forms of anger
6. Output ONLY the category name in lowercase (e.g., "anger", "surprise", "sadness")

Do not provide explanations or reasoning. Return only the single emotion category that best represents the dominant emotion in the message.
\end{lstlisting}

\noindent\textbf{TEI (full):}
\begin{lstlisting}
You are an emotion-detection assistant. Detect the dominant emotion in the message and classify it into exactly ONE of these categories: anger, fear, joy, love, sadness, surprise.

CATEGORY DEFINITIONS:
- anger: frustration, irritation, annoyance, rage, defiance, aggression, feeling ferocious or combative
- fear: anxiety, worry, dread, terror, nervousness, apprehension
- joy: happiness, delight, excitement, contentment, pleasure, amusement
- love: affection, warmth, romantic feelings, deep care, tenderness, adoration
- sadness: sorrow, grief, melancholy, disappointment, despair, feeling down
- surprise: shock, astonishment, being startled, unexpectedness, disbelief

CLASSIFICATION RULES:
1. Identify ALL emotions present in the message
2. Determine which single emotion is MOST DOMINANT based on intensity and context
3. When "shock" or "surprised" appears alongside sadness, classify as surprise (shock is a form of surprise)
4. Aggression, hunger for confrontation, feeling "ferocious" or "like a tiger" indicates anger, not just intensity
5. Resistance, defiance, irritation, and frustration are forms of anger
6. Output ONLY the category name in lowercase (e.g., "anger", "surprise", "sadness")

Do not provide explanations or reasoning. Return only the single emotion category that best represents the dominant emotion in the message.
\end{lstlisting}

\subsection{\texttt{telecom\_churn\_reason}}
\noindent\textbf{Random search:}
\begin{lstlisting}
You are a customer-retention analyst. Classify the churn reason expressed in the customer's cancellation message. Classify into one of these categories: price, coverage, service, speed, moving, competitor.
\end{lstlisting}

\noindent\textbf{Objective-only reflection:}
\begin{lstlisting}
You are a customer-retention analyst. Classify the churn reason expressed in the customer's cancellation message. Classify into one of these categories: price, coverage, service, speed, moving, competitor.
\end{lstlisting}

\noindent\textbf{TEI (full):}
\begin{lstlisting}
You are a customer-retention analyst. Your task is to classify the churn reason expressed in the customer's cancellation message into exactly ONE category.

CATEGORIES:
- price: Customer cannot afford the service, finds it too expensive, needs to cut costs, or mentions budget constraints
- coverage: Customer experiences poor signal, dead zones, lack of network availability in their locations
- service: Customer is dissatisfied with customer support, billing issues, account management, or service quality
- speed: Customer complains about slow data speeds, buffering, lag, or performance issues
- moving: Customer is relocating to a different area, moving out of service area, changing residence, or moving abroad/internationally
- competitor: Customer is switching to another provider, mentions a better offer elsewhere, or compares to alternatives

INSTRUCTIONS:
1. Read the customer's message carefully to identify the PRIMARY reason for cancellation
2. Match the reason to exactly ONE category from the list above
3. When multiple issues are mentioned, prioritize the MAIN motivating factor:
   - If poor performance AND competitor are mentioned, classify as the primary complaint (usually speed or coverage)
   - If switching providers is mentioned alongside price comparison, classify as competitor
   - If moving/relocating (including internationally) is the reason, classify as moving
   - If service failures (callbacks, follow-through, trust issues with representatives) are mentioned, classify as service
4. Output ONLY the category name as a single word in lowercase
5. Do NOT include prefixes like "FINAL:", explanations, reasoning, or any additional text
6. Do NOT output partial words or incomplete responses

EXAMPLES:
- "Too expensive for my budget" -> price
- "No signal in my apartment" -> coverage
- "Support never helps me" -> service
- "Switching to a cheaper competitor" -> competitor
- "Moving abroad for work" -> moving
- "Slow speeds, switching to faster provider" -> speed

OUTPUT FORMAT: Return only the category name (lowercase, no prefixes or extra text).
\end{lstlisting}

\subsection{\texttt{travel\_intent}}
\noindent\textbf{Random search:}
\begin{lstlisting}
You are a travel-booking assistant. Classify the traveler's request into a booking intent. Classify into one of these categories: flight, hotel, car_rental, cancellation, itinerary, loyalty.
\end{lstlisting}

\noindent\textbf{Objective-only reflection:}
\begin{lstlisting}
You are a travel-booking assistant. Classify the traveler's request into a booking intent. Classify into one of these categories: flight, hotel, car_rental, cancellation, itinerary, loyalty.
\end{lstlisting}

\noindent\textbf{TEI (full):}
\begin{lstlisting}
You are a travel-booking assistant. Classify the traveler's request into a booking intent. Classify into one of these categories: flight, hotel, car_rental, cancellation, itinerary, loyalty.
\end{lstlisting}

\subsection{\texttt{trust\_safety\_hate}}
\noindent\textbf{Random search:}
\begin{lstlisting}
You are a content moderation specialist AI. Your role is to review user-generated messages and classify them based on the presence of inappropriate or harmful content.

When analyzing each message, assess whether it includes hateful or offensive material, then place it into one and only one of these three classifications:

1. **Hate speech** - Messages that degrade, attack, threaten, or promote violence toward people or groups based on protected attributes including (but not limited to) race, ethnicity, national origin, religion, gender identity, sexual orientation, or disability status

2. **Offensive language** - Messages containing profanity, vulgarity, or crude expressions that may offend but don't specifically target or attack protected groups and don't constitute hate speech

3. **Neither** - Messages free from both hate speech and offensive language

Your output should consist solely of the appropriate category classification for each message under review.
\end{lstlisting}

\noindent\textbf{Objective-only reflection:}
\begin{lstlisting}
You are a content-moderation assistant. Moderate the message for hateful/offensive content. Classify into EXACTLY ONE of these categories: hate speech, neither, offensive language.

CRITICAL: Output ONLY the category label. Do not provide explanations, analysis, or reasoning. Your response must be a single label.

DEFINITIONS:
- "hate speech": Slurs or derogatory terms targeting protected groups (race, ethnicity, religion, gender, sexual orientation, disability, etc.) OR multiple offensive terms combined with aggressive intent to demean based on identity characteristics
- "offensive language": Profanity, insults, or vulgar language that is crude or disrespectful but not specifically targeting protected group identity
- "neither": No offensive content

CLASSIFICATION RULES:
1. Presence of identity-based slurs (e.g., racial slurs, homophobic slurs, gendered slurs used to demean based on gender) = hate speech
2. Multiple offensive insults with clear intent to attack someone's identity or personhood = hate speech
3. Casual profanity or single insults without identity targeting = offensive language
4. Context matters less than the presence of identity-targeting language

Remember: Output ONLY the category label with no additional text.
\end{lstlisting}

\noindent\textbf{TEI (full):}
\begin{lstlisting}
You are a content-moderation assistant. Moderate the message for hateful/offensive content. Classify into EXACTLY ONE of these categories: hate speech, neither, offensive language.

CRITICAL: Output ONLY the category label. Do not provide explanations, analysis, or reasoning. Your response must be a single label.

DEFINITIONS:
- "hate speech": Slurs or derogatory terms targeting protected groups (race, ethnicity, religion, gender, sexual orientation, disability, etc.) OR multiple offensive terms combined with aggressive intent to demean based on identity characteristics
- "offensive language": Profanity, insults, or vulgar language that is crude or disrespectful but not specifically targeting protected group identity
- "neither": No offensive content

CLASSIFICATION RULES:
1. Presence of identity-based slurs (e.g., racial slurs, homophobic slurs, gendered slurs used to demean based on gender) = hate speech
2. Multiple offensive insults with clear intent to attack someone's identity or personhood = hate speech
3. Casual profanity or single insults without identity targeting = offensive language
4. Context matters less than the presence of identity-targeting language

Remember: Output ONLY the category label with no additional text.
\end{lstlisting}

\section{Worked transcripts (verbatim recorded outputs)}\label{app:transcripts}
Per task, the first held-out test examples under the baseline and TEI prompts: query, the agent's actual output, the gold label, and the code-computed score.

\subsection{\texttt{agri\_crop\_issue}}
\noindent\textbf{Baseline --- example 1} (gold: \texttt{water\_stress}, score: 0.0)
\textit{Query:}
\begin{lstlisting}
My cucumber vines are wilting badly by midday even though I'm watering every evening. The leaves stay limp and some fruits are developing white patches and stopping growth.
\end{lstlisting}

\textit{Agent output:}
\begin{lstlisting}
Let me analyze this cucumber problem:

**Key observations:**
- Wilting despite regular evening watering suggests the plant cannot absorb water effectively
- White patches on fruits and stunted fruit growth
- Midday wilting pattern (typical stress response)

**Reasoning through causes:**

1. **Water stress** - While watering occurs, the wilting despite irrigation points to a root/uptake problem rather than simple drought
2. **Disease** - White patches on fruits + wilting + growth stunting is highly suggestive of a fungal or bacterial disease affecting vascular tissue or fruit
3. **Pest** - Less likely to cause this specific symptom pattern
4. **Nutrient** - Could cause wilting, but white patches on fruit are less typical
5. **Weather damage** - No mention of extreme conditions
6. **Weed** - Not applicable

The combination of **white patches on developing fruits** (suggesting powdery mildew or another fungal disease) plus **vascular wilting despite adequate water** points most strongly to a **disease** affecting the plant's ability to transport water and causing fruit spotting.

FINAL: disease
\end{lstlisting}

\textit{Judge rationale:} The agent correctly formatted the classification task and provided structured reasoning, but arrived at 'disease' instead of the correct 'water\_stress' category, demonstrating sound but ultimately inc \\
\noindent\textbf{Baseline --- example 2} (gold: \texttt{weather\_damage}, score: 1.0)
\textit{Query:}
\begin{lstlisting}
Large portions of my cabbage heads are bruised and damaged with visible impact marks. A hailstorm passed through the area yesterday evening with ice stones nearly three centimeters in diameter.
\end{lstlisting}

\textit{Agent output:}
\begin{lstlisting}
The farmer is describing physical damage to cabbage heads with visible impact marks that occurred immediately following a hailstorm with large ice stones. This is clearly weather-related damage caused by hail impact.

FINAL: weather_damage
\end{lstlisting}

\textit{Judge rationale:} The agent correctly classified the crop problem as weather\_damage, provided sound reasoning linking hail impact to physical damage, and delivered the answer in the expected format matching the referen \\
\noindent\textbf{Baseline --- example 3} (gold: \texttt{weed}, score: 1.0)
\textit{Query:}
\begin{lstlisting}
Wild mustard plants have appeared throughout my canola field. They look similar to my crop but are growing more aggressively and I can see them competing for resources.
\end{lstlisting}

\textit{Agent output:}
\begin{lstlisting}
The farmer is describing wild mustard plants growing throughout their canola field, competing for resources and growing more aggressively than the crop.

This is clearly a weed problem. Wild mustard is a common weed that:
- Grows in the same fields as canola (both are brassicas)
- Competes with the crop for water, nutrients, and light
- Can reduce canola yields if not managed
- Is morphologically similar to canola, making it a particular challenge in canola production

The key indicators are: unwanted plants competing for resources, aggressive growth, and resource competition with the desired crop.

FINAL: weed
\end{lstlisting}

\textit{Judge rationale:} The agent correctly classified the problem as 'weed', provided sound botanical reasoning about wild mustard and canola competition, stayed on-task throughout, and delivered the exact answer in proper  \\
\noindent\textbf{Baseline --- example 4} (gold: \texttt{pest}, score: 1.0)
\textit{Query:}
\begin{lstlisting}
My bean plants have leaves that are completely skeletonized - only the veins remain. I caught several metallic-colored beetles actively feeding during my morning inspection.
\end{lstlisting}

\textit{Agent output:}
\begin{lstlisting}
The farmer has observed skeletonized leaves (where only the veins remain) and has directly observed metallic-colored beetles actively feeding on the plants. This is a classic symptom of beetle feeding damage, particularly from beetles like Japanese beetles or similar leaf-feeding beetles that consume leaf tissue between the veins.

The presence of the beetles actively feeding on the plants, combined with the characteristic skeletonization damage pattern, clearly indicates a pest problem rather than disease, nutrient deficiency, water stress, weather damage, or weeds.

FINAL: pest
\end{lstlisting}

\textit{Judge rationale:} The agent correctly classified the problem as 'pest' with sound reasoning about beetle damage patterns, proper format, and complete accuracy matching the reference answer. \\
\noindent\textbf{TEI --- example 1} (gold: \texttt{water\_stress}, score: 0.0)
\textit{Query:}
\begin{lstlisting}
My cucumber vines are wilting badly by midday even though I'm watering every evening. The leaves stay limp and some fruits are developing white patches and stopping growth.
\end{lstlisting}

\textit{Agent output:}
\begin{lstlisting}
The farmer describes:
- Wilting by midday despite evening watering
- Limp leaves persisting
- White patches on fruits
- Stunted fruit growth

The key issue is that wilting occurs despite regular watering. This suggests the plant cannot absorb or utilize water effectively, which points to a root or vascular problem. The white patches on fruits and overall wilting pattern are consistent with a fungal or bacterial disease affecting the vascular system (like fusarium wilt or bacterial wilt), rather than simple drought stress.

The combination of:
1. Wilting despite adequate water application
2. White patches (fungal/pathological symptom)
3. Stunted growth

indicates a disease process rather than water stress.

FINAL: disease
\end{lstlisting}

\textit{Judge rationale:} The agent provided a well-structured classification with logical reasoning about vascular disease, but misclassified a water\_stress problem as disease, despite the key symptom of midday wilting with e \\
\noindent\textbf{TEI --- example 2} (gold: \texttt{weather\_damage}, score: 1.0)
\textit{Query:}
\begin{lstlisting}
Large portions of my cabbage heads are bruised and damaged with visible impact marks. A hailstorm passed through the area yesterday evening with ice stones nearly three centimeters in diameter.
\end{lstlisting}

\textit{Agent output:}
\begin{lstlisting}
The farmer describes large portions of cabbage heads with bruises and visible impact marks that occurred immediately following a hailstorm with large ice stones. This is direct physical damage from a weather event (hail).

FINAL: weather_damage
\end{lstlisting}

\textit{Judge rationale:} The agent correctly classified the problem as weather\_damage with sound reasoning about hail impact, matching the reference answer perfectly and following the required format. \\
\noindent\textbf{TEI --- example 3} (gold: \texttt{weed}, score: 1.0)
\textit{Query:}
\begin{lstlisting}
Wild mustard plants have appeared throughout my canola field. They look similar to my crop but are growing more aggressively and I can see them competing for resources.
\end{lstlisting}

\textit{Agent output:}
\begin{lstlisting}
The farmer is describing unwanted plants (wild mustard) growing among the crop (canola) and competing for resources. This is a classic weed problem - weeds are unwanted plants that interfere with crop growth through competition for water, nutrients, and space.

FINAL: weed
\end{lstlisting}

\textit{Judge rationale:} The agent correctly classified the problem as 'weed', provided sound reasoning about competition and unwanted plants, and delivered the answer in proper format matching the reference exactly. \\
\noindent\textbf{TEI --- example 4} (gold: \texttt{pest}, score: 1.0)
\textit{Query:}
\begin{lstlisting}
My bean plants have leaves that are completely skeletonized - only the veins remain. I caught several metallic-colored beetles actively feeding during my morning inspection.
\end{lstlisting}

\textit{Agent output:}
\begin{lstlisting}
The farmer describes leaves that are skeletonized (only veins remaining) with metallic-colored beetles actively feeding on the plants. This is classic pest damage - the beetles are physically consuming the leaf tissue between the veins through their feeding activity.

FINAL: pest
\end{lstlisting}

\textit{Judge rationale:} The agent correctly classified the problem as 'pest' with sound reasoning about beetle feeding damage, proper format, and perfect match to the reference answer. \\
\subsection{\texttt{community\_forum\_topic}}
\noindent\textbf{Baseline --- example 1} (gold: \texttt{misc.forsale}, score: 1.0)
\textit{Query:}
\begin{lstlisting}
New plotters, 2 of 'em, straight out of the box, but docs have all been lost... make offer, COD shipping...
\end{lstlisting}

\textit{Agent output:}
\begin{lstlisting}
This post is about selling items (plotters - which are computer graphics output devices). The seller is offering two new plotters for sale and requesting offers with COD (Cash on Delivery) shipping.

This is clearly a classified advertisement for items being sold.

FINAL: misc.forsale
\end{lstlisting}

\textit{Judge rationale:} The agent correctly categorized the post as misc.forsale with sound reasoning about the selling context, matching the reference answer perfectly. \\
\noindent\textbf{Baseline --- example 2} (gold: \texttt{sci.space}, score: 1.0)
\textit{Query:}
\begin{lstlisting}
Even better. Make up pete conrad in a Martian Suit, and have him get ou;t and throw a football to the refs.
\end{lstlisting}

\textit{Agent output:}
\begin{lstlisting}
This post appears to be discussing a creative or humorous scenario involving Pete Conrad (an astronaut) in a Martian suit throwing a football to referees. This is clearly space-related content, referencing astronauts and Mars exploration concepts.

The topic fits within space exploration and astronomy discussions.

FINAL: sci.space
\end{lstlisting}

\textit{Judge rationale:} The agent correctly identified the newsgroup as sci.space, provided sound reasoning about the space-related content (astronaut Pete Conrad, Martian suit), and delivered the answer in the expected form \\
\noindent\textbf{Baseline --- example 3} (gold: \texttt{alt.atheism}, score: 1.0)
\textit{Query:}
\begin{lstlisting}
I don't have a stat, but, unfortunately, I did read generally that both smoking and belief in the supernatural (occultish garbage) is on the rise here. This is very interesting. I thing the principle is sort of the same though: all "philosophical" ideas are generally tried out and tested mostly during college years. Whether the idea is christian or atheist doesn't always matter. But I'd like to say it's because atheists are more intelligent :) -- regards, --Adam
\end{lstlisting}

\textit{Agent output:}
\begin{lstlisting}
This post discusses philosophical and religious beliefs, specifically comparing atheism with supernatural/occult beliefs and Christianity. The author is making observations about belief systems and intelligence, touching on atheism, religious philosophy, and supernatural beliefs.

The core topic is about religion, atheism, and belief systems in general - this falls squarely into discussion of religious and anti-religious philosophy.

FINAL: alt.atheism
\end{lstlisting}

\textit{Judge rationale:} The agent correctly categorized the post as alt.atheism with sound reasoning about the content's focus on atheism and religious philosophy, matching the reference answer perfectly. \\
\noindent\textbf{Baseline --- example 4} (gold: \texttt{talk.politics.guns}, score: 0.0)
\textit{Query:}
\begin{lstlisting}
This is my first time on the network, but I am very concerned with this incident in Waco. I will refrain from stating my opinions until after I have read the FAQ. Could someone tell me how to get this? I must say that I believe the Govt. was wrong in the actions that they took in this situation.
\end{lstlisting}

\textit{Agent output:}
\begin{lstlisting}
This post discusses concerns about the Waco incident and government actions, which is a political/policy matter. The poster is asking for FAQ information and expressing a political viewpoint about government conduct.

The Waco siege was a significant political event involving government actions and civil liberties concerns. This type of discussion about government policy, political events, and civil liberties falls under political discussion categories.

Given the focus on government actions, political controversy, and policy matters related to a domestic political incident, this belongs in the general politics/miscellaneous politics category rather than a specific political issue like guns, Middle East, or religion.

FINAL: talk.politics.misc
\end{lstlisting}

\textit{Judge rationale:} The agent correctly identified the political nature and provided a valid newsgroup format, but misclassified the Waco incident (which centrally involved firearms and ATF enforcement) as general politi \\
\noindent\textbf{TEI --- example 1} (gold: \texttt{misc.forsale}, score: 1.0)
\textit{Query:}
\begin{lstlisting}
New plotters, 2 of 'em, straight out of the box, but docs have all been lost... make offer, COD shipping...
\end{lstlisting}

\textit{Agent output:}
\begin{lstlisting}
This is a straightforward classified ad offering to sell two plotters (computer peripherals/hardware devices) with an offer to negotiate price and ship cash-on-delivery. The post is a sales listing with no technical discussion, no debate, and no context suggesting it's part of a larger discussion thread.

This belongs in the for-sale category.

FINAL: misc.forsale
\end{lstlisting}

\textit{Judge rationale:} The agent correctly identified the newsgroup as misc.forsale with sound reasoning about the sales nature of the post, proper format, and perfect match to the reference answer. \\
\noindent\textbf{TEI --- example 2} (gold: \texttt{sci.space}, score: 1.0)
\textit{Query:}
\begin{lstlisting}
Even better. Make up pete conrad in a Martian Suit, and have him get ou;t and throw a football to the refs.
\end{lstlisting}

\textit{Agent output:}
\begin{lstlisting}
This post appears to be a creative/humorous suggestion about Pete Conrad (Apollo astronaut) in a fictional scenario involving Mars and throwing a football. The context suggests this is part of a discussion about space exploration activities or scenarios.

The mention of Pete Conrad, Mars, and astronaut suits points to a space exploration discussion. This is a casual, imaginative post likely part of a broader conversation in a space-focused community.

FINAL: sci.space
\end{lstlisting}

\textit{Judge rationale:} The agent correctly categorized the post as sci.space, provided sound reasoning about Pete Conrad and space context, and delivered the answer in the expected format matching the reference answer. \\
\noindent\textbf{TEI --- example 3} (gold: \texttt{alt.atheism}, score: 1.0)
\textit{Query:}
\begin{lstlisting}
I don't have a stat, but, unfortunately, I did read generally that both smoking and belief in the supernatural (occultish garbage) is on the rise here. This is very interesting. I thing the principle is sort of the same though: all "philosophical" ideas are generally tried out and tested mostly during college years. Whether the idea is christian or atheist doesn't always matter. But I'd like to say it's because atheists are more intelligent :) -- regards, --Adam
\end{lstlisting}

\textit{Agent output:}
\begin{lstlisting}
This post is part of a discussion about atheism, intelligence, and religious/supernatural beliefs. The author is making comparative claims about atheists versus believers, discussing philosophical ideas and their adoption patterns, and making a tongue-in-cheek claim about atheist intelligence. The tone and content indicate this is part of an ongoing debate in a community focused on atheism versus religion discussions.

FINAL: alt.atheism
\end{lstlisting}

\textit{Judge rationale:} The agent correctly categorized the post as 'alt.atheism' with sound reasoning about the content's focus on atheism, supernatural beliefs, and religious discussion, perfectly matching the reference an \\
\noindent\textbf{TEI --- example 4} (gold: \texttt{talk.politics.guns}, score: 0.0)
\textit{Query:}
\begin{lstlisting}
This is my first time on the network, but I am very concerned with this incident in Waco. I will refrain from stating my opinions until after I have read the FAQ. Could someone tell me how to get this? I must say that I believe the Govt. was wrong in the actions that they took in this situation.
\end{lstlisting}

\textit{Agent output:}
\begin{lstlisting}
This post is discussing the Waco siege incident and expressing concern about government actions. The poster is new to the network and asking for FAQ information while stating a political opinion about government conduct.

The Waco siege (1993) was a significant political/governmental controversy in the United States. The poster is expressing a political viewpoint about government overreach and misconduct. This is a political discussion about domestic U.S. government policy and actions, not about Middle Eastern conflicts, guns as a political issue per se, religious doctrine, or other specialized topics.

This falls under general political discussion about U.S. government actions and policy.

FINAL: talk.politics.misc
\end{lstlisting}

\textit{Judge rationale:} The agent correctly identified the task and provided a valid newsgroup category with sound reasoning about political discussion, but misclassified the post as 'talk.politics.misc' when the Waco incide \\
\subsection{\texttt{consumer\_sentiment\_5}}
\noindent\textbf{Baseline --- example 1} (gold: \texttt{very positive}, score: 0.0)
\textit{Query:}
\begin{lstlisting}
celebrated at sundance , this slight comedy of manners has winning performances and a glossy , glib charm that 's hard to beat .
\end{lstlisting}

\textit{Agent output:}
\begin{lstlisting}
This statement contains several positive indicators:

1. "celebrated at Sundance" - indicates recognition and acclaim
2. "winning performances" - explicitly positive descriptor
3. "glossy, glib charm" - positive qualities being highlighted
4. "hard to beat" - strong positive endorsement

The phrase "slight comedy of manners" is mildly diminishing (suggesting it's not weighty), but this is outweighed by the enthusiastic praise throughout. The overall tone is clearly complimentary and favorable toward the film.

FINAL: positive
\end{lstlisting}

\textit{Judge rationale:} The agent correctly identified positive sentiment with sound reasoning, but failed to capture the fine-grained distinction of 'very positive' versus merely 'positive,' which was the core task requirem \\
\noindent\textbf{Baseline --- example 2} (gold: \texttt{very negative}, score: 0.0)
\textit{Query:}
\begin{lstlisting}
what 's at stake in this film is nothing more than an obsolete , if irritating , notion of class .
\end{lstlisting}

\textit{Agent output:}
\begin{lstlisting}
This statement is critiquing a film's thematic content. Let me analyze the sentiment:

The speaker is saying that what the film addresses (class notions) is:
- "obsolete" - outdated, no longer relevant
- "irritating" - annoying
- "nothing more than" - dismissive phrasing suggesting the film's stakes are limited/trivial

The tone is critical and dismissive. The speaker views the film's central concern as both outdated and bothersome, which indicates a negative evaluation of the film's thematic substance.

However, this isn't extremely harsh - it's a measured critique rather than a scathing attack. The speaker isn't calling the film terrible overall, just criticizing what it chooses to focus on.

FINAL: negative
\end{lstlisting}

\textit{Judge rationale:} The agent correctly identified negative sentiment with sound reasoning about dismissive language and critical tone, but failed to capture the 'very negative' intensity required by the reference answer \\
\noindent\textbf{Baseline --- example 3} (gold: \texttt{negative}, score: 0.0)
\textit{Query:}
\begin{lstlisting}
if myers decides to make another austin powers movie , maybe he should just stick with austin and dr evil .
\end{lstlisting}

\textit{Agent output:}
\begin{lstlisting}
This statement expresses a conditional suggestion about a potential future Austin Powers movie. The speaker is offering constructive advice that if Myers were to make another film, he should focus on the main characters (Austin and Dr. Evil) rather than potentially diluting the story with other characters.

The tone is:
- Not strongly negative (not criticizing what exists)
- Not strongly positive (not enthusiastically praising)
- Somewhat advisory/suggestive in nature
- Mildly critical of the idea of including too many characters

This falls into a neutral-to-slightly-negative range, as it implies concern about potential creative decisions while offering a measured suggestion.

FINAL: neutral
\end{lstlisting}

\textit{Judge rationale:} The agent provided extensive analysis but misclassified the sentiment as 'neutral' when the reference answer is 'negative', and the phrasing 'maybe he should just stick with' implies dissatisfaction w \\
\noindent\textbf{Baseline --- example 4} (gold: \texttt{positive}, score: 1.0)
\textit{Query:}
\begin{lstlisting}
in its ragged , cheap and unassuming way , the movie works .
\end{lstlisting}

\textit{Agent output:}
\begin{lstlisting}
This statement contains mixed sentiment elements that need careful analysis:

**Positive aspects:**
- "the movie works" - a clear positive statement indicating the film functions/succeeds
- The phrase "in its...way" suggests acceptance despite flaws

**Negative/Limiting aspects:**
- "ragged" - suggests rough, unpolished quality
- "cheap" - implies low budget or poor production values
- "unassuming" - suggests lack of pretension, but in context with other descriptors, emphasizes modesty/limitation

**Overall interpretation:**
The statement is fundamentally positive (the movie works), but heavily qualified by acknowledgment of the film's limitations and rough qualities. The speaker is saying the film succeeds despite--or perhaps because of--its humble, low-budget nature. This is a backhanded compliment that ultimately endorses the film while being honest about its shortcomings.

This represents a **positive** sentiment with reservations, rather than a strong endorsement. The core message is approval, though tempered by criticism.

FINAL: positive
\end{lstlisting}

\textit{Judge rationale:} The agent correctly classified the sentiment as positive, provided thorough and logically sound analysis of the mixed elements, delivered the answer in the required format, and matched the reference a \\
\noindent\textbf{TEI --- example 1} (gold: \texttt{very positive}, score: 1.0)
\textit{Query:}
\begin{lstlisting}
celebrated at sundance , this slight comedy of manners has winning performances and a glossy , glib charm that 's hard to beat .
\end{lstlisting}

\textit{Agent output:}
\begin{lstlisting}
This statement contains multiple positive descriptors: "celebrated," "winning performances," "glossy," "glib charm," and "hard to beat." The phrase "hard to beat" is an intensifier emphasizing the strength of the charm. The overall tone is enthusiastically favorable with strong praise for the film's qualities.

FINAL: very positive
\end{lstlisting}

\textit{Judge rationale:} The agent correctly classified the sentiment as 'very positive' with sound reasoning identifying multiple positive descriptors and intensifiers, matching the reference answer perfectly. \\
\noindent\textbf{TEI --- example 2} (gold: \texttt{very negative}, score: 1.0)
\textit{Query:}
\begin{lstlisting}
what 's at stake in this film is nothing more than an obsolete , if irritating , notion of class .
\end{lstlisting}

\textit{Agent output:}
\begin{lstlisting}
This statement criticizes the film for dealing with an outdated and annoying concept (class). The phrase "nothing more than" suggests the film's central concern is trivial or limited. "Obsolete" and "irritating" are two negative descriptors combined, indicating strong criticism.

FINAL: very negative
\end{lstlisting}

\textit{Judge rationale:} The agent correctly identified the sentiment as 'very negative' with sound reasoning about the critical language ('obsolete,' 'irritating,' 'nothing more than'), properly addressing the fine-grained s \\
\noindent\textbf{TEI --- example 3} (gold: \texttt{negative}, score: 0.0)
\textit{Query:}
\begin{lstlisting}
if myers decides to make another austin powers movie , maybe he should just stick with austin and dr evil .
\end{lstlisting}

\textit{Agent output:}
\begin{lstlisting}
This statement suggests that if Myers makes another Austin Powers movie, he should focus on the main characters (Austin and Dr. Evil) rather than other elements. This is a qualified suggestion with an implicit criticism--it implies that other aspects of previous films were problematic or unnecessary. The "maybe he should just stick with" phrasing indicates a reservation or limitation about what would work, rather than enthusiastic praise.

This is mixed sentiment: there's a mild positive (acknowledging these characters work) combined with an implicit negative (suggesting other elements don't).

FINAL: neutral
\end{lstlisting}

\textit{Judge rationale:} The agent provided extensive reasoning but classified the sentiment as 'neutral' when the reference answer is 'negative,' and the statement's implicit criticism ('just stick with') clearly conveys neg \\
\noindent\textbf{TEI --- example 4} (gold: \texttt{positive}, score: 0.0)
\textit{Query:}
\begin{lstlisting}
in its ragged , cheap and unassuming way , the movie works .
\end{lstlisting}

\textit{Agent output:}
\begin{lstlisting}
This statement contains multiple negative descriptors ("ragged," "cheap," "unassuming") combined with a positive assertion ("the movie works"). The negative words describe the film's qualities, yet the overall verdict is favorable despite these limitations. This is a qualified praise with reservations--the movie succeeds despite its flaws.

FINAL: neutral
\end{lstlisting}

\textit{Judge rationale:} The agent misclassified a clearly positive sentiment as neutral by overweighting surface-level negative descriptors while ignoring the affirming conclusion 'the movie works,' which signals overall app \\
\subsection{\texttt{cybersec\_alert}}
\noindent\textbf{Baseline --- example 1} (gold: \texttt{data\_exfiltration}, score: 1.0)
\textit{Query:}
\begin{lstlisting}
Firewall logs show an employee workstation established 47 consecutive HTTPS connections to a newly registered domain in the past hour, transferring approximately 2.3 GB of data to an external IP address located in a non-business region. The destination server has no prior communication history with our organization.
\end{lstlisting}

\textit{Agent output:}
\begin{lstlisting}
**Analysis:**

This alert exhibits several key indicators:

1. **High-volume data transfer** (2.3 GB in one hour) to an external destination
2. **Newly registered domain** - suggests attacker infrastructure or compromised domain
3. **Non-business region** - outside normal operational geography
4. **No prior communication history** - anomalous destination
5. **Multiple consecutive connections** - systematic data exfiltration pattern
6. **Employee workstation as source** - suggests either compromised device or insider activity

The pattern is consistent with **data exfiltration** - the systematic unauthorized transfer of large volumes of data to an external, untrusted location. While this could involve malware (which may be present), the primary threat being demonstrated here is the actual exfiltration of data itself.

The characteristics rule out:
- **Phishing** (no email/social engineering component evident)
- **Brute force** (no authentication attack pattern)
- **DDoS** (no volumetric attack to disrupt service)
- **Insider threat** (while possible, the use of newly registered domain and non-business region suggests external attacker or compromised device rather than intentional insider action)
- **Malware** (while malware may be the delivery mechanism, the alert describes the outcome: data exfiltration)

FINAL: data_exfiltration
\end{lstlisting}

\textit{Judge rationale:} The agent correctly classified the alert as data\_exfiltration with sound reasoning that systematically analyzed all indicators and ruled out alternative categories, perfectly matching the reference an \\
\noindent\textbf{Baseline --- example 2} (gold: \texttt{insider\_threat}, score: 0.0)
\textit{Query:}
\begin{lstlisting}
Finance department user executed 23 failed attempts to access restricted merger and acquisition folders on the executive file server within a 10-minute window, despite having no M&A project assignments.
\end{lstlisting}

\textit{Agent output:}
\begin{lstlisting}
**Analysis:**

This alert exhibits characteristics of a brute force attack:

1. **Multiple failed attempts**: 23 failed access attempts in a short timeframe (10 minutes) indicates systematic trial-and-error behavior
2. **Rapid succession**: The concentrated 10-minute window shows automated or deliberate repeated attempts
3. **Unauthorized access target**: Attempts directed at restricted M&A folders the user shouldn't need
4. **Anomalous behavior**: User has no legitimate project assignment requiring this access

While this could potentially indicate an insider threat (compromised account or malicious insider), the pattern of **repeated failed authentication/authorization attempts** is the defining characteristic of brute force activity. The attacker is attempting to bypass access controls through multiple tries.

FINAL: brute_force
\end{lstlisting}

\textit{Judge rationale:} The agent provided well-structured reasoning and correct format, but misclassified an insider threat scenario as brute\_force--the user is authenticated (finance department user) attempting unauthorized \\
\noindent\textbf{Baseline --- example 3} (gold: \texttt{ddos}, score: 1.0)
\textit{Query:}
\begin{lstlisting}
Database connection pool exhausted due to 15,000 simultaneous TCP connection attempts from distributed sources. Application tier reporting complete service outage for past 8 minutes.
\end{lstlisting}

\textit{Agent output:}
\begin{lstlisting}
**Analysis:**

This alert describes a scenario where:
- 15,000 simultaneous TCP connection attempts are being made
- Connections originate from distributed sources (multiple IPs)
- The attack exhausts the database connection pool
- Result is complete service outage (unavailability)

This is a classic Distributed Denial of Service (DDoS) attack pattern. The attacker is flooding the database with connection requests from multiple sources to overwhelm resources and prevent legitimate users from accessing the service.

Key indicators:
- High volume of connections (15,000)
- Distributed sources (hallmark of DDoS)
- Resource exhaustion (connection pool)
- Service unavailability as the outcome

This is not brute force (which targets authentication), malware, phishing, data exfiltration, or insider threat. It's a volumetric attack designed to deny service.

FINAL: ddos
\end{lstlisting}

\textit{Judge rationale:} The agent correctly classified the threat as 'ddos', provided sound reasoning identifying all key attack indicators (distributed sources, volume, resource exhaustion, service denial), and delivered th \\
\noindent\textbf{Baseline --- example 4} (gold: \texttt{phishing}, score: 1.0)
\textit{Query:}
\begin{lstlisting}
Email security system quarantined messages claiming to be password expiration notices from IT. The emails contain QR codes that, when scanned, direct mobile devices to fraudulent authentication pages designed to harvest both passwords and MFA codes in real-time.
\end{lstlisting}

\textit{Agent output:}
\begin{lstlisting}
**Analysis:**

This alert describes a sophisticated phishing attack with the following indicators:

1. **Deceptive Sender Spoofing**: Emails falsely claim to be from IT department, impersonating a trusted internal source
2. **Social Engineering**: Uses urgency/authority (password expiration notices) to pressure victims into action
3. **Credential Harvesting**: QR codes direct to fraudulent authentication pages designed to steal:
   - Passwords
   - MFA codes (multi-factor authentication)
4. **Attack Vector**: Email-based with mobile targeting

This is a classic phishing attack pattern - the primary goal is to deceive users into voluntarily providing authentication credentials through fraudulent means. While the attack may lead to account compromise or data exfiltration, the attack itself is fundamentally a phishing campaign.

**FINAL: phishing**
\end{lstlisting}

\textit{Judge rationale:} The agent correctly classified the threat as phishing with comprehensive reasoning that accurately identifies all key attack indicators (spoofing, social engineering, credential harvesting), matches t \\
\noindent\textbf{TEI --- example 1} (gold: \texttt{data\_exfiltration}, score: 1.0)
\textit{Query:}
\begin{lstlisting}
Firewall logs show an employee workstation established 47 consecutive HTTPS connections to a newly registered domain in the past hour, transferring approximately 2.3 GB of data to an external IP address located in a non-business region. The destination server has no prior communication history with our organization.
\end{lstlisting}

\textit{Agent output:}
\begin{lstlisting}
This alert describes unauthorized transfer of a large volume of sensitive data (2.3 GB) from an employee workstation to an external, previously unknown destination. The indicators include:

- Suspicious pattern of connections to a newly registered domain
- Large data transfer volume (2.3 GB)
- External IP in a non-business region
- No prior legitimate communication history
- Employee workstation as the source

This represents unauthorized exfiltration of organizational data outside the network, which is the defining characteristic of data exfiltration.

FINAL: data_exfiltration
\end{lstlisting}

\textit{Judge rationale:} The agent correctly classified the threat as data\_exfiltration with sound reasoning identifying all key indicators, though the task only required the classification itself without extended explanation \\
\noindent\textbf{TEI --- example 2} (gold: \texttt{insider\_threat}, score: 0.0)
\textit{Query:}
\begin{lstlisting}
Finance department user executed 23 failed attempts to access restricted merger and acquisition folders on the executive file server within a 10-minute window, despite having no M&A project assignments.
\end{lstlisting}

\textit{Agent output:}
\begin{lstlisting}
This alert describes repeated failed authentication attempts to access resources by a user who lacks authorization for those resources. The key indicators are:

- Multiple failed login/access attempts (23 in 10 minutes)
- Systematic attempts to access specific restricted folders
- User has no legitimate business need for M&A access
- Rapid succession of attempts suggests automated or deliberate enumeration

This is a brute-force attack pattern - the user is making repeated attempts to gain access to restricted resources, likely trying different credentials, tokens, or access methods.

FINAL: brute_force
\end{lstlisting}

\textit{Judge rationale:} The agent correctly identified the classification task and provided proper formatting, but misclassified an insider threat scenario as brute\_force by focusing on the failed attempts pattern rather tha \\
\noindent\textbf{TEI --- example 3} (gold: \texttt{ddos}, score: 1.0)
\textit{Query:}
\begin{lstlisting}
Database connection pool exhausted due to 15,000 simultaneous TCP connection attempts from distributed sources. Application tier reporting complete service outage for past 8 minutes.
\end{lstlisting}

\textit{Agent output:}
\begin{lstlisting}
This alert describes a large volume of simultaneous connection attempts from distributed sources causing service unavailability. This is a classic Distributed Denial of Service (DDoS) attack pattern - overwhelming system resources (connection pool) with traffic from multiple sources to make the service unavailable.

FINAL: ddos
\end{lstlisting}

\textit{Judge rationale:} The agent correctly classified the alert as 'ddos', provided sound reasoning about distributed sources overwhelming resources, and the output matches the reference answer exactly. \\
\noindent\textbf{TEI --- example 4} (gold: \texttt{phishing}, score: 1.0)
\textit{Query:}
\begin{lstlisting}
Email security system quarantined messages claiming to be password expiration notices from IT. The emails contain QR codes that, when scanned, direct mobile devices to fraudulent authentication pages designed to harvest both passwords and MFA codes in real-time.
\end{lstlisting}

\textit{Agent output:}
\begin{lstlisting}
This alert describes a phishing attack. The emails use social engineering (impersonation of IT department), deceptive content (fake password expiration notices), and credential harvesting mechanisms (fraudulent authentication pages designed to steal passwords and MFA codes). The QR codes serve as the delivery mechanism for the phishing attack.

While the ultimate goal is credential theft, the primary threat vector is the phishing attack itself - the deceptive communication designed to trick users into visiting malicious pages.

FINAL: phishing
\end{lstlisting}

\textit{Judge rationale:} The agent correctly classified the threat as phishing with sound reasoning about social engineering, credential harvesting, and deceptive mechanisms, matching the reference answer perfectly. \\
\subsection{\texttt{edtech\_question\_router}}
\noindent\textbf{Baseline --- example 1} (gold: \texttt{Chemistry}, score: 0.0)
\textit{Query:}
\begin{lstlisting}
The compounds shown below have similar weights but significantly different boiling points. Match the compound with its boiling point. (Boiling points \( \left.\left(^{o} C\right) 28,57,100,141\right) \) A. Methyl acetate=100, 2-butanol=141, 2-methyl butane=28, Propanoic acid=57 B. Methyl acetate=57, 2-butanol=100, 2-methyl butane=28, Propanoic acid=141 c. Methyl acetate=28, 2-butanol=100, 2-methyl butane=57, Propanoic acid=141 D. Methyl acetate=141, 2-butanol=57, 2-methy butane=28, Propanoic acid=100
\end{lstlisting}

\textit{Agent output:}
\begin{lstlisting}
I need to match compounds with similar molecular weights to their boiling points based on intermolecular forces.

**Analyzing the compounds:**

1. **2-methyl butane (C5H12)**: MW ~ 72
   - Only dispersion forces (nonpolar hydrocarbon)
   - Weakest intermolecular forces -> lowest boiling point

2. **Methyl acetate (C3H6O2)**: MW ~ 74
   - Polar molecule with dipole-dipole interactions
   - No hydrogen bonding
   - Moderate boiling point

3. **Propanoic acid (C3H6O2)**: MW ~ 74
   - Carboxylic acid with strong hydrogen bonding
   - Can form dimers through H-bonding
   - High boiling point

4. **2-butanol (C4H10O)**: MW ~ 74
   - Alcohol with hydrogen bonding capability
   - Strong intermolecular forces
   - High boiling point

**Expected boiling point order:**
- 2-methyl butane: **28C** (only dispersion forces)
- Methyl acetate: **57C** (dipole-dipole, no H-bonding)
- 2-butanol: **100C** (hydrogen bonding from -OH)
- Propanoic acid: **141C** (strong hydrogen bonding, dimer formation)

**Checking the options:**
- Option A: Methyl acetate=100  (too high)
- Option B: Methyl acetate=57 , 2-butanol=100 , 2-methyl butane=28 , Propanoic acid=141 
- Option C: 2-methyl butane=57  (too high)
- Option D: Methyl acetate=141  (too high)

FINAL: B
\end{lstlisting}

\textit{Judge rationale:} The agent was asked to route/categorize the question but instead solved the chemistry problem itself, completely missing the task of classification. \\
\noindent\textbf{Baseline --- example 2} (gold: \texttt{Biology}, score: 1.0)
\textit{Query:}
\begin{lstlisting}
Which of the following is true for Krebs cycle? A. Isocitrate \( \rightarrow \) Fumrate \( \rightarrow \alpha \) ketoglutarate B. \( \alpha \) Ketoglutarate \( \rightarrow \) Oxalosuccinate \( \rightarrow \) Isocitrate c. oxalosuccinate \( \rightarrow \) Isocitrate \( \rightarrow \alpha \) ketoglutarate D. Isocitrate \( \rightarrow \alpha \) ketoglutarate \( \rightarrow \) Succinate
\end{lstlisting}

\textit{Agent output:}
\begin{lstlisting}
This question is asking about the sequence of intermediates in the Krebs cycle (also called the citric acid cycle or TCA cycle).

Let me trace through the relevant portion of the Krebs cycle:

The Krebs cycle sequence includes:
- Citrate -> Isocitrate -> -ketoglutarate -> Succinyl-CoA -> Succinate -> Fumarate -> Malate -> Oxaloacetate

Now let me evaluate each option:

**A. Isocitrate -> Fumarate -> -ketoglutarate**
- This is incorrect. The sequence should be Isocitrate -> -ketoglutarate (not directly to Fumarate), and Fumarate comes much later.

**B. -Ketoglutarate -> Oxalosuccinate -> Isocitrate**
- This is incorrect. This reverses the direction. The correct direction is Isocitrate -> Oxalosuccinate -> -ketoglutarate (going forward in the cycle).

**C. Oxalosuccinate -> Isocitrate -> -ketoglutarate**
- This is incorrect. Oxalosuccinate is an intermediate that forms between Isocitrate and -ketoglutarate, but the sequence shown here is backwards.

**D. Isocitrate -> -ketoglutarate -> Succinate**
- This is correct! The sequence is: Isocitrate -> -ketoglutarate -> Succinyl-CoA -> Succinate. While Succinyl-CoA is technically an intermediate between -ketoglutarate and Succinate, this option correctly shows the major intermediates in the proper order.

FINAL: Biology
\end{lstlisting}

\textit{Judge rationale:} The agent correctly identified the category as 'Biology' and provided sound biochemical reasoning, but the task was simply to route/categorize the question, not to solve it--the extensive answer analys \\
\noindent\textbf{Baseline --- example 3} (gold: \texttt{Maths}, score: 1.0)
\textit{Query:}
\begin{lstlisting}
Find the value of \( x \) if \( \sin \left\{\sin ^{-1} \frac{1}{5}+\cos ^{-1} x\right\}=1 \)
\end{lstlisting}

\textit{Agent output:}
\begin{lstlisting}
I need to find the value of x given that $\sin \left\{\sin ^{-1} \frac{1}{5}+\cos ^{-1} x\right\}=1$.

**Reasoning:**

For $\sin(\theta) = 1$, we need $\theta = \frac{\pi}{2}$ (within the principal range).

So: $\sin ^{-1} \frac{1}{5}+\cos ^{-1} x = \frac{\pi}{2}$

This gives us: $\cos ^{-1} x = \frac{\pi}{2} - \sin ^{-1} \frac{1}{5}$

**Using the complementary angle identity:**

I know that $\sin^{-1}(a) + \cos^{-1}(a) = \frac{\pi}{2}$ for $a \in [-1,1]$.

This means: $\cos^{-1}(a) = \frac{\pi}{2} - \sin^{-1}(a)$

Therefore: $\frac{\pi}{2} - \sin ^{-1} \frac{1}{5} = \cos^{-1}\left(\frac{1}{5}\right)$

**Solving for x:**

From $\cos ^{-1} x = \cos^{-1}\left(\frac{1}{5}\right)$

We get: $x = \frac{1}{5}$

**Verification:**
- $\sin^{-1}\frac{1}{5} + \cos^{-1}\frac{1}{5} = \frac{\pi}{2}$ 
- $\sin\left(\frac{\pi}{2}\right) = 1$ 

This is a mathematics problem involving inverse trigonometric functions and trigonometric equations.

FINAL: Maths
\end{lstlisting}

\textit{Judge rationale:} The agent correctly identified the category as 'Maths', provided sound mathematical reasoning to solve the problem, and delivered the exact answer matching the reference. \\
\noindent\textbf{Baseline --- example 4} (gold: \texttt{Physics}, score: 1.0)
\textit{Query:}
\begin{lstlisting}
A uniform solid cylinder of mass \( 2 k g \) and radius \( 0.2 m \) is released from the rest at the top of a semicircular track of radius \( 0.7 m \) cut in block of mass \( M= \) \( 3 k g \) as shown. The block is resting on a smooth horizontal surface and the cylinder rolls down without slipping. The distance moved by the block when the cylinder reaches the bottom of the track is: \( \mathbf{A} \cdot 0.3 m \) B. \( 0.5 m \) c. \( 0.7 m \) D. \( 0.2 m \)
\end{lstlisting}

\textit{Agent output:}
\begin{lstlisting}
I need to analyze this problem involving a cylinder rolling down a semicircular track on a movable block.

**Given Information:**
- Cylinder: mass m = 2 kg, radius r = 0.2 m
- Block: mass M = 3 kg
- Track radius: R = 0.7 m
- Cylinder rolls without slipping
- Block is on a smooth (frictionless) horizontal surface
- Cylinder starts from rest at the top

**Key Physics Concepts:**

Since the horizontal surface is smooth (frictionless), there are no external horizontal forces on the system. Therefore, the center of mass of the system doesn't move horizontally.

**Using Conservation of Momentum:**

Initial momentum = 0 (everything at rest)

When the cylinder reaches the bottom:
- Let v_c = velocity of cylinder (relative to ground)
- Let v_b = velocity of block (relative to ground)

Horizontal momentum conservation:
$$m \cdot v_c + M \cdot v_b = 0$$
$$2v_c + 3v_b = 0$$
$$v_c = -1.5v_b$$

The negative sign indicates opposite directions.

**Using Center of Mass Position:**

The center of mass of the system doesn't move horizontally. If the block moves distance x to the right, the cylinder's center of mass must move such that:

$$m \cdot x_c + M \cdot x_b = 0$$

where x_c is the displacement of the cylinder's center and x_b is the displacement of the block.

The cylinder descends a vertical distance R = 0.7 m along the track. Relative to the block, the cylinder moves from the top to the bottom of the semicircle.

If the block moves distance x to the right, the cylinder's center moves distance (R - x) to the left relative to the ground (approximately, considering the geometry).

Actually, more precisely: the cylinder's center moves horizontally by some amount relative to the block. If the block moves x to the right, and the cylinder's center moves distance d to the left relative to ground:

$$2(-d) + 3(x) = 0$$
$$3x = 2d$$

The cylinder descends R = 0.7 m vertically. The horizontal displacement of the cylinder relative to the block is approximately R (from the geometry of the
\end{lstlisting}

\textit{Judge rationale:} The agent correctly identified this as a Physics problem, provided sound reasoning using conservation principles, and arrived at the correct category matching the reference answer. \\
\noindent\textbf{TEI --- example 1} (gold: \texttt{Chemistry}, score: 0.0)
\textit{Query:}
\begin{lstlisting}
The compounds shown below have similar weights but significantly different boiling points. Match the compound with its boiling point. (Boiling points \( \left.\left(^{o} C\right) 28,57,100,141\right) \) A. Methyl acetate=100, 2-butanol=141, 2-methyl butane=28, Propanoic acid=57 B. Methyl acetate=57, 2-butanol=100, 2-methyl butane=28, Propanoic acid=141 c. Methyl acetate=28, 2-butanol=100, 2-methyl butane=57, Propanoic acid=141 D. Methyl acetate=141, 2-butanol=57, 2-methy butane=28, Propanoic acid=100
\end{lstlisting}

\textit{Agent output:}
\begin{lstlisting}
This question asks you to match chemical compounds with their boiling points based on understanding intermolecular forces and molecular structure.

The key concept is that boiling points depend on intermolecular forces:
- **2-methyl butane** (alkane): only weak London dispersion forces -> lowest boiling point (28C)
- **Methyl acetate** (ester): dipole-dipole interactions -> moderate boiling point (57C)
- **Propanoic acid** (carboxylic acid): strong hydrogen bonding -> high boiling point (141C)
- **2-butanol** (alcohol): hydrogen bonding -> high boiling point (100C)

Propanoic acid has the strongest intermolecular forces due to extensive hydrogen bonding in carboxylic acids, so it has the highest boiling point (141C). 2-butanol also has hydrogen bonding but is less extensive than carboxylic acids (100C). Methyl acetate has weaker dipole-dipole forces (57C). 2-methyl butane has only dispersion forces (28C).

The correct answer is **B**.

FINAL: Physics
\end{lstlisting}

\textit{Judge rationale:} The agent was asked to route the question to a category and should have output 'Chemistry', but instead provided a detailed chemistry explanation followed by the incorrect category 'Physics', complete \\
\noindent\textbf{TEI --- example 2} (gold: \texttt{Biology}, score: 1.0)
\textit{Query:}
\begin{lstlisting}
Which of the following is true for Krebs cycle? A. Isocitrate \( \rightarrow \) Fumrate \( \rightarrow \alpha \) ketoglutarate B. \( \alpha \) Ketoglutarate \( \rightarrow \) Oxalosuccinate \( \rightarrow \) Isocitrate c. oxalosuccinate \( \rightarrow \) Isocitrate \( \rightarrow \alpha \) ketoglutarate D. Isocitrate \( \rightarrow \alpha \) ketoglutarate \( \rightarrow \) Succinate
\end{lstlisting}

\textit{Agent output:}
\begin{lstlisting}
FINAL: Biology
\end{lstlisting}

\textit{Judge rationale:} The agent correctly identified the question as Biology, matching the reference answer exactly with proper formatting and no extraneous content. \\
\noindent\textbf{TEI --- example 3} (gold: \texttt{Maths}, score: 1.0)
\textit{Query:}
\begin{lstlisting}
Find the value of \( x \) if \( \sin \left\{\sin ^{-1} \frac{1}{5}+\cos ^{-1} x\right\}=1 \)
\end{lstlisting}

\textit{Agent output:}
\begin{lstlisting}
FINAL: Maths
\end{lstlisting}

\textit{Judge rationale:} The agent correctly identified the mathematical question and routed it to 'Maths', matching the reference answer exactly with proper formatting. \\
\noindent\textbf{TEI --- example 4} (gold: \texttt{Physics}, score: 1.0)
\textit{Query:}
\begin{lstlisting}
A uniform solid cylinder of mass \( 2 k g \) and radius \( 0.2 m \) is released from the rest at the top of a semicircular track of radius \( 0.7 m \) cut in block of mass \( M= \) \( 3 k g \) as shown. The block is resting on a smooth horizontal surface and the cylinder rolls down without slipping. The distance moved by the block when the cylinder reaches the bottom of the track is: \( \mathbf{A} \cdot 0.3 m \) B. \( 0.5 m \) c. \( 0.7 m \) D. \( 0.2 m \)
\end{lstlisting}

\textit{Agent output:}
\begin{lstlisting}
This question involves a cylinder rolling down a semicircular track on a movable block, requiring analysis of:
- Rolling motion (rotational and translational dynamics)
- Conservation of momentum (horizontal direction)
- Conservation of energy
- Constraint of rolling without slipping
- Motion of the block on a frictionless surface

These are all concepts from classical mechanics and physics.

FINAL: Physics
\end{lstlisting}

\textit{Judge rationale:} The agent correctly categorized the question as Physics with proper identification of relevant mechanics concepts, matching the reference answer exactly. \\
\subsection{\texttt{edu\_gsm8k}}
\noindent\textbf{Baseline --- example 1} (gold: \texttt{49}, score: 1.0)
\textit{Query:}
\begin{lstlisting}
There are 96 fourth-graders at Small Tree School. 43 of them are girls. On Friday, 5 fourth-grade girls and 4 fourth-grade boys were absent. How many fourth grade boys were at Small Tree School on Friday?
\end{lstlisting}

\textit{Agent output:}
\begin{lstlisting}
I need to find how many fourth-grade boys were at school on Friday.

**Step 1: Find the total number of fourth-grade boys at the school.**
- Total fourth-graders: 96
- Fourth-grade girls: 43
- Fourth-grade boys: 96 - 43 = 53

**Step 2: Find how many fourth-grade boys were absent on Friday.**
- Boys absent: 4

**Step 3: Find how many fourth-grade boys were present on Friday.**
- Boys present on Friday: 53 - 4 = 49

FINAL: 49
\end{lstlisting}

\textit{Judge rationale:} The agent correctly identified the total boys (53), subtracted the absent boys (4), arrived at the correct answer (49), with clear logical steps and perfect alignment to the task. \\
\noindent\textbf{Baseline --- example 2} (gold: \texttt{150}, score: 1.0)
\textit{Query:}
\begin{lstlisting}
Ali is a dean of a private school where he teaches one class. John is also a dean of a public school. John has two classes in his school. Each class has 1/8 the capacity of Ali's class which has the capacity of 120 students. What is the combined capacity of both schools?
\end{lstlisting}

\textit{Agent output:}
\begin{lstlisting}
I need to find the combined capacity of both schools.

**Ali's school:**
- Ali teaches one class
- Ali's class capacity = 120 students
- Ali's school capacity = 120 students

**John's school:**
- John has two classes
- Each of John's classes has 1/8 the capacity of Ali's class
- Capacity of each of John's classes = 1/8 x 120 = 15 students
- John's school capacity = 2 x 15 = 30 students

**Combined capacity:**
- Combined capacity = Ali's school + John's school
- Combined capacity = 120 + 30 = 150 students

FINAL: 150
\end{lstlisting}

\textit{Judge rationale:} The agent correctly identified all given information, performed accurate calculations (1/8 x 120 = 15, 2 x 15 = 30, 120 + 30 = 150), showed clear logical steps, and arrived at the correct answer of 15 \\
\noindent\textbf{Baseline --- example 3} (gold: \texttt{19}, score: 1.0)
\textit{Query:}
\begin{lstlisting}
3 customers were kicked out of the Walmart for refusing to wear masks. A number equals to four times that many minus 5 were kicked out for shoplifting. Three times the number of shoplifters were kicked out for physical violence over goods on sale. If a total of 50 people were kicked out of the Walmart, how many were kicked out for other reasons?
\end{lstlisting}

\textit{Agent output:}
\begin{lstlisting}
I need to find how many people were kicked out for other reasons.

Let me define the groups:
- Mask refusal: 3 customers
- Shoplifting: 4 times the mask refusers minus 5
- Physical violence: 3 times the shoplifters
- Other reasons: unknown

**Step 1: Calculate shoplifters**
Shoplifters = 4 x 3 - 5 = 12 - 5 = 7

**Step 2: Calculate physical violence**
Physical violence = 3 x 7 = 21

**Step 3: Calculate other reasons**
Total kicked out = 50
Other reasons = 50 - (mask refusal + shoplifting + physical violence)
Other reasons = 50 - (3 + 7 + 21)
Other reasons = 50 - 31
Other reasons = 19

FINAL: 19
\end{lstlisting}

\textit{Judge rationale:} The agent correctly identified all groups, performed accurate arithmetic calculations step-by-step, and arrived at the correct answer of 19 matching the reference. \\
\noindent\textbf{Baseline --- example 4} (gold: \texttt{159}, score: 1.0)
\textit{Query:}
\begin{lstlisting}
Nate is feeding his livestock hay. Each goat needs 5 pounds, and each sheep needs 3 pounds less than twice the amount each goat needs. If there are 15 goats and 12 sheep, how many pounds of hay does Nate need?
\end{lstlisting}

\textit{Agent output:}
\begin{lstlisting}
I need to find the total amount of hay needed for all the livestock.

**Step 1: Find how much hay each goat needs.**
Each goat needs 5 pounds.

**Step 2: Find how much hay each sheep needs.**
Each sheep needs 3 pounds less than twice the amount each goat needs.
- Twice the amount each goat needs = 2 x 5 = 10 pounds
- 3 pounds less than that = 10 - 3 = 7 pounds per sheep

**Step 3: Calculate total hay for all goats.**
15 goats x 5 pounds per goat = 75 pounds

**Step 4: Calculate total hay for all sheep.**
12 sheep x 7 pounds per sheep = 84 pounds

**Step 5: Calculate total hay needed.**
75 + 84 = 159 pounds

FINAL: 159
\end{lstlisting}

\textit{Judge rationale:} The agent correctly interpreted the problem, applied sound arithmetic reasoning at each step, followed proper format, and arrived at the exact reference answer of 159 pounds. \\
\noindent\textbf{TEI --- example 1} (gold: \texttt{49}, score: 1.0)
\textit{Query:}
\begin{lstlisting}
There are 96 fourth-graders at Small Tree School. 43 of them are girls. On Friday, 5 fourth-grade girls and 4 fourth-grade boys were absent. How many fourth grade boys were at Small Tree School on Friday?
\end{lstlisting}

\textit{Agent output:}
\begin{lstlisting}
I need to find how many fourth-grade boys were at school on Friday.

**Step 1: Find the total number of boys at the school**
- Total fourth-graders: 96
- Girls: 43
- Boys: 96 - 43 = 53 boys total

**Step 2: Find how many boys were absent on Friday**
- Boys absent: 4

**Step 3: Find how many boys were present on Friday**
- Boys present on Friday: 53 - 4 = 49 boys

FINAL: 49
\end{lstlisting}

\textit{Judge rationale:} The agent correctly identified the total boys (53), subtracted the absent boys (4), arrived at the correct answer (49), with clear logical steps and proper formatting. \\
\noindent\textbf{TEI --- example 2} (gold: \texttt{150}, score: 1.0)
\textit{Query:}
\begin{lstlisting}
Ali is a dean of a private school where he teaches one class. John is also a dean of a public school. John has two classes in his school. Each class has 1/8 the capacity of Ali's class which has the capacity of 120 students. What is the combined capacity of both schools?
\end{lstlisting}

\textit{Agent output:}
\begin{lstlisting}
I need to find the combined capacity of both schools.

**Given information:**
- Ali's class capacity: 120 students
- Ali teaches 1 class
- John has 2 classes
- Each of John's classes has 1/8 the capacity of Ali's class

**Step 1: Find the capacity of each of John's classes**
Each of John's classes = 1/8 x 120 = 15 students

**Step 2: Find the total capacity of John's school**
John has 2 classes, so:
Total capacity of John's school = 2 x 15 = 30 students

**Step 3: Find the total capacity of Ali's school**
Ali has 1 class, so:
Total capacity of Ali's school = 1 x 120 = 120 students

**Step 4: Find the combined capacity of both schools**
Combined capacity = 120 + 30 = 150 students

FINAL: 150
\end{lstlisting}

\textit{Judge rationale:} The agent correctly interpreted the problem, applied sound arithmetic reasoning through all steps, followed proper format, and arrived at the correct answer of 150 matching the reference. \\
\noindent\textbf{TEI --- example 3} (gold: \texttt{19}, score: 1.0)
\textit{Query:}
\begin{lstlisting}
3 customers were kicked out of the Walmart for refusing to wear masks. A number equals to four times that many minus 5 were kicked out for shoplifting. Three times the number of shoplifters were kicked out for physical violence over goods on sale. If a total of 50 people were kicked out of the Walmart, how many were kicked out for other reasons?
\end{lstlisting}

\textit{Agent output:}
\begin{lstlisting}
I need to find how many people were kicked out for other reasons.

**Given information:**
- Mask refusal: 3 customers
- Shoplifting: 4 times the mask refusers minus 5
- Physical violence: 3 times the shoplifters
- Total kicked out: 50
- Find: kicked out for other reasons

**Step 1: Calculate shoplifters**
Shoplifters = 4 x 3 - 5
Shoplifters = 12 - 5
Shoplifters = 7

**Step 2: Calculate physical violence**
Physical violence = 3 x 7
Physical violence = 21

**Step 3: Calculate total for these three reasons**
Mask refusal + Shoplifting + Physical violence
= 3 + 7 + 21
= 31

**Step 4: Calculate other reasons**
Other reasons = Total - (Mask refusal + Shoplifting + Physical violence)
Other reasons = 50 - 31
Other reasons = 19

FINAL: 19
\end{lstlisting}

\textit{Judge rationale:} The agent correctly identified all given values, performed accurate arithmetic calculations step-by-step, and arrived at the correct answer of 19 matching the reference. \\
\noindent\textbf{TEI --- example 4} (gold: \texttt{159}, score: 1.0)
\textit{Query:}
\begin{lstlisting}
Nate is feeding his livestock hay. Each goat needs 5 pounds, and each sheep needs 3 pounds less than twice the amount each goat needs. If there are 15 goats and 12 sheep, how many pounds of hay does Nate need?
\end{lstlisting}

\textit{Agent output:}
\begin{lstlisting}
I need to find the total pounds of hay needed for all the livestock.

**Given information:**
- Each goat needs 5 pounds
- Each sheep needs 3 pounds less than twice the amount each goat needs
- There are 15 goats
- There are 12 sheep

**Step 1: Find how much hay each sheep needs.**
- Twice the amount each goat needs = 2 x 5 = 10 pounds
- Each sheep needs 3 pounds less than that = 10 - 3 = 7 pounds

**Step 2: Calculate total hay for goats.**
- 15 goats x 5 pounds per goat = 75 pounds

**Step 3: Calculate total hay for sheep.**
- 12 sheep x 7 pounds per sheep = 84 pounds

**Step 4: Calculate total hay needed.**
- Total = 75 + 84 = 159 pounds

FINAL: 159
\end{lstlisting}

\textit{Judge rationale:} The agent correctly interpreted the problem, applied sound arithmetic reasoning through all steps, followed proper format, and arrived at the exact reference answer of 159 pounds. \\
\subsection{\texttt{edu\_math\_word}}
\noindent\textbf{Baseline --- example 1} (gold: \texttt{144}, score: 1.0)
\textit{Query:}
\begin{lstlisting}
There are 9 crates of 16 oranges. How many oranges total?
\end{lstlisting}

\textit{Agent output:}
\begin{lstlisting}
I need to find the total number of oranges.

Given information:
- 9 crates
- 16 oranges per crate

Calculation:
Total oranges = 9 x 16 = 144

FINAL: 144
\end{lstlisting}

\textit{Judge rationale:} The agent correctly identified the problem, applied the appropriate multiplication operation (9 x 16 = 144), and provided the exact answer matching the reference with clear reasoning and proper format \\
\noindent\textbf{Baseline --- example 2} (gold: \texttt{390}, score: 1.0)
\textit{Query:}
\begin{lstlisting}
A phone costs $450 with a $60 trade-in credit. What is the final price?
\end{lstlisting}

\textit{Agent output:}
\begin{lstlisting}
I need to find the final price of the phone after applying the trade-in credit.

Starting price: $450
Trade-in credit: $60

Final price = Starting price - Trade-in credit
Final price = $450 - $60
Final price = $390

FINAL: 390
\end{lstlisting}

\textit{Judge rationale:} The agent correctly identified the problem, applied valid arithmetic reasoning (subtraction of trade-in credit), executed the calculation accurately, and produced the correct final answer of 390 match \\
\noindent\textbf{Baseline --- example 3} (gold: \texttt{126}, score: 1.0)
\textit{Query:}
\begin{lstlisting}
A worker earns $18 per hour for 7 hours. How much is earned?
\end{lstlisting}

\textit{Agent output:}
\begin{lstlisting}
To find the total earnings, I need to multiply the hourly rate by the number of hours worked.

Hourly rate: $18
Hours worked: 7

Total earnings = $18 x 7 = $126

FINAL: 126
\end{lstlisting}

\textit{Judge rationale:} The agent correctly identified the problem as multiplication, showed clear reasoning, followed proper format, and arrived at the correct answer of 126 matching the reference. \\
\noindent\textbf{Baseline --- example 4} (gold: \texttt{220}, score: 1.0)
\textit{Query:}
\begin{lstlisting}
A bus travels 55 km/h for 4 hours. How far does it go?
\end{lstlisting}

\textit{Agent output:}
\begin{lstlisting}
I need to find the distance traveled by the bus.

Using the formula: Distance = Speed x Time

Given information:
- Speed = 55 km/h
- Time = 4 hours

Calculation:
Distance = 55 km/h x 4 hours
Distance = 220 km

FINAL: 220
\end{lstlisting}

\textit{Judge rationale:} The agent correctly identified the problem type, applied the distance formula accurately, showed clear work, and arrived at the correct answer of 220 matching the reference. \\
\noindent\textbf{TEI --- example 1} (gold: \texttt{144}, score: 1.0)
\textit{Query:}
\begin{lstlisting}
There are 9 crates of 16 oranges. How many oranges total?
\end{lstlisting}

\textit{Agent output:}
\begin{lstlisting}
I need to find the total number of oranges.

Given information:
- 9 crates
- 16 oranges per crate

Calculation:
Total oranges = 9 x 16 = 144

FINAL: 144
\end{lstlisting}

\textit{Judge rationale:} The agent correctly identified the problem, applied the appropriate multiplication operation (9 x 16 = 144), and provided the exact answer matching the reference with clear reasoning and proper format \\
\noindent\textbf{TEI --- example 2} (gold: \texttt{390}, score: 1.0)
\textit{Query:}
\begin{lstlisting}
A phone costs $450 with a $60 trade-in credit. What is the final price?
\end{lstlisting}

\textit{Agent output:}
\begin{lstlisting}
I need to find the final price of the phone after applying the trade-in credit.

Starting price: $450
Trade-in credit: $60

Final price = Starting price - Trade-in credit
Final price = $450 - $60
Final price = $390

FINAL: 390
\end{lstlisting}

\textit{Judge rationale:} The agent correctly identified the problem, applied valid arithmetic reasoning (subtraction of trade-in credit), executed the calculation accurately, and produced the correct final answer of 390 match \\
\noindent\textbf{TEI --- example 3} (gold: \texttt{126}, score: 1.0)
\textit{Query:}
\begin{lstlisting}
A worker earns $18 per hour for 7 hours. How much is earned?
\end{lstlisting}

\textit{Agent output:}
\begin{lstlisting}
To find the total earnings, I need to multiply the hourly rate by the number of hours worked.

Hourly rate: $18
Hours worked: 7

Total earnings = $18 x 7 = $126

FINAL: 126
\end{lstlisting}

\textit{Judge rationale:} The agent correctly identified the problem as multiplication, showed clear reasoning, followed proper format, and arrived at the correct answer of 126 matching the reference. \\
\noindent\textbf{TEI --- example 4} (gold: \texttt{220}, score: 1.0)
\textit{Query:}
\begin{lstlisting}
A bus travels 55 km/h for 4 hours. How far does it go?
\end{lstlisting}

\textit{Agent output:}
\begin{lstlisting}
I need to find the distance traveled by the bus.

Using the formula: Distance = Speed x Time

Given information:
- Speed = 55 km/h
- Time = 4 hours

Calculation:
Distance = 55 km/h x 4 hours
Distance = 220 km

FINAL: 220
\end{lstlisting}

\textit{Judge rationale:} The agent correctly identified the problem type, applied the distance formula accurately, showed clear work, and arrived at the correct answer of 220 matching the reference. \\
\subsection{\texttt{energy\_meter\_event}}
\noindent\textbf{Baseline --- example 1} (gold: \texttt{voltage\_anomaly}, score: 1.0)
\textit{Query:}
\begin{lstlisting}
Residential meter 4782-B registered a sustained drop to 198V at 14:23 hours, remaining below the 220V nominal specification for approximately 47 seconds before self-correction. Customer appliances may have experienced performance degradation during this period.
\end{lstlisting}

\textit{Agent output:}
\begin{lstlisting}
This event describes a temporary voltage drop below the nominal specification (220V down to 198V) that lasted approximately 47 seconds before self-correcting. The customer experienced potential appliance performance issues as a result.

This is a voltage anomaly - a deviation from normal voltage levels that affects grid stability and customer equipment operation. While it self-corrected, it represents an electrical quality issue rather than an outage (which would be a complete loss of power), tampering, overconsumption, billing dispute, or maintenance activity.

FINAL: voltage_anomaly
\end{lstlisting}

\textit{Judge rationale:} The agent correctly classified the event as 'voltage\_anomaly' matching the reference answer, provided sound reasoning distinguishing it from other event types, and delivered the answer in proper forma \\
\noindent\textbf{Baseline --- example 2} (gold: \texttt{billing\_dispute}, score: 0.0)
\textit{Query:}
\begin{lstlisting}
Account representative forwarded case where customer's January statement reflects warm-weather consumption levels despite freezing temperatures and documented furnace usage.
\end{lstlisting}

\textit{Agent output:}
\begin{lstlisting}
Let me analyze this event:

The customer is reporting that their January bill shows consumption levels typical of warm weather (low usage), but they experienced freezing temperatures and actively used their furnace. This is a discrepancy between expected consumption (should be high due to heating needs) and what the meter recorded (appears low).

This could indicate:
- A meter malfunction or failure to properly register usage
- Potential tampering with the meter to underreport consumption
- A billing error in how the meter data was processed

The key indicator here is the mismatch between documented actual usage (furnace running during freezing weather) and the meter reading/bill. This points to either a meter integrity issue or deliberate manipulation.

Given the context of a billing discrepancy tied to meter accuracy concerns, this is most consistent with a **tamper** classification, as it suggests the meter may not be functioning correctly or may have been interfered with to underreport consumption.

FINAL: tamper
\end{lstlisting}

\textit{Judge rationale:} The agent correctly formatted the response and addressed the classification task, but misclassified a billing dispute as 'tamper' through speculative reasoning that leapt beyond what the input describ \\
\noindent\textbf{Baseline --- example 3} (gold: \texttt{maintenance}, score: 1.0)
\textit{Query:}
\begin{lstlisting}
Routine vegetation management along transmission corridor from mile marker 14 to 28 needed before storm season. Tree trimming crews to be deployed.
\end{lstlisting}

\textit{Agent output:}
\begin{lstlisting}
This event describes a planned activity to maintain the transmission infrastructure by managing vegetation along a corridor before storm season. Tree trimming is a preventive maintenance operation to ensure grid reliability and safety.

FINAL: maintenance
\end{lstlisting}

\textit{Judge rationale:} The agent correctly classified the event as 'maintenance' with sound reasoning about preventive vegetation management, properly formatted output, and perfect alignment with the reference answer. \\
\noindent\textbf{Baseline --- example 4} (gold: \texttt{outage}, score: 1.0)
\textit{Query:}
\begin{lstlisting}
Apartment complex lost all electricity at 13:07 following loud transformer noise. Meter communications still active but showing no energy delivery to units.
\end{lstlisting}

\textit{Agent output:}
\begin{lstlisting}
This event describes a complete loss of electrical power to an apartment complex following an audible transformer failure (loud noise), with meters still communicating but registering zero energy delivery. This is a clear indication of a power outage affecting the distribution infrastructure.

The key indicators are:
- Loss of all electricity to multiple units
- Transformer noise (suggesting equipment failure)
- No energy delivery despite active meter communications
- Sudden onset at a specific time

This is not a tamper (no unauthorized access), not overconsumption (power is completely lost), not a voltage anomaly (complete loss, not anomalous voltage), not billing dispute (operational issue), and not routine maintenance (unexpected failure event).

FINAL: outage
\end{lstlisting}

\textit{Judge rationale:} The agent correctly classified the event as 'outage' with sound reasoning that systematically eliminated alternative classifications, though the task only required the classification itself without ex \\
\noindent\textbf{TEI --- example 1} (gold: \texttt{voltage\_anomaly}, score: 1.0)
\textit{Query:}
\begin{lstlisting}
Residential meter 4782-B registered a sustained drop to 198V at 14:23 hours, remaining below the 220V nominal specification for approximately 47 seconds before self-correction. Customer appliances may have experienced performance degradation during this period.
\end{lstlisting}

\textit{Agent output:}
\begin{lstlisting}
The event describes a sustained voltage drop below the nominal specification (198V vs 220V nominal), which is an abnormal voltage level condition. This is a voltage anomaly - specifically an undervoltage situation that affected the meter and customer appliances.

FINAL: voltage_anomaly
\end{lstlisting}

\textit{Judge rationale:} The agent correctly classified the event as 'voltage\_anomaly' matching the reference answer, provided sound reasoning about the undervoltage condition, and delivered the response in proper format with \\
\noindent\textbf{TEI --- example 2} (gold: \texttt{billing\_dispute}, score: 1.0)
\textit{Query:}
\begin{lstlisting}
Account representative forwarded case where customer's January statement reflects warm-weather consumption levels despite freezing temperatures and documented furnace usage.
\end{lstlisting}

\textit{Agent output:}
\begin{lstlisting}
The customer is reporting that their energy consumption on the January bill appears unusually low (warm-weather levels) when they actually used their furnace due to freezing temperatures. This suggests either:

1. A billing error where consumption was incorrectly recorded/calculated
2. A meter malfunction that didn't properly register the actual usage
3. A dispute about the charges shown on the statement

The core issue here is a discrepancy between expected consumption (high, due to furnace heating) and what appears on the bill (low, warm-weather levels). This is fundamentally a customer complaint about billing accuracy and charges.

FINAL: billing_dispute
\end{lstlisting}

\textit{Judge rationale:} The agent correctly classified the event as 'billing\_dispute' with sound reasoning about the consumption-billing discrepancy, though the extensive explanation exceeded what was strictly necessary for  \\
\noindent\textbf{TEI --- example 3} (gold: \texttt{maintenance}, score: 1.0)
\textit{Query:}
\begin{lstlisting}
Routine vegetation management along transmission corridor from mile marker 14 to 28 needed before storm season. Tree trimming crews to be deployed.
\end{lstlisting}

\textit{Agent output:}
\begin{lstlisting}
This describes a planned, routine maintenance activity - vegetation management and tree trimming along a transmission corridor as preventive preparation before storm season.

FINAL: maintenance
\end{lstlisting}

\textit{Judge rationale:} The agent correctly classified the event as 'maintenance', provided sound reasoning about routine preventive work, and delivered the answer in proper format matching the reference exactly. \\
\noindent\textbf{TEI --- example 4} (gold: \texttt{outage}, score: 1.0)
\textit{Query:}
\begin{lstlisting}
Apartment complex lost all electricity at 13:07 following loud transformer noise. Meter communications still active but showing no energy delivery to units.
\end{lstlisting}

\textit{Agent output:}
\begin{lstlisting}
The event describes a complete loss of electricity to an apartment complex following transformer noise, with no energy delivery to the units despite meter communications remaining active. This is a power loss/service interruption scenario.

FINAL: outage
\end{lstlisting}

\textit{Judge rationale:} The agent correctly classified the event as 'outage' matching the reference answer, provided sound reasoning about the power loss scenario, and followed the classification task format properly. \\
\subsection{\texttt{fin\_banking\_intent}}
\noindent\textbf{Baseline --- example 1} (gold: \texttt{dispute\_charge}, score: 1.0)
\textit{Query:}
\begin{lstlisting}
There's a transaction on my statement I don't recognize and want reversed.
\end{lstlisting}

\textit{Agent output:}
\begin{lstlisting}
The customer is reporting an unrecognized transaction on their statement and wants it reversed. This is a request to dispute or challenge a charge that they don't believe they authorized or recognize.

FINAL: dispute_charge
\end{lstlisting}

\textit{Judge rationale:} The agent correctly classified the intent as 'dispute\_charge', matching the reference answer exactly, with sound reasoning about the unrecognized transaction requiring reversal. \\
\noindent\textbf{Baseline --- example 2} (gold: \texttt{card\_arrival}, score: 1.0)
\textit{Query:}
\begin{lstlisting}
When will the replacement card get to my house?
\end{lstlisting}

\textit{Agent output:}
\begin{lstlisting}
The customer is asking about when they will receive a replacement card at their address. This indicates they are inquiring about the arrival status of a card.

FINAL: card_arrival
\end{lstlisting}

\textit{Judge rationale:} The agent correctly classified the intent as 'card\_arrival', provided sound reasoning about the customer's inquiry, and the output matches the reference answer perfectly. \\
\noindent\textbf{Baseline --- example 3} (gold: \texttt{card\_lost}, score: 1.0)
\textit{Query:}
\begin{lstlisting}
Someone took my wallet with the card inside.
\end{lstlisting}

\textit{Agent output:}
\begin{lstlisting}
The customer is reporting that their wallet was stolen, which included their card. This indicates the card is lost or has been taken without authorization.

FINAL: card_lost
\end{lstlisting}

\textit{Judge rationale:} The agent correctly classified the intent as 'card\_lost', provided sound reasoning about the stolen wallet scenario, and delivered the exact label format required. \\
\noindent\textbf{Baseline --- example 4} (gold: \texttt{balance\_inquiry}, score: 1.0)
\textit{Query:}
\begin{lstlisting}
What's my current balance?
\end{lstlisting}

\textit{Agent output:}
\begin{lstlisting}
The customer is asking about their current account balance. This is a straightforward balance inquiry request.

FINAL: balance_inquiry
\end{lstlisting}

\textit{Judge rationale:} The agent correctly classified the message as 'balance\_inquiry', matching the reference answer exactly, with sound reasoning and proper task execution. \\
\noindent\textbf{TEI --- example 1} (gold: \texttt{dispute\_charge}, score: 1.0)
\textit{Query:}
\begin{lstlisting}
There's a transaction on my statement I don't recognize and want reversed.
\end{lstlisting}

\textit{Agent output:}
\begin{lstlisting}
The customer is reporting an unrecognized transaction on their statement that they want reversed. This is a case of a charge they don't recognize, which falls under disputing a charge.

FINAL: dispute_charge
\end{lstlisting}

\textit{Judge rationale:} The agent correctly classified the intent as 'dispute\_charge', provided sound reasoning about the unrecognized transaction, and delivered the exact label format required, matching the reference answer \\
\noindent\textbf{TEI --- example 2} (gold: \texttt{card\_arrival}, score: 1.0)
\textit{Query:}
\begin{lstlisting}
When will the replacement card get to my house?
\end{lstlisting}

\textit{Agent output:}
\begin{lstlisting}
The customer is asking about the arrival time of a replacement card they are expecting to receive.

FINAL: card_arrival
\end{lstlisting}

\textit{Judge rationale:} The agent correctly classified the intent as 'card\_arrival' matching the reference, provided sound reasoning about the customer's question, and followed the task format properly. \\
\noindent\textbf{TEI --- example 3} (gold: \texttt{card\_lost}, score: 1.0)
\textit{Query:}
\begin{lstlisting}
Someone took my wallet with the card inside.
\end{lstlisting}

\textit{Agent output:}
\begin{lstlisting}
The customer is reporting that their wallet containing their card was taken, which means the card is lost or stolen.

FINAL: card_lost
\end{lstlisting}

\textit{Judge rationale:} The agent correctly classified the intent as 'card\_lost', provided sound reasoning about the stolen wallet scenario, and delivered the exact label matching the reference answer. \\
\noindent\textbf{TEI --- example 4} (gold: \texttt{balance\_inquiry}, score: 1.0)
\textit{Query:}
\begin{lstlisting}
What's my current balance?
\end{lstlisting}

\textit{Agent output:}
\begin{lstlisting}
FINAL: balance_inquiry
\end{lstlisting}

\textit{Judge rationale:} The agent correctly classified the balance question as 'balance\_inquiry', matching the reference answer exactly with proper format and no extraneous content. \\
\subsection{\texttt{gaming\_support}}
\noindent\textbf{Baseline --- example 1} (gold: \texttt{gameplay}, score: 1.0)
\textit{Query:}
\begin{lstlisting}
I'm trying to figure out how the crafting system works. I have all the materials in my inventory but the combine button is greyed out. What am I missing?
\end{lstlisting}

\textit{Agent output:}
\begin{lstlisting}
The player is asking about the crafting system mechanics - specifically why they can't combine materials even though they have them in inventory. This is a question about how game systems work and what might be preventing an action.

This could be either a bug (if the combine button should be available but isn't) or a gameplay question (if there's a mechanic they don't understand). However, the core issue is that they're trying to understand game mechanics and troubleshoot why a feature isn't working as expected. This falls under gameplay support, as it relates to understanding and using game systems.

FINAL: gameplay
\end{lstlisting}

\textit{Judge rationale:} The agent correctly classified the ticket as 'gameplay', provided sound reasoning about the mechanics question, properly formatted the output with the required label, and the answer matches the refere \\
\noindent\textbf{Baseline --- example 2} (gold: \texttt{toxicity\_report}, score: 1.0)
\textit{Query:}
\begin{lstlisting}
There was someone in my squad earlier using a racist username and posting hate symbols in the team chat. Their tag was something like SS_Commander and they were making everyone uncomfortable.
\end{lstlisting}

\textit{Agent output:}
\begin{lstlisting}
This ticket describes a player reporting another player's offensive behavior, including a racist username and hate speech in team chat. This is clearly a report about player conduct and community standards violations.

FINAL: toxicity_report
\end{lstlisting}

\textit{Judge rationale:} The agent correctly identified the ticket as a toxicity report, provided sound reasoning about offensive behavior and community violations, and the final classification matches the reference answer ex \\
\noindent\textbf{Baseline --- example 3} (gold: \texttt{connectivity}, score: 1.0)
\textit{Query:}
\begin{lstlisting}
Experiencing massive rubber-banding where my character teleports backward every few seconds. Makes it impossible to aim or move properly in any match.
\end{lstlisting}

\textit{Agent output:}
\begin{lstlisting}
The player is describing a technical issue where their character is experiencing rubber-banding (teleporting backward repeatedly), which prevents normal gameplay functionality. This is a network-related problem where the client and server are getting out of sync, causing the character position to jump around.

This is a connectivity issue - the rubber-banding effect is typically caused by lag, packet loss, or network synchronization problems between the player's client and the game server.

FINAL: connectivity
\end{lstlisting}

\textit{Judge rationale:} The agent correctly identified rubber-banding as a connectivity issue, provided sound technical reasoning about network synchronization problems, and delivered the exact classification matching the re \\
\noindent\textbf{Baseline --- example 4} (gold: \texttt{billing}, score: 1.0)
\textit{Query:}
\begin{lstlisting}
The promotional discount code wasn't applied to my order but I was charged the full price. Can you honor the advertised discount?
\end{lstlisting}

\textit{Agent output:}
\begin{lstlisting}
The player is reporting that a promotional discount code failed to apply during checkout, resulting in being charged full price instead of the discounted amount. This is a financial transaction issue related to payment and pricing.

FINAL: billing
\end{lstlisting}

\textit{Judge rationale:} The agent correctly classified the ticket as 'billing', provided sound reasoning about the financial/pricing nature of the issue, and the output matches the reference answer perfectly. \\
\noindent\textbf{TEI --- example 1} (gold: \texttt{gameplay}, score: 1.0)
\textit{Query:}
\begin{lstlisting}
I'm trying to figure out how the crafting system works. I have all the materials in my inventory but the combine button is greyed out. What am I missing?
\end{lstlisting}

\textit{Agent output:}
\begin{lstlisting}
The player is asking about how the crafting system works and why they can't combine materials despite having them. This is a question about understanding game mechanics and how to use a feature.

FINAL: gameplay
\end{lstlisting}

\textit{Judge rationale:} The agent correctly classified the ticket as 'gameplay', provided sound reasoning about game mechanics, and the final answer perfectly matches the reference. \\
\noindent\textbf{TEI --- example 2} (gold: \texttt{toxicity\_report}, score: 1.0)
\textit{Query:}
\begin{lstlisting}
There was someone in my squad earlier using a racist username and posting hate symbols in the team chat. Their tag was something like SS_Commander and they were making everyone uncomfortable.
\end{lstlisting}

\textit{Agent output:}
\begin{lstlisting}
This is a report of abusive behavior and harassment by another player, including use of hate symbols and racist conduct in team chat. This falls under toxicity_report.

FINAL: toxicity_report
\end{lstlisting}

\textit{Judge rationale:} The agent correctly classified the ticket as 'toxicity\_report', provided sound reasoning about hate symbols and racist conduct, and the output matches the reference answer exactly. \\
\noindent\textbf{TEI --- example 3} (gold: \texttt{connectivity}, score: 1.0)
\textit{Query:}
\begin{lstlisting}
Experiencing massive rubber-banding where my character teleports backward every few seconds. Makes it impossible to aim or move properly in any match.
\end{lstlisting}

\textit{Agent output:}
\begin{lstlisting}
The player is describing rubber-banding and character teleportation, which are classic symptoms of network lag and position desynchronization between the client and server. This is a connectivity issue, not a local performance bug.

FINAL: connectivity
\end{lstlisting}

\textit{Judge rationale:} The agent correctly classified the ticket as 'connectivity', provided sound reasoning about network lag and desynchronization, and the output matches the reference answer perfectly. \\
\noindent\textbf{TEI --- example 4} (gold: \texttt{billing}, score: 1.0)
\textit{Query:}
\begin{lstlisting}
The promotional discount code wasn't applied to my order but I was charged the full price. Can you honor the advertised discount?
\end{lstlisting}

\textit{Agent output:}
\begin{lstlisting}
The player is reporting that a promotional discount code failed to apply to their purchase, resulting in being charged full price instead of the discounted amount. This is a payment/billing issue involving a discount that should have been applied but wasn't.

FINAL: billing
\end{lstlisting}

\textit{Judge rationale:} The agent correctly classified the ticket as 'billing', provided sound reasoning about the discount/payment issue, and the final answer matches the reference exactly. \\
\subsection{\texttt{gov\_service\_router}}
\noindent\textbf{Baseline --- example 1} (gold: \texttt{infrastructure}, score: 1.0)
\textit{Query:}
\begin{lstlisting}
Can someone please come repair the sidewalk in front of 342 Washington Street? Several concrete slabs are cracked and lifted, creating a tripping hazard for people walking by.
\end{lstlisting}

\textit{Agent output:}
\begin{lstlisting}
This is a request to repair public infrastructure (a sidewalk) that poses a safety hazard. The citizen is reporting a maintenance issue with municipal property that needs attention from the appropriate public works or infrastructure department.

FINAL: infrastructure
\end{lstlisting}

\textit{Judge rationale:} The agent correctly routed the sidewalk repair request to 'infrastructure' with sound reasoning about public works maintenance, matching the reference answer exactly. \\
\noindent\textbf{Baseline --- example 2} (gold: \texttt{records}, score: 1.0)
\textit{Query:}
\begin{lstlisting}
I'm looking for the official paperwork showing when my son was born here in 2015, as his school is requesting it.
\end{lstlisting}

\textit{Agent output:}
\begin{lstlisting}
You need an official birth certificate for your son. This is a vital record that schools commonly request for enrollment verification.

This request should be routed to the government department that handles vital records and document issuance.

FINAL: records
\end{lstlisting}

\textit{Judge rationale:} The agent correctly identified the request as needing vital records/birth certificate and routed it to 'records' matching the reference answer exactly, with sound reasoning about the document type and \\
\noindent\textbf{Baseline --- example 3} (gold: \texttt{complaint}, score: 1.0)
\textit{Query:}
\begin{lstlisting}
The staff member I spoke with yesterday was incredibly rude and dismissive when I asked a simple question about my application status. I expect better treatment from public servants.
\end{lstlisting}

\textit{Agent output:}
\begin{lstlisting}
I understand your frustration. Being treated dismissively by a government staff member is a legitimate concern, and you have every right to expect professional and courteous service.

This is a complaint about the quality of service you received from a government employee. The appropriate channel for this is the complaint/feedback system, which allows citizens to formally document service issues and helps departments improve their staff conduct and customer service standards.

FINAL: complaint
\end{lstlisting}

\textit{Judge rationale:} The agent correctly identified this as a complaint, provided sound reasoning about why it's a complaint rather than a service request, and delivered the exact answer matching the reference, though wit \\
\noindent\textbf{Baseline --- example 4} (gold: \texttt{taxation}, score: 1.0)
\textit{Query:}
\begin{lstlisting}
I made charitable donations throughout the year and want to make sure I'm claiming them correctly to reduce what I owe.
\end{lstlisting}

\textit{Agent output:}
\begin{lstlisting}
The citizen is asking about how to properly claim charitable donations on their tax return to minimize their tax liability. This is clearly a taxation matter, as it involves:
- Tax deductions
- Reducing tax owed
- Proper tax filing procedures

FINAL: taxation
\end{lstlisting}

\textit{Judge rationale:} The agent correctly routed to 'taxation' with sound reasoning about tax deductions and liability, perfectly matching the reference answer with proper format. \\
\noindent\textbf{TEI --- example 1} (gold: \texttt{infrastructure}, score: 1.0)
\textit{Query:}
\begin{lstlisting}
Can someone please come repair the sidewalk in front of 342 Washington Street? Several concrete slabs are cracked and lifted, creating a tripping hazard for people walking by.
\end{lstlisting}

\textit{Agent output:}
\begin{lstlisting}
This request is about infrastructure maintenance and repair - specifically a public sidewalk that needs repair due to safety hazards.

FINAL: infrastructure
\end{lstlisting}

\textit{Judge rationale:} The agent correctly identified the department as 'infrastructure', provided sound reasoning about sidewalk maintenance being an infrastructure issue, and delivered the answer in the expected format ma \\
\noindent\textbf{TEI --- example 2} (gold: \texttt{records}, score: 1.0)
\textit{Query:}
\begin{lstlisting}
I'm looking for the official paperwork showing when my son was born here in 2015, as his school is requesting it.
\end{lstlisting}

\textit{Agent output:}
\begin{lstlisting}
FINAL: records
\end{lstlisting}

\textit{Judge rationale:} The agent correctly routed the birth certificate request to 'records' department, matching the reference answer exactly with proper format and no extraneous content. \\
\noindent\textbf{TEI --- example 3} (gold: \texttt{complaint}, score: 1.0)
\textit{Query:}
\begin{lstlisting}
The staff member I spoke with yesterday was incredibly rude and dismissive when I asked a simple question about my application status. I expect better treatment from public servants.
\end{lstlisting}

\textit{Agent output:}
\begin{lstlisting}
This is a complaint about the quality of service and treatment received from a government employee.

FINAL: complaint
\end{lstlisting}

\textit{Judge rationale:} The agent correctly identified the input as a complaint, provided sound reasoning about service quality issues, and delivered the exact label matching the reference answer. \\
\noindent\textbf{TEI --- example 4} (gold: \texttt{taxation}, score: 1.0)
\textit{Query:}
\begin{lstlisting}
I made charitable donations throughout the year and want to make sure I'm claiming them correctly to reduce what I owe.
\end{lstlisting}

\textit{Agent output:}
\begin{lstlisting}
FINAL: taxation
\end{lstlisting}

\textit{Judge rationale:} The agent correctly routed a tax-related query about charitable donation deductions to the taxation department, matching the reference answer exactly with proper formatting. \\
\subsection{\texttt{health\_triage}}
\noindent\textbf{Baseline --- example 1} (gold: \texttt{routine}, score: 1.0)
\textit{Query:}
\begin{lstlisting}
I want a dental cleaning appointment.
\end{lstlisting}

\textit{Agent output:}
\begin{lstlisting}
This is a straightforward request for preventive dental care. A routine dental cleaning is a standard, non-urgent procedure that helps maintain oral health and is typically scheduled in advance.

FINAL: routine
\end{lstlisting}

\textit{Judge rationale:} The agent correctly classified the dental cleaning request as 'routine' matching the reference answer, provided sound reasoning about preventive care being non-urgent, and followed the task format app \\
\noindent\textbf{Baseline --- example 2} (gold: \texttt{emergency}, score: 1.0)
\textit{Query:}
\begin{lstlisting}
Unconscious friend who is not responding to voice or touch.
\end{lstlisting}

\textit{Agent output:}
\begin{lstlisting}
This is a critical situation requiring immediate medical intervention.

An unconscious person who is not responding to voice or touch indicates:
- Altered level of consciousness (potentially severe)
- Unknown underlying cause (could be stroke, head injury, overdose, severe infection, cardiac event, etc.)
- Risk of airway compromise
- Potential for rapid deterioration

This patient needs emergency medical services immediately, including:
- Activation of emergency services (911 or local equivalent)
- Assessment of airway, breathing, and circulation
- Rapid diagnostic evaluation
- Potential life-saving interventions

FINAL: emergency
\end{lstlisting}

\textit{Judge rationale:} The agent correctly classified the situation as 'emergency', provided sound clinical reasoning for the urgency level, and fully addressed the triage classification task with accurate execution and com \\
\noindent\textbf{Baseline --- example 3} (gold: \texttt{emergency}, score: 1.0)
\textit{Query:}
\begin{lstlisting}
Heavy bleeding that soaks through a bandage in minutes.
\end{lstlisting}

\textit{Agent output:}
\begin{lstlisting}
This patient is experiencing heavy, rapid bleeding that saturates a bandage within minutes. This indicates:

- Significant blood loss occurring quickly
- Potential for hemorrhagic shock if not controlled
- Possible injury to a major blood vessel
- Risk of life-threatening complications without immediate intervention

This requires immediate medical evaluation and intervention to control bleeding, assess the wound, and prevent shock.

FINAL: emergency
\end{lstlisting}

\textit{Judge rationale:} The agent correctly classified the situation as 'emergency' matching the reference answer, provided sound medical reasoning for the urgency level, and followed the task format appropriately with only  \\
\noindent\textbf{Baseline --- example 4} (gold: \texttt{self\_care}, score: 1.0)
\textit{Query:}
\begin{lstlisting}
Minor bruise on my knee from bumping a table.
\end{lstlisting}

\textit{Agent output:}
\begin{lstlisting}
This is a minor injury with no signs of serious damage. A small bruise from a table bump typically resolves on its own with basic home care like ice, rest, and elevation if needed. No medical evaluation is required unless swelling, severe pain, or other concerning symptoms develop.

FINAL: self_care
\end{lstlisting}

\textit{Judge rationale:} The agent correctly classified the minor bruise as 'self\_care', provided sound reasoning about why medical attention is unnecessary, and fully matched the reference answer with proper format. \\
\noindent\textbf{TEI --- example 1} (gold: \texttt{routine}, score: 1.0)
\textit{Query:}
\begin{lstlisting}
I want a dental cleaning appointment.
\end{lstlisting}

\textit{Agent output:}
\begin{lstlisting}
FINAL: routine
\end{lstlisting}

\textit{Judge rationale:} The agent correctly classified a dental cleaning request as 'routine' triage urgency, matching the reference answer with proper format and sound implicit reasoning about non-urgent preventive care. \\
\noindent\textbf{TEI --- example 2} (gold: \texttt{emergency}, score: 1.0)
\textit{Query:}
\begin{lstlisting}
Unconscious friend who is not responding to voice or touch.
\end{lstlisting}

\textit{Agent output:}
\begin{lstlisting}
An unconscious person who is not responding to voice or touch is a life-threatening emergency requiring immediate medical intervention. This could indicate severe conditions such as stroke, head injury, cardiac event, severe intoxication, or other critical medical emergencies.

FINAL: emergency
\end{lstlisting}

\textit{Judge rationale:} The agent correctly classified the situation as 'emergency', provided sound medical reasoning for the urgency level, and fully aligned with the reference answer without any errors or fabrications. \\
\noindent\textbf{TEI --- example 3} (gold: \texttt{emergency}, score: 1.0)
\textit{Query:}
\begin{lstlisting}
Heavy bleeding that soaks through a bandage in minutes.
\end{lstlisting}

\textit{Agent output:}
\begin{lstlisting}
This is a life-threatening condition requiring immediate medical attention. Heavy bleeding that saturates a bandage within minutes indicates significant blood loss and potential for rapid deterioration.

FINAL: emergency
\end{lstlisting}

\textit{Judge rationale:} The agent correctly classified the situation as 'emergency' matching the reference answer, provided sound medical reasoning about blood loss severity, and followed the required format perfectly. \\
\noindent\textbf{TEI --- example 4} (gold: \texttt{self\_care}, score: 1.0)
\textit{Query:}
\begin{lstlisting}
Minor bruise on my knee from bumping a table.
\end{lstlisting}

\textit{Agent output:}
\begin{lstlisting}
FINAL: self_care
\end{lstlisting}

\textit{Judge rationale:} The agent correctly classified a minor bruise as 'self\_care', matching the reference answer exactly with proper format and no extraneous content. \\
\subsection{\texttt{hospitality\_review\_stars}}
\noindent\textbf{Baseline --- example 1} (gold: \texttt{5 stars}, score: 1.0)
\textit{Query:}
\begin{lstlisting}
I have been coming to this particular Filiberto's for a while now and the food and friendly service are consistently top notch. The bar is stocked with good tequila and they usually have something going on like karaoke in the evening to keep things interesting. It fills an interesting niche of taco shop fast food along with a nice bar that I didn't think would gel together but surprisingly does. Food and drink specials for the most part are also pretty great as well.
\end{lstlisting}

\textit{Agent output:}
\begin{lstlisting}
This review is clearly positive and enthusiastic. The customer highlights:

- Consistent quality of food and service ("consistently top notch")
- Friendly service
- Well-stocked bar with good tequila
- Entertainment options (karaoke)
- Appreciation for the unique concept (taco shop + bar combination)
- Good food and drink specials

The reviewer is a repeat customer ("coming to this particular Filiberto's for a while now") and expresses satisfaction with multiple aspects of the establishment. There are no complaints or negative points mentioned. The tone is appreciative and recommending.

This clearly warrants a 5-star rating.

FINAL: 5 stars
\end{lstlisting}

\textit{Judge rationale:} The agent correctly predicted 5 stars with sound reasoning that accurately identified all positive indicators in the review, addressed the exact task, and matched the reference answer perfectly. \\
\noindent\textbf{Baseline --- example 2} (gold: \texttt{4 stars}, score: 0.0)
\textit{Query:}
\begin{lstlisting}
With 500+ Reviews, Most of all it has been said but I will offer this:\nFox Concepts Restaurants consistently pulls off a memorable, pleasant experience that is tasty, full of quality and will, occasionally, have a memorable service experience. \nAll restaurants have their ups & downs, but Olive & Ivy is consistently good, innovative and the atmosphere is quite nice; inside or outside.\nDrinks are tasty, Scallopini and Cannelloni are both just Yum! Never had a bad experience here.\n\nOverall Rating: 4.5 Stars\nService: 4.1 Stars\nFood: 4.6 Stars \nQuality: 4.5 Stars\nCleanliness: 4.2 Stars\nValue: 3.9 Stars
\end{lstlisting}

\textit{Agent output:}
\begin{lstlisting}
This review is highly positive overall. The reviewer praises:

- Consistent quality and memorable experiences
- Innovative menu items (specifically praising Scallopini and Cannelloni)
- Pleasant atmosphere (both indoor and outdoor)
- Good drinks
- Never had a bad experience

The reviewer's own detailed ratings show:
- Food: 4.6 stars (excellent)
- Quality: 4.5 stars (very good)
- Service: 4.1 stars (good)
- Cleanliness: 4.2 stars (good)
- Value: 3.9 stars (acceptable, slightly lower)
- Overall stated: 4.5 stars

The only minor criticism is the value rating at 3.9, suggesting prices are somewhat high, but this doesn't diminish the overall very positive sentiment. The reviewer explicitly states they've "never had a bad experience" and consistently recommends the restaurant.

This clearly warrants a 5-star rating, as the reviewer is expressing strong satisfaction with the establishment despite the slightly lower value score.

FINAL: 5 stars
\end{lstlisting}

\textit{Judge rationale:} The agent correctly formatted the prediction and reasoned through the review systematically, but predicted 5 stars when the reference answer is 4 stars, and the reviewer's own stated overall rating wa \\
\noindent\textbf{Baseline --- example 3} (gold: \texttt{1 star}, score: 1.0)
\textit{Query:}
\begin{lstlisting}
Just sad\nOrdered online paid on line when I arrived at the pick up time they told me that their fax machine ran out of paper and they did not get my order\n\nI am now sitting in the dining room writing this review on my phone \n\nNo attempt to make this right \n\nI hope the pizza is good
\end{lstlisting}

\textit{Agent output:}
\begin{lstlisting}
This review describes a significant service failure. The customer:
- Placed and paid for an online order
- Arrived at the scheduled pickup time
- Was told the order was never received due to a technical issue (fax machine out of paper)
- Experienced no attempt by the restaurant to resolve the situation
- Is now waiting, frustrated enough to write a negative review

The customer's tone is clearly disappointed and frustrated. While they haven't explicitly stated the pizza was bad, the phrase "I hope the pizza is good" suggests skepticism and that this experience has already soured their perception. The lack of any service recovery attempt (offering to expedite the order, provide a discount, etc.) makes this worse.

This is a clear operational failure with poor customer service recovery.

FINAL: 1 star
\end{lstlisting}

\textit{Judge rationale:} The agent correctly predicted 1 star with sound reasoning about service failure, proper format, and complete alignment with the reference answer. \\
\noindent\textbf{Baseline --- example 4} (gold: \texttt{2 star}, score: 1.0)
\textit{Query:}
\begin{lstlisting}
Hmmmmm really just not sure about this place at all, went for a lunch and just wasn't impressed. I had the quesadilla with carnitas which was ok, we also got chicken wings and they were well awful! Cold, greasy and just not very good.\n\nThe search for good tex-mex goes on!
\end{lstlisting}

\textit{Agent output:}
\begin{lstlisting}
The reviewer expresses clear dissatisfaction with their experience. While they found the quesadilla "ok," they were notably disappointed with the chicken wings (described as "well awful," cold, and greasy). The overall tone is negative - they state they "just weren't impressed" and are still searching for better tex-mex elsewhere.

The presence of one acceptable dish prevents this from being a 1-star review, but the significant complaints about quality and the overall lack of satisfaction point to a below-average experience.

FINAL: 2 star
\end{lstlisting}

\textit{Judge rationale:} The agent correctly predicted '2 star' matching the reference, provided sound reasoning about mixed-but-negative sentiment, followed the required format, and delivered an accurate answer without fabri \\
\noindent\textbf{TEI --- example 1} (gold: \texttt{5 stars}, score: 1.0)
\textit{Query:}
\begin{lstlisting}
I have been coming to this particular Filiberto's for a while now and the food and friendly service are consistently top notch. The bar is stocked with good tequila and they usually have something going on like karaoke in the evening to keep things interesting. It fills an interesting niche of taco shop fast food along with a nice bar that I didn't think would gel together but surprisingly does. Food and drink specials for the most part are also pretty great as well.
\end{lstlisting}

\textit{Agent output:}
\begin{lstlisting}
This review is overwhelmingly positive. The reviewer:

- Explicitly states they are a repeat customer ("coming to this particular Filiberto's for a while now")
- Praises both food and service as "consistently top notch"
- Highlights additional positive amenities (well-stocked bar, entertainment like karaoke)
- Expresses pleasant surprise at how well the concept works ("surprisingly does")
- Appreciates the specials ("pretty great")
- Uses enthusiastic language throughout with no significant complaints or reservations

There are no negative issues mentioned, no drawbacks acknowledged, and the overall tone is one of genuine satisfaction and recommendation. This is a clear case of strong, consistent praise across multiple dimensions of the experience.

FINAL: 5 stars
\end{lstlisting}

\textit{Judge rationale:} The agent correctly predicted 5 stars with sound reasoning that systematically identified positive sentiment indicators, properly formatted the answer, and matched the reference exactly. \\
\noindent\textbf{TEI --- example 2} (gold: \texttt{4 stars}, score: 1.0)
\textit{Query:}
\begin{lstlisting}
With 500+ Reviews, Most of all it has been said but I will offer this:\nFox Concepts Restaurants consistently pulls off a memorable, pleasant experience that is tasty, full of quality and will, occasionally, have a memorable service experience. \nAll restaurants have their ups & downs, but Olive & Ivy is consistently good, innovative and the atmosphere is quite nice; inside or outside.\nDrinks are tasty, Scallopini and Cannelloni are both just Yum! Never had a bad experience here.\n\nOverall Rating: 4.5 Stars\nService: 4.1 Stars\nFood: 4.6 Stars \nQuality: 4.5 Stars\nCleanliness: 4.2 Stars\nValue: 3.9 Stars
\end{lstlisting}

\textit{Agent output:}
\begin{lstlisting}
**Reasoning:**

This review is overwhelmingly positive with strong, consistent praise across multiple dimensions:

- **Explicit positive language**: "memorable, pleasant experience," "tasty, full of quality," "consistently good, innovative," "quite nice," "Yum!"
- **No significant complaints**: The reviewer acknowledges that "all restaurants have their ups & downs" but explicitly states "Never had a bad experience here"
- **Strong food praise**: Multiple dishes called out as excellent (Scallopini and Cannelloni described as "Yum!")
- **Atmosphere and drinks**: Both praised as tasty and nice
- **Tone**: Enthusiastic and recommending, with clear intent to return
- **Numerical ratings provided**: All sub-ratings are 3.9-4.6 stars, with an overall 4.5 rating

The reviewer acknowledges minor variations in service quality ("occasionally" memorable) but this is framed as normal variation, not a significant detraction. The overall sentiment is strongly positive with no deal-breaker issues, warnings, or substantial disappointments.

This is a clear 4-star review: mostly positive with minor acknowledgments of normal restaurant variability, but strong praise dominating and no hesitation about returning.

FINAL: 4 stars
\end{lstlisting}

\textit{Judge rationale:} The agent correctly predicted 4 stars matching the reference, with sound reasoning analyzing sentiment, explicit ratings, and review tone, properly formatted output, and no fabrication. \\
\noindent\textbf{TEI --- example 3} (gold: \texttt{1 star}, score: 1.0)
\textit{Query:}
\begin{lstlisting}
Just sad\nOrdered online paid on line when I arrived at the pick up time they told me that their fax machine ran out of paper and they did not get my order\n\nI am now sitting in the dining room writing this review on my phone \n\nNo attempt to make this right \n\nI hope the pizza is good
\end{lstlisting}

\textit{Agent output:}
\begin{lstlisting}
**Analysis:**

This review describes a fundamental service failure:
- Customer paid in advance for an online order
- Restaurant failed to receive the order due to a technical issue (fax machine out of paper)
- Customer arrived at the agreed pickup time to find no order prepared
- **Critical issue**: "No attempt to make this right" - the restaurant did not offer service recovery, compensation, or solutions
- The reviewer is sitting in the dining room, clearly frustrated and disappointed
- Opening phrase "Just sad" sets a deeply disappointed tone
- The sarcastic closing ("I hope the pizza is good") suggests skepticism and resignation rather than optimism

This is a deal-breaker situation: the customer was let down by a core service failure (order not received/prepared), and the restaurant compounded the problem by failing to attempt any recovery or compensation. The lack of effort to resolve the issue is explicitly called out as a major problem.

While the reviewer hasn't explicitly said "never again," the combination of being stood up on a paid order, no service recovery attempt, and the overall "just sad" tone indicates a significantly negative experience that would deter future visits.

FINAL: 1 star
\end{lstlisting}

\textit{Judge rationale:} The agent correctly predicted 1 star with sound reasoning about service failure, lack of recovery attempt, and negative tone, perfectly matching the reference answer. \\
\noindent\textbf{TEI --- example 4} (gold: \texttt{2 star}, score: 1.0)
\textit{Query:}
\begin{lstlisting}
Hmmmmm really just not sure about this place at all, went for a lunch and just wasn't impressed. I had the quesadilla with carnitas which was ok, we also got chicken wings and they were well awful! Cold, greasy and just not very good.\n\nThe search for good tex-mex goes on!
\end{lstlisting}

\textit{Agent output:}
\begin{lstlisting}
**Analysis:**

The reviewer expresses clear disappointment with their experience:

- **Overall tone**: Uncertain and unimpressed ("really just not sure about this place at all," "wasn't impressed")
- **Food quality issues**: 
  - Quesadilla described as merely "ok" (lukewarm assessment)
  - Chicken wings explicitly criticized as "well awful" with specific problems (cold, greasy, poor quality)
- **No redeeming elements**: No positive comments about service, atmosphere, or any menu items
- **Conclusion**: Actively searching elsewhere ("The search for good tex-mex goes on!") suggests they won't return
- **Ratio**: Multiple substantive complaints with no genuine positives to balance them

This is a predominantly negative experience where problems outweigh any positives. The reviewer had at least two dishes, and both disappointed (one mediocre, one explicitly bad). The tone is disappointed rather than enthusiastic, and they're clearly moving on to find better options elsewhere.

This does not rise to 1 star (no explicit "never again" or strong warning language), but it clearly falls below 3 stars (which requires balanced positives and negatives or mostly positive with minor drawbacks).

FINAL: 2 star
\end{lstlisting}

\textit{Judge rationale:} The agent correctly predicted '2 star' matching the reference, with sound reasoning analyzing tone and food quality, proper format, and no fabrication. \\
\subsection{\texttt{hr\_resume\_fit}}
\noindent\textbf{Baseline --- example 1} (gold: \texttt{overqualified}, score: 1.0)
\textit{Query:}
\begin{lstlisting}
Candidate has 15 years leading enterprise sales teams with P&L responsibility for $50M+ organizations. Our sales associate role is an individual contributor position focused on inside sales to SMB clients. Worried they'll be bored within months and leave for a leadership role elsewhere.
\end{lstlisting}

\textit{Agent output:}
\begin{lstlisting}
**Reasoning:**

This candidate has extensive executive-level experience managing large enterprise sales organizations with significant P&L responsibility. The role in question is an individual contributor position at the SMB level--a substantial step down in scope, complexity, and seniority.

Key concerns:
- **Significant overqualification**: 15 years of leadership experience managing $50M+ portfolios far exceeds what an inside sales associate role requires
- **Flight risk**: The concern about boredom and departure is well-founded. Candidates at this level typically seek growth, autonomy, and strategic responsibility
- **Role mismatch**: An IC position lacks the leadership, decision-making authority, and scope this candidate is accustomed to
- **Retention likelihood**: Low probability of long-term retention once the candidate realizes the limited growth trajectory

While the candidate clearly has strong sales skills, the role doesn't align with their experience level or likely career expectations.

FINAL: overqualified
\end{lstlisting}

\textit{Judge rationale:} The agent correctly classified the candidate as 'overqualified' with thorough, logically sound reasoning that directly addresses the task requirements and matches the reference answer perfectly. \\
\noindent\textbf{Baseline --- example 2} (gold: \texttt{underqualified}, score: 1.0)
\textit{Query:}
\begin{lstlisting}
Candidate's background is entirely in retail management at small stores. The position requires a licensed attorney with 10+ years of corporate law experience and expertise in M&A transactions. No law degree or legal credentials present.
\end{lstlisting}

\textit{Agent output:}
\begin{lstlisting}
This candidate lacks the fundamental qualifications for the role. They have retail management experience but the position requires:
- Law degree and active bar license (not present)
- 10+ years of corporate law experience (candidate has none)
- Specialized M&A transaction expertise (not applicable to retail background)

The candidate is missing both the required credentials and relevant professional experience entirely. This is a clear mismatch.

FINAL: underqualified
\end{lstlisting}

\textit{Judge rationale:} The agent correctly classified the candidate as 'underqualified' with sound reasoning that systematically identified all missing qualifications (law degree, bar license, corporate law experience, M\&A  \\
\noindent\textbf{Baseline --- example 3} (gold: \texttt{wrong\_role}, score: 1.0)
\textit{Query:}
\begin{lstlisting}
This person has deep expertise in financial auditing and compliance, which is valuable but not aligned with our finance operations analyst position. We need someone to handle FP&A, budgeting forecasts, and operational metrics rather than regulatory reporting.
\end{lstlisting}

\textit{Agent output:}
\begin{lstlisting}
**Reasoning:**
The candidate possesses strong technical expertise in financial auditing and compliance--valuable skills in finance. However, the role requires FP&A, budgeting, forecasting, and operational metrics analysis. These are distinct skill sets. While there's overlap in financial acumen, auditing/compliance focuses on historical accuracy and regulatory requirements, whereas FP&A emphasizes forward-looking analysis, strategic planning, and business operations metrics.

The candidate's background suggests they may lack direct experience in the core competencies needed (forecasting models, operational analytics, budget planning). This is a skills mismatch rather than a cultural or qualification level issue.

FINAL: wrong_role
\end{lstlisting}

\textit{Judge rationale:} The agent correctly classified as 'wrong\_role' with sound reasoning distinguishing auditing/compliance from FP\&A skills, proper format, and perfect match to reference answer. \\
\noindent\textbf{Baseline --- example 4} (gold: \texttt{strong\_fit}, score: 1.0)
\textit{Query:}
\begin{lstlisting}
Ideal background for our content marketing manager opening. Grew organic traffic by 300% over two years through SEO optimization and strategic content planning. Writing samples demonstrate versatility across technical and consumer audiences. Analytics fluency evident in discussion of attribution modeling and conversion optimization. Team collaboration stories highlighted empathy and cross-functional partnership. Genuine excitement about our brand and growth trajectory.
\end{lstlisting}

\textit{Agent output:}
\begin{lstlisting}
This candidate demonstrates strong alignment with the content marketing manager role:

**Strengths:**
- Proven track record of measurable impact (300% organic traffic growth)
- Technical competency in SEO, analytics, and attribution modeling
- Demonstrated versatility in writing across different audiences
- Evidence of soft skills (empathy, collaboration) essential for cross-functional work
- Authentic enthusiasm for the company and mission

**Assessment:**
The candidate shows expertise directly relevant to the role, concrete results that validate their capabilities, and cultural fit indicators. No red flags regarding overqualification, underqualification, or misalignment with the position.

FINAL: strong_fit
\end{lstlisting}

\textit{Judge rationale:} The agent correctly classified the candidate as 'strong\_fit' with well-structured supporting reasoning that directly addresses the task requirements, though the additional explanation exceeds minimal  \\
\noindent\textbf{TEI --- example 1} (gold: \texttt{overqualified}, score: 1.0)
\textit{Query:}
\begin{lstlisting}
Candidate has 15 years leading enterprise sales teams with P&L responsibility for $50M+ organizations. Our sales associate role is an individual contributor position focused on inside sales to SMB clients. Worried they'll be bored within months and leave for a leadership role elsewhere.
\end{lstlisting}

\textit{Agent output:}
\begin{lstlisting}
The candidate's extensive leadership experience (15 years managing enterprise teams with significant P&L responsibility) far exceeds the requirements of an individual contributor sales associate role focused on SMB inside sales. This represents a multi-level seniority mismatch in the same functional area (sales). While the candidate clearly has strong sales capabilities, their career trajectory and experience level suggest they would likely find this role understimulating and seek advancement or opportunities elsewhere relatively quickly, creating retention risk.

FINAL: overqualified
\end{lstlisting}

\textit{Judge rationale:} The agent correctly classified the candidate as 'overqualified' with sound reasoning about the seniority mismatch and retention risk, fully addressing the task requirements. \\
\noindent\textbf{TEI --- example 2} (gold: \texttt{underqualified}, score: 1.0)
\textit{Query:}
\begin{lstlisting}
Candidate's background is entirely in retail management at small stores. The position requires a licensed attorney with 10+ years of corporate law experience and expertise in M&A transactions. No law degree or legal credentials present.
\end{lstlisting}

\textit{Agent output:}
\begin{lstlisting}
The candidate lacks the essential qualifications for this role. They have retail management experience but no law degree, legal credentials, or corporate law background. The position requires a licensed attorney with specialized M&A experience--a fundamentally different functional area and professional discipline.

FINAL: underqualified
\end{lstlisting}

\textit{Judge rationale:} The agent correctly classified the candidate as 'underqualified' with sound reasoning about the mismatch between retail management background and attorney requirements, perfectly matching the referenc \\
\noindent\textbf{TEI --- example 3} (gold: \texttt{wrong\_role}, score: 1.0)
\textit{Query:}
\begin{lstlisting}
This person has deep expertise in financial auditing and compliance, which is valuable but not aligned with our finance operations analyst position. We need someone to handle FP&A, budgeting forecasts, and operational metrics rather than regulatory reporting.
\end{lstlisting}

\textit{Agent output:}
\begin{lstlisting}
The candidate's core expertise is in financial auditing and compliance--a different functional area than the finance operations analyst role, which requires FP&A, budgeting, forecasting, and operational metrics focus. While both are finance-related, they represent distinct career paths and skill sets. The candidate's specialization doesn't align with the position's core functions.

FINAL: wrong_role
\end{lstlisting}

\textit{Judge rationale:} The agent correctly classified the candidate as 'wrong\_role' with sound reasoning that accurately explains the mismatch between auditing/compliance expertise and the FP\&A/operations analyst requiremen \\
\noindent\textbf{TEI --- example 4} (gold: \texttt{strong\_fit}, score: 1.0)
\textit{Query:}
\begin{lstlisting}
Ideal background for our content marketing manager opening. Grew organic traffic by 300% over two years through SEO optimization and strategic content planning. Writing samples demonstrate versatility across technical and consumer audiences. Analytics fluency evident in discussion of attribution modeling and conversion optimization. Team collaboration stories highlighted empathy and cross-functional partnership. Genuine excitement about our brand and growth trajectory.
\end{lstlisting}

\textit{Agent output:}
\begin{lstlisting}
This candidate demonstrates strong alignment with the content marketing manager role:

- **Core skills match**: Proven SEO expertise, strategic content planning, and analytics proficiency directly address the position's requirements
- **Measurable impact**: 300% organic traffic growth shows concrete results in the functional area
- **Communication ability**: Writing samples across diverse audiences indicate versatility needed for the role
- **Technical competency**: Attribution modeling and conversion optimization knowledge demonstrates analytical depth
- **Behavioral indicators**: Team collaboration stories and empathy suggest positive cultural fit
- **Engagement**: Genuine enthusiasm about the brand and growth trajectory indicates motivation and alignment

No significant gaps, concerns, or red flags are present. The candidate's experience level appears appropriately matched to the role (not overqualified or underqualified).

FINAL: strong_fit
\end{lstlisting}

\textit{Judge rationale:} The agent correctly classified the candidate as 'strong\_fit' matching the reference answer, provided well-structured reasoning that directly addresses the hiring-stage classification task, followed pr \\
\subsection{\texttt{insurance\_claim\_type}}
\noindent\textbf{Baseline --- example 1} (gold: \texttt{theft}, score: 1.0)
\textit{Query:}
\begin{lstlisting}
Someone broke into my car last night while it was parked on the street and took my laptop, GPS device, and about $200 in cash from the center console. The driver's side window was smashed.
\end{lstlisting}

\textit{Agent output:}
\begin{lstlisting}
This claim involves items being taken from a vehicle without authorization. While there is property damage (the broken window), the primary loss is the theft of personal items from inside the car.

The key elements are:
- Unauthorized entry into the vehicle
- Stolen items (laptop, GPS, cash)
- This is a theft claim rather than an auto collision claim (which would involve vehicle damage from an accident)

FINAL: theft
\end{lstlisting}

\textit{Judge rationale:} The agent correctly classified the claim as 'theft' with sound reasoning about unauthorized entry and stolen items, though the task only required the classification itself without extended explanation \\
\noindent\textbf{Baseline --- example 2} (gold: \texttt{liability}, score: 1.0)
\textit{Query:}
\begin{lstlisting}
While hosting a dinner party, my chandelier fell from the ceiling and injured one of my guests, causing a head laceration requiring surgery. The guest is pursuing compensation for their injuries.
\end{lstlisting}

\textit{Agent output:}
\begin{lstlisting}
This claim involves a guest being injured at the policyholder's property due to a falling chandelier. The guest is seeking compensation for medical expenses and damages resulting from the injury.

This is a liability claim because:
- It involves bodily injury to a third party (the guest)
- The policyholder may be held legally responsible for the injury
- The guest is pursuing compensation against the policyholder
- This falls under premises liability coverage

While the guest will have medical expenses, the claim itself is being routed based on the liability exposure and the third-party injury claim.

FINAL: liability
\end{lstlisting}

\textit{Judge rationale:} The agent correctly classified the claim as 'liability' with sound reasoning about third-party injury and premises liability, though the task only required classification without explanation. \\
\noindent\textbf{Baseline --- example 3} (gold: \texttt{medical}, score: 1.0)
\textit{Query:}
\begin{lstlisting}
Submitted bills for ambulance transport, ER visit, CT scan, and follow-up appointments related to head trauma sustained during a fall at a retail store.
\end{lstlisting}

\textit{Agent output:}
\begin{lstlisting}
This claim involves medical expenses resulting from an injury sustained at a retail location. The submitted bills cover ambulance transport, emergency room visit, CT scan, and follow-up medical appointments - all medical services related to treating head trauma from a fall.

While the fall occurred at a retail store (which could potentially involve a liability component), the claim itself is being submitted for the medical costs incurred for treatment and care.

FINAL: medical
\end{lstlisting}

\textit{Judge rationale:} The agent correctly classified the claim as 'medical', provided sound reasoning distinguishing medical costs from liability aspects, and delivered the exact answer matching the reference with proper f \\
\noindent\textbf{Baseline --- example 4} (gold: \texttt{auto\_collision}, score: 1.0)
\textit{Query:}
\begin{lstlisting}
I was stopped at a red light when the vehicle behind me failed to brake and rear-ended my car. The trunk won't close properly now, the rear bumper is cracked, and the taillights on both sides are broken.
\end{lstlisting}

\textit{Agent output:}
\begin{lstlisting}
This claim describes a vehicle collision where the claimant's car was struck from behind while stationary at a traffic light. The damage includes structural damage (trunk, bumper) and broken taillights resulting from the impact.

This is clearly an auto collision claim, as it involves damage to a vehicle caused by a collision with another vehicle.

FINAL: auto_collision
\end{lstlisting}

\textit{Judge rationale:} The agent correctly classified the claim as auto\_collision with sound reasoning about the rear-end collision scenario, though the task only required the classification label itself. \\
\noindent\textbf{TEI --- example 1} (gold: \texttt{theft}, score: 1.0)
\textit{Query:}
\begin{lstlisting}
Someone broke into my car last night while it was parked on the street and took my laptop, GPS device, and about $200 in cash from the center console. The driver's side window was smashed.
\end{lstlisting}

\textit{Agent output:}
\begin{lstlisting}
This claim involves stolen property taken from a vehicle due to criminal activity (break-in and theft). While the vehicle itself sustained damage (broken window), the primary claim is for the stolen items (laptop, GPS device, cash) resulting from a burglary.

FINAL: theft
\end{lstlisting}

\textit{Judge rationale:} The agent correctly classified the claim as 'theft' matching the reference answer, provided sound reasoning distinguishing theft from vehicle damage, and delivered the answer in proper format with app \\
\noindent\textbf{TEI --- example 2} (gold: \texttt{liability}, score: 1.0)
\textit{Query:}
\begin{lstlisting}
While hosting a dinner party, my chandelier fell from the ceiling and injured one of my guests, causing a head laceration requiring surgery. The guest is pursuing compensation for their injuries.
\end{lstlisting}

\textit{Agent output:}
\begin{lstlisting}
This claim involves a guest being injured at the policyholder's home due to the chandelier falling, and the guest is seeking compensation. This is a third-party injury claim where the policyholder may be held liable for the harm caused to another person.

FINAL: liability
\end{lstlisting}

\textit{Judge rationale:} The agent correctly classified the claim as 'liability' with sound reasoning about third-party injury and policyholder responsibility, fully matching the reference answer. \\
\noindent\textbf{TEI --- example 3} (gold: \texttt{medical}, score: 1.0)
\textit{Query:}
\begin{lstlisting}
Submitted bills for ambulance transport, ER visit, CT scan, and follow-up appointments related to head trauma sustained during a fall at a retail store.
\end{lstlisting}

\textit{Agent output:}
\begin{lstlisting}
This claim involves medical expenses for the policyholder's own healthcare treatment following an injury. The bills are for ambulance transport, emergency room visit, CT scan, and follow-up appointments - all medical services related to treating the head trauma.

FINAL: medical
\end{lstlisting}

\textit{Judge rationale:} The agent correctly classified the claim as 'medical' with sound reasoning about the medical services involved, properly formatted output, and perfect alignment with the reference answer. \\
\noindent\textbf{TEI --- example 4} (gold: \texttt{auto\_collision}, score: 1.0)
\textit{Query:}
\begin{lstlisting}
I was stopped at a red light when the vehicle behind me failed to brake and rear-ended my car. The trunk won't close properly now, the rear bumper is cracked, and the taillights on both sides are broken.
\end{lstlisting}

\textit{Agent output:}
\begin{lstlisting}
This claim describes damage to a vehicle resulting from a collision (rear-end accident). The damage to the trunk, bumper, and taillights are all consequences of the vehicle being struck by another car.

FINAL: auto_collision
\end{lstlisting}

\textit{Judge rationale:} The agent correctly classified the claim as auto\_collision with sound reasoning about the rear-end accident, matching the reference answer perfectly despite including explanatory context. \\
\subsection{\texttt{it\_email\_spam}}
\noindent\textbf{Baseline --- example 1} (gold: \texttt{ham}, score: 1.0)
\textit{Query:}
\begin{lstlisting}
re : park n ride linda , i talked with linda and michelle this week on this issue . what michelle wanted to know was does tw have operators who have pnr agreements . currently , tw does not have any operators with pnr agreements because they have mandatory obas on tw . with our 637 compliance filing tw is requesting that operators move their imbalance to a pnr if it is over a certin volume . also eog does pool on tw so as a shipper we could get them to enter into a pnr . they would need to have a shipper transport it to and from the pnr point . i also told michelle that all points the flow any volumes on tw are required to be an an oba . michelle said that they were able to resolve this issue with eog . thanks ramona
\end{lstlisting}

\textit{Agent output:}
\begin{lstlisting}
This email appears to be a legitimate business communication regarding operational matters.

**Reasoning:**
- The email discusses specific business topics (Park and Ride agreements, operators, compliance filings, OBA requirements)
- References real people (Linda, Michelle, Ramona) and specific companies (TW, EOG)
- Contains detailed technical/operational information relevant to energy/transportation industry
- Uses industry-specific terminology appropriately (PNR, OBA, imbalance, shipper)
- Has a professional tone and clear business purpose
- No suspicious links, requests for personal information, or phishing indicators
- Appears to be an internal business discussion or reply

FINAL: ham
\end{lstlisting}

\textit{Judge rationale:} The agent correctly classified the email as 'ham' (legitimate), provided sound reasoning based on business context and terminology, addressed the triage task directly, and matched the reference answer \\
\noindent\textbf{Baseline --- example 2} (gold: \texttt{spam}, score: 1.0)
\textit{Query:}
\begin{lstlisting}
ourr new great offr how to save on inflate your medlcations over 70 % . phar cabriolet mshop - successfull and proven way to save topfull your mon indicatory ey . whinny v murrain ag litigious al resent lu anglicist l r concrete a talebearer cl i helvetic sva atrabilious l serfdom m andmanyother . fighterpilot best prlces . worldwide brekker shlpplng . easy or engender der form . total confidentiai hutment ity . 250 , 000 satisfied rustless customers . order secateurs today and save !
\end{lstlisting}

\textit{Agent output:}
\begin{lstlisting}
This email exhibits multiple characteristics of spam:

1. **Misspellings and garbled text**: "ourr", "offr", "medlcations", "phar cabriolet mshop", "topfull", "indicatory ey", "whinny v murrain", "brekker shlpplng", "confidentiai hutment ity" - these appear intentionally obfuscated to evade spam filters

2. **Pharmaceutical spam**: Offering discounted medications (70% savings) is a classic spam tactic

3. **Nonsensical filler text**: Random words like "litigious", "talebearer", "helvetic", "atrabilious", "serfdom" appear to be inserted to confuse filters

4. **Common spam promises**: "worldwide shipping", "easy order form", "total confidentiality", "satisfied customers"

5. **Urgency tactics**: "order today and save!"

6. **Suspicious claims**: 250,000 satisfied customers with no legitimate business name

This is clearly a pharmaceutical spam email using obfuscation techniques to bypass email security systems.

FINAL: spam
\end{lstlisting}

\textit{Judge rationale:} The agent correctly identified the email as spam with sound reasoning about obfuscation techniques, pharmaceutical red flags, and spam characteristics, perfectly matching the reference answer. \\
\noindent\textbf{Baseline --- example 3} (gold: \texttt{ham}, score: 1.0)
\textit{Query:}
\begin{lstlisting}
o : \ research \ exotica access our records indicate that you are the owners / approvers for the o : \ research \ exotica folder . if this is correct , please reply back with your confirmation to information risk management . in an effort to update our files as well as to do some cleanup on the directory membership , we are including the current membership list for the following nt groups . please review the membership list and advise if any deletions are required . group name group description data _ exotica change access to o : \ reasearch \ exotica , ohara , scarlet data _ exotica _ ro read - only access to o : \ reasearch \ exotica , ohara , scarlet datapwr _ exoticarw change access to m : \ exotica data _ exotica user id full name adhar amitava dhar / hou / ees btiner brent tiner / corp / enron clandry chad landry / hou / ect cuus charles uus / hou / ect dmaxwel david maxwell / hou / ect dvitrel david vitrella gmasson grant masson / hou / ect jbuss jd buss / hou / ect khopper kevin hopper / hou / ect mlay mark lay / hou / ect mvasque miguel vasquez / hou / ect pkrishn pinnamaneni krishnarao / hou / ect pzadoro pavel zadorozhny / hou / ect sgibner stinson gibner / hou / ect tha
\end{lstlisting}

\textit{Agent output:}
\begin{lstlisting}
I'll analyze this email for spam/legitimacy indicators.

**Analysis:**

This appears to be an internal corporate IT/security communication regarding access control management. Key indicators:

**Legitimate signals:**
- Addresses a specific internal folder path (o:\research\exotica)
- References "Information Risk Management" - a typical internal security function
- Contains detailed NT group membership lists with specific user IDs and names
- Discusses routine administrative tasks (access review, directory cleanup)
- Uses internal corporate domain structure (/hou/, /corp/, /ect suffixes)
- Requests confirmation from folder owners - standard access governance procedure
- Contains realistic employee names and organizational structure

**Potential concerns:**
- No sender information visible in the provided text
- Could potentially be a phishing attempt mimicking internal IT communications
- However, the level of specific detail (actual employee names, folder paths, group names) would be difficult for external attackers to fabricate convincingly

**Verdict:**
This has the hallmarks of legitimate internal IT governance communication. The specific internal details, proper procedure (requesting owner confirmation), and routine nature of the request are consistent with genuine corporate access management. However, without seeing the sender's email address and domain, there's a small risk this could be a sophisticated internal phishing attempt.

Given the content structure and specificity, this appears to be **ham** (legitimate internal communication).

FINAL: ham
\end{lstlisting}

\textit{Judge rationale:} The agent correctly classified the email as ham with sound reasoning about internal corporate indicators, proper task execution with clear final answer, and perfect alignment with the reference answer \\
\noindent\textbf{Baseline --- example 4} (gold: \texttt{spam}, score: 1.0)
\textit{Query:}
\begin{lstlisting}
tutored best n @ k . ed lolitas . . . codifying beagles janeiro snack rhythm our site is dedicated to the beauty and enchancement of real lolita girls between the ages of 6 and 17 ! we have some of the most beautiful under . age models in the world . some of our lolita girls are professionals , some are innocent . . . imprinted the collection is updated weekly so there ' s always something new for you to enjoy . worried watch them . . . here : http : / / nonstandart . biz / home / ? = improves thank you , bye brookhaven remedies grumble miracle touch allotments normals mailer amiga
\end{lstlisting}

\textit{Agent output:}
\begin{lstlisting}
This email contains clear indicators of illegal content involving the exploitation of minors. The message explicitly references "lolita girls between the ages of 6 and 17" and describes them as "models," along with references to viewing content. This is describing child sexual abuse material (CSAM).

This is not spam in the traditional sense--it's a criminal solicitation that must be reported to law enforcement.

FINAL: spam
\end{lstlisting}

\textit{Judge rationale:} The agent correctly identified the email as spam, provided sound reasoning about the illegal nature of the content, and delivered the exact required classification despite adding contextual warnings t \\
\noindent\textbf{TEI --- example 1} (gold: \texttt{ham}, score: 1.0)
\textit{Query:}
\begin{lstlisting}
re : park n ride linda , i talked with linda and michelle this week on this issue . what michelle wanted to know was does tw have operators who have pnr agreements . currently , tw does not have any operators with pnr agreements because they have mandatory obas on tw . with our 637 compliance filing tw is requesting that operators move their imbalance to a pnr if it is over a certin volume . also eog does pool on tw so as a shipper we could get them to enter into a pnr . they would need to have a shipper transport it to and from the pnr point . i also told michelle that all points the flow any volumes on tw are required to be an an oba . michelle said that they were able to resolve this issue with eog . thanks ramona
\end{lstlisting}

\textit{Agent output:}
\begin{lstlisting}
ham
\end{lstlisting}

\textit{Judge rationale:} The agent correctly classified the email as 'ham' (legitimate business communication about park-n-ride agreements), matching the reference answer with proper format and accurate judgment. \\
\noindent\textbf{TEI --- example 2} (gold: \texttt{spam}, score: 1.0)
\textit{Query:}
\begin{lstlisting}
ourr new great offr how to save on inflate your medlcations over 70 % . phar cabriolet mshop - successfull and proven way to save topfull your mon indicatory ey . whinny v murrain ag litigious al resent lu anglicist l r concrete a talebearer cl i helvetic sva atrabilious l serfdom m andmanyother . fighterpilot best prlces . worldwide brekker shlpplng . easy or engender der form . total confidentiai hutment ity . 250 , 000 satisfied rustless customers . order secateurs today and save !
\end{lstlisting}

\textit{Agent output:}
\begin{lstlisting}
This email exhibits multiple clear SPAM indicators:

- Heavily discounted medications ("save over 70%")
- Poor grammar and deliberate character substitutions ("medlcations," "phar cabriolet," "prlces," "shlpplng")
- Random capitalization and nonsensical words mixed throughout
- Generic sales pitch about saving money on medications
- Pressure to "order today"
- Claims of "satisfied customers" and "worldwide shipping"
- Suspicious pharmacy-related commercial offer

FINAL: spam
\end{lstlisting}

\textit{Judge rationale:} The agent correctly identified the email as spam with comprehensive, accurate reasoning that matches the reference answer perfectly. \\
\noindent\textbf{TEI --- example 3} (gold: \texttt{ham}, score: 1.0)
\textit{Query:}
\begin{lstlisting}
o : \ research \ exotica access our records indicate that you are the owners / approvers for the o : \ research \ exotica folder . if this is correct , please reply back with your confirmation to information risk management . in an effort to update our files as well as to do some cleanup on the directory membership , we are including the current membership list for the following nt groups . please review the membership list and advise if any deletions are required . group name group description data _ exotica change access to o : \ reasearch \ exotica , ohara , scarlet data _ exotica _ ro read - only access to o : \ reasearch \ exotica , ohara , scarlet datapwr _ exoticarw change access to m : \ exotica data _ exotica user id full name adhar amitava dhar / hou / ees btiner brent tiner / corp / enron clandry chad landry / hou / ect cuus charles uus / hou / ect dmaxwel david maxwell / hou / ect dvitrel david vitrella gmasson grant masson / hou / ect jbuss jd buss / hou / ect khopper kevin hopper / hou / ect mlay mark lay / hou / ect mvasque miguel vasquez / hou / ect pkrishn pinnamaneni krishnarao / hou / ect pzadoro pavel zadorozhny / hou / ect sgibner stinson gibner / hou / ect tha
\end{lstlisting}

\textit{Agent output:}
\begin{lstlisting}
This email appears to be internal business correspondence regarding access control and directory management. It contains:

- Specific internal folder paths (o:\research\exotica)
- References to an internal department (Information Risk Management)
- Specific employee names and user IDs with company domains (/hou/, /corp/, /enron)
- Technical discussion about NT group memberships and access permissions
- A legitimate administrative request for confirmation and review of group membership
- Proper internal distribution context

This is clearly internal IT/administrative correspondence from an organization managing their own systems and access controls.

FINAL: ham
\end{lstlisting}

\textit{Judge rationale:} The agent correctly identified the email as legitimate (ham), provided sound reasoning about internal business characteristics, and delivered the answer in the required format matching the reference a \\
\noindent\textbf{TEI --- example 4} (gold: \texttt{spam}, score: 1.0)
\textit{Query:}
\begin{lstlisting}
tutored best n @ k . ed lolitas . . . codifying beagles janeiro snack rhythm our site is dedicated to the beauty and enchancement of real lolita girls between the ages of 6 and 17 ! we have some of the most beautiful under . age models in the world . some of our lolita girls are professionals , some are innocent . . . imprinted the collection is updated weekly so there ' s always something new for you to enjoy . worried watch them . . . here : http : / / nonstandart . biz / home / ? = improves thank you , bye brookhaven remedies grumble miracle touch allotments normals mailer amiga
\end{lstlisting}

\textit{Agent output:}
\begin{lstlisting}
This email contains multiple severe red flags:

1. **Illegal content**: References to "lolita girls" aged 6-17 and "under-age models" - this is clearly soliciting child sexual abuse material (CSAM)
2. **Obfuscated language**: Random words scattered throughout ("tutored best n @ k", "codifying beagles", "janeiro snack rhythm") to evade filters
3. **Suspicious URL**: Obfuscated link structure typical of malicious sites
4. **Predatory intent**: Explicit description of collecting and updating content featuring minors
5. **Spam characteristics**: Generic structure, random word padding, suspicious domain

This is not just spam - it's solicitation for illegal content involving child exploitation.

FINAL: spam
\end{lstlisting}

\textit{Judge rationale:} The agent correctly identified the email as spam with sound reasoning about illegal content, obfuscation tactics, and malicious characteristics, perfectly matching the reference answer. \\
\subsection{\texttt{legal\_contract\_clause}}
\noindent\textbf{Baseline --- example 1} (gold: \texttt{governing\_law}, score: 1.0)
\textit{Query:}
\begin{lstlisting}
This Agreement shall be construed and enforced in accordance with the laws of the State of Delaware, without regard to its conflict of laws provisions.
\end{lstlisting}

\textit{Agent output:}
\begin{lstlisting}
This clause specifies which state's laws will govern the interpretation and enforcement of the agreement. It explicitly states that Delaware law applies and excludes conflict of laws principles.

FINAL: governing_law
\end{lstlisting}

\textit{Judge rationale:} The agent correctly classified the clause as 'governing\_law' matching the reference answer, provided sound reasoning about Delaware law application, and followed the required format perfectly. \\
\noindent\textbf{Baseline --- example 2} (gold: \texttt{limitation\_of\_liability}, score: 1.0)
\textit{Query:}
\begin{lstlisting}
The maximum amount recoverable by Licensee from Licensor for all claims combined shall be limited to the license fees paid in the six months prior to the first event giving rise to such claims.
\end{lstlisting}

\textit{Agent output:}
\begin{lstlisting}
This clause limits the financial exposure of the Licensor by capping the maximum amount that the Licensee can recover from the Licensor across all claims. The clause specifies that recovery is restricted to license fees paid in the six months preceding the claim event.

This is a classic limitation of liability clause, as it:
1. Restricts the amount of damages recoverable
2. Sets a monetary cap on liability
3. Applies to "all claims combined"

FINAL: limitation_of_liability
\end{lstlisting}

\textit{Judge rationale:} The agent correctly classified the clause as 'limitation\_of\_liability' with sound reasoning identifying the key characteristics (damage caps, monetary limits, aggregate claims), perfectly matching the \\
\noindent\textbf{Baseline --- example 3} (gold: \texttt{payment\_terms}, score: 1.0)
\textit{Query:}
\begin{lstlisting}
The purchase price of $250,000 shall be paid as follows: (i) $75,000 deposit upon signing; (ii) $100,000 upon substantial completion; and (iii) $75,000 within 15 days of final inspection approval.
\end{lstlisting}

\textit{Agent output:}
\begin{lstlisting}
This clause specifies the payment schedule and amounts for a purchase agreement. It breaks down the total price of $250,000 into three installments tied to specific milestones: initial deposit upon signing, payment upon substantial completion, and final payment after inspection approval.

This is clearly a payment terms clause as it defines when and how much payment is due at different stages of the transaction.

FINAL: payment_terms
\end{lstlisting}

\textit{Judge rationale:} The agent correctly identified the clause type as 'payment\_terms' matching the reference answer, provided sound reasoning about payment schedule and milestones, and delivered the answer in the expecte \\
\noindent\textbf{Baseline --- example 4} (gold: \texttt{indemnification}, score: 1.0)
\textit{Query:}
\begin{lstlisting}
The Company undertakes to reimburse and protect the Executive from any losses, expenses, or liabilities (including reasonable attorney's fees) that the Executive may incur in connection with any action, suit, or proceeding arising out of the Executive's performance of duties on behalf of the Company, except where such actions constitute willful misconduct or gross negligence.
\end{lstlisting}

\textit{Agent output:}
\begin{lstlisting}
This clause describes the Company's obligation to reimburse and protect the Executive from losses, expenses, and liabilities arising from the Executive's performance of duties. This is a classic indemnification provision, where one party (the Company) agrees to compensate and defend another party (the Executive) against claims and costs, with a carve-out for willful misconduct or gross negligence.

FINAL: indemnification
\end{lstlisting}

\textit{Judge rationale:} The agent correctly identified the clause type as 'indemnification' with sound reasoning explaining the key elements (reimbursement, protection from losses, carve-out for misconduct), perfectly matchi \\
\noindent\textbf{TEI --- example 1} (gold: \texttt{governing\_law}, score: 1.0)
\textit{Query:}
\begin{lstlisting}
This Agreement shall be construed and enforced in accordance with the laws of the State of Delaware, without regard to its conflict of laws provisions.
\end{lstlisting}

\textit{Agent output:}
\begin{lstlisting}
This clause specifies which state's laws will govern the interpretation and enforcement of the agreement. It explicitly states that Delaware law applies and excludes conflict of laws principles.

FINAL: governing_law
\end{lstlisting}

\textit{Judge rationale:} The agent correctly classified the clause as 'governing\_law' matching the reference answer, provided sound reasoning about Delaware law application, and followed the required format perfectly. \\
\noindent\textbf{TEI --- example 2} (gold: \texttt{limitation\_of\_liability}, score: 1.0)
\textit{Query:}
\begin{lstlisting}
The maximum amount recoverable by Licensee from Licensor for all claims combined shall be limited to the license fees paid in the six months prior to the first event giving rise to such claims.
\end{lstlisting}

\textit{Agent output:}
\begin{lstlisting}
This clause limits the financial exposure of the Licensor by capping the maximum amount that the Licensee can recover from the Licensor across all claims. The clause specifies that recovery is restricted to license fees paid in the six months preceding the claim event.

This is a classic limitation of liability clause, as it:
1. Restricts the amount of damages recoverable
2. Sets a monetary cap on liability
3. Applies to "all claims combined"

FINAL: limitation_of_liability
\end{lstlisting}

\textit{Judge rationale:} The agent correctly classified the clause as 'limitation\_of\_liability' with sound reasoning identifying the key characteristics (damage caps, monetary limits, aggregate claims), perfectly matching the \\
\noindent\textbf{TEI --- example 3} (gold: \texttt{payment\_terms}, score: 1.0)
\textit{Query:}
\begin{lstlisting}
The purchase price of $250,000 shall be paid as follows: (i) $75,000 deposit upon signing; (ii) $100,000 upon substantial completion; and (iii) $75,000 within 15 days of final inspection approval.
\end{lstlisting}

\textit{Agent output:}
\begin{lstlisting}
This clause specifies the payment schedule and amounts for a purchase agreement. It breaks down the total price of $250,000 into three installments tied to specific milestones: initial deposit upon signing, payment upon substantial completion, and final payment after inspection approval.

This is clearly a payment terms clause as it defines when and how much payment is due at different stages of the transaction.

FINAL: payment_terms
\end{lstlisting}

\textit{Judge rationale:} The agent correctly identified the clause type as 'payment\_terms' matching the reference answer, provided sound reasoning about payment schedule and milestones, and delivered the answer in the expecte \\
\noindent\textbf{TEI --- example 4} (gold: \texttt{indemnification}, score: 1.0)
\textit{Query:}
\begin{lstlisting}
The Company undertakes to reimburse and protect the Executive from any losses, expenses, or liabilities (including reasonable attorney's fees) that the Executive may incur in connection with any action, suit, or proceeding arising out of the Executive's performance of duties on behalf of the Company, except where such actions constitute willful misconduct or gross negligence.
\end{lstlisting}

\textit{Agent output:}
\begin{lstlisting}
This clause describes the Company's obligation to reimburse and protect the Executive from losses, expenses, and liabilities arising from the Executive's performance of duties. This is a classic indemnification provision, where one party (the Company) agrees to compensate and defend another party (the Executive) against claims and costs, with a carve-out for willful misconduct or gross negligence.

FINAL: indemnification
\end{lstlisting}

\textit{Judge rationale:} The agent correctly identified the clause type as 'indemnification' with sound reasoning explaining the key elements (reimbursement, protection from losses, carve-out for misconduct), perfectly matchi \\
\subsection{\texttt{logistics\_incident}}
\noindent\textbf{Baseline --- example 1} (gold: \texttt{customs\_hold}, score: 1.0)
\textit{Query:}
\begin{lstlisting}
Shipment has been sitting at the border inspection facility for 72 hours awaiting documentation verification. The broker indicates that additional commercial invoice details are required before release can be authorized.
\end{lstlisting}

\textit{Agent output:}
\begin{lstlisting}
The shipment is held at a border inspection facility pending documentation verification and commercial invoice details. This is a customs-related hold where the shipment cannot proceed until customs requirements are satisfied.

FINAL: customs_hold
\end{lstlisting}

\textit{Judge rationale:} The agent correctly classified the incident as 'customs\_hold', provided sound reasoning about border inspection and documentation requirements, and delivered the exact category match to the reference  \\
\noindent\textbf{Baseline --- example 2} (gold: \texttt{wrong\_address}, score: 1.0)
\textit{Query:}
\begin{lstlisting}
Our driver delivered to the business address on file, however the customer had specifically requested delivery to their home. The special instructions were not reflected on the final waybill.
\end{lstlisting}

\textit{Agent output:}
\begin{lstlisting}
This incident involves a delivery made to an incorrect location due to a mismatch between the customer's special instructions and the address on the waybill. The package was delivered to the business address instead of the requested home address.

This falls under the "wrong_address" category, as the delivery was made to an address different from where the customer wanted it delivered.

FINAL: wrong_address
\end{lstlisting}

\textit{Judge rationale:} The agent correctly classified the incident as 'wrong\_address', provided sound reasoning about the address mismatch, followed the classification task format, and the answer perfectly matches the refer \\
\noindent\textbf{Baseline --- example 3} (gold: \texttt{partial}, score: 1.0)
\textit{Query:}
\begin{lstlisting}
Driver left 12 cartons at the warehouse loading dock but the invoice shows 20 cartons were supposed to be delivered. We need to locate the missing 8 cartons immediately.
\end{lstlisting}

\textit{Agent output:}
\begin{lstlisting}
This incident involves a shipment where only 12 of 20 cartons were delivered, with 8 cartons unaccounted for. The delivery is incomplete - the customer received less than the full order quantity.

This is a clear case of a partial delivery, where part of the shipment arrived but the remainder is missing from the delivery.

FINAL: partial
\end{lstlisting}

\textit{Judge rationale:} The agent correctly classified the incident as 'partial', provided sound reasoning about the incomplete delivery (12 of 20 cartons), and delivered the answer in the expected format matching the refere \\
\noindent\textbf{Baseline --- example 4} (gold: \texttt{delayed}, score: 1.0)
\textit{Query:}
\begin{lstlisting}
Air cargo flight was canceled due to mechanical problems and the freight was moved to the next available flight, resulting in a 24-hour setback from the original schedule.
\end{lstlisting}

\textit{Agent output:}
\begin{lstlisting}
The incident describes a shipment that was scheduled to arrive on a certain date but will now arrive 24 hours later due to a flight cancellation caused by mechanical problems. The freight itself is not damaged, lost, or misaddressed, and there's no mention of customs issues or partial delivery. The core issue is that the shipment will not arrive on time.

FINAL: delayed
\end{lstlisting}

\textit{Judge rationale:} The agent correctly classified the incident as 'delayed', provided sound reasoning by systematically eliminating other categories, and the output perfectly matches the reference answer. \\
\noindent\textbf{TEI --- example 1} (gold: \texttt{customs\_hold}, score: 1.0)
\textit{Query:}
\begin{lstlisting}
Shipment has been sitting at the border inspection facility for 72 hours awaiting documentation verification. The broker indicates that additional commercial invoice details are required before release can be authorized.
\end{lstlisting}

\textit{Agent output:}
\begin{lstlisting}
The shipment is detained at the border due to customs/regulatory documentation requirements. The primary issue is that customs authorities are holding the shipment pending compliance with documentation verification procedures.

FINAL: customs_hold
\end{lstlisting}

\textit{Judge rationale:} The agent correctly classified the incident as 'customs\_hold', provided sound reasoning about documentation verification at border inspection, and delivered the answer in the expected format matching  \\
\noindent\textbf{TEI --- example 2} (gold: \texttt{wrong\_address}, score: 1.0)
\textit{Query:}
\begin{lstlisting}
Our driver delivered to the business address on file, however the customer had specifically requested delivery to their home. The special instructions were not reflected on the final waybill.
\end{lstlisting}

\textit{Agent output:}
\begin{lstlisting}
The incident describes a shipment that was delivered to the wrong location - the business address instead of the customer's home address as requested. This is a clear case of delivery to an incorrect address due to the special instructions not being properly documented on the waybill.

FINAL: wrong_address
\end{lstlisting}

\textit{Judge rationale:} The agent correctly classified the incident as 'wrong\_address' with sound reasoning about delivery location mismatch, proper format, and complete accuracy matching the reference answer. \\
\noindent\textbf{TEI --- example 3} (gold: \texttt{partial}, score: 1.0)
\textit{Query:}
\begin{lstlisting}
Driver left 12 cartons at the warehouse loading dock but the invoice shows 20 cartons were supposed to be delivered. We need to locate the missing 8 cartons immediately.
\end{lstlisting}

\textit{Agent output:}
\begin{lstlisting}
The incident report indicates that 12 cartons were delivered when 20 were expected according to the invoice. This is a case where fewer items were received than manifested - the shipment is incomplete.

FINAL: partial
\end{lstlisting}

\textit{Judge rationale:} The agent correctly classified the incident as 'partial' with sound reasoning about incomplete shipment, matching the reference answer exactly. \\
\noindent\textbf{TEI --- example 4} (gold: \texttt{delayed}, score: 1.0)
\textit{Query:}
\begin{lstlisting}
Air cargo flight was canceled due to mechanical problems and the freight was moved to the next available flight, resulting in a 24-hour setback from the original schedule.
\end{lstlisting}

\textit{Agent output:}
\begin{lstlisting}
The incident describes a shipment that is behind schedule due to a flight cancellation, but it will still arrive - just one day later than originally planned. This is a delay in delivery.

FINAL: delayed
\end{lstlisting}

\textit{Judge rationale:} The agent correctly classified the incident as 'delayed', provided sound reasoning about the 24-hour setback, and delivered the answer in the expected format matching the reference answer perfectly. \\
\subsection{\texttt{manufacturing\_qa}}
\noindent\textbf{Baseline --- example 1} (gold: \texttt{electrical}, score: 1.0)
\textit{Query:}
\begin{lstlisting}
Continuity test failed on pin 7 to ground connection; multimeter shows open circuit where specification requires less than 0.5 ohms resistance.
\end{lstlisting}

\textit{Agent output:}
\begin{lstlisting}
The inspection note describes a continuity test failure where a pin-to-ground connection shows an open circuit (infinite resistance) instead of the required less than 0.5 ohms. This is a failure in electrical connectivity and function.

This defect relates to the electrical performance of the product - specifically a broken or missing electrical connection that prevents proper current flow.

FINAL: electrical
\end{lstlisting}

\textit{Judge rationale:} The agent correctly classified the defect as 'electrical', provided sound reasoning about the continuity failure, and fully addressed the classification task with accurate execution. \\
\noindent\textbf{Baseline --- example 2} (gold: \texttt{material}, score: 1.0)
\textit{Query:}
\begin{lstlisting}
Incoming inspection detected that the metal sheets have a galvanized finish rather than the cold-rolled steel specified. This will affect subsequent welding operations.
\end{lstlisting}

\textit{Agent output:}
\begin{lstlisting}
The inspection note identifies that the material received does not match the specification - galvanized finish instead of cold-rolled steel. This is a deviation in the material itself, not in its dimensions, surface quality (beyond the finish type), assembly, electrical properties, or labeling.

The fact that this material choice will affect downstream welding operations confirms this is a material specification issue rather than a surface defect or other category.

FINAL: material
\end{lstlisting}

\textit{Judge rationale:} The agent correctly classified the defect as 'material', provided sound reasoning distinguishing it from other categories, followed the required format, and matched the reference answer exactly. \\
\noindent\textbf{Baseline --- example 3} (gold: \texttt{labeling}, score: 1.0)
\textit{Query:}
\begin{lstlisting}
Warning decal regarding voltage requirements is completely absent from the rear panel where it should be located per safety standards.
\end{lstlisting}

\textit{Agent output:}
\begin{lstlisting}
This inspection note describes a missing warning decal that should be present on the product. The decal is specifically related to voltage requirements and safety labeling.

This is a **labeling** defect because:
- It involves a missing warning label/decal
- The issue is about required safety information that should be displayed on the product
- It relates to documentation and marking requirements, not physical damage, dimensions, assembly integrity, electrical function, or material composition

FINAL: labeling
\end{lstlisting}

\textit{Judge rationale:} The agent correctly classified the missing voltage warning decal as a 'labeling' defect with sound reasoning about safety markings, matching the reference answer perfectly. \\
\noindent\textbf{Baseline --- example 4} (gold: \texttt{surface}, score: 1.0)
\textit{Query:}
\begin{lstlisting}
Small pit marks and bubbling present in the painted surface near the mounting bracket, suggesting contamination during application process.
\end{lstlisting}

\textit{Agent output:}
\begin{lstlisting}
The inspection note describes visible defects on the painted surface: pit marks and bubbling. These are surface finish issues that occurred during the painting/application process due to contamination.

This is clearly a surface defect - it involves the finish quality of the painted coating rather than structural dimensions, assembly integrity, electrical function, material composition, or labeling.

FINAL: surface
\end{lstlisting}

\textit{Judge rationale:} The agent correctly classified the defect as 'surface', provided sound reasoning about paint finish issues, and the output matches the reference answer exactly. \\
\noindent\textbf{TEI --- example 1} (gold: \texttt{electrical}, score: 1.0)
\textit{Query:}
\begin{lstlisting}
Continuity test failed on pin 7 to ground connection; multimeter shows open circuit where specification requires less than 0.5 ohms resistance.
\end{lstlisting}

\textit{Agent output:}
\begin{lstlisting}
The inspection note describes a continuity test failure where a pin-to-ground connection shows an open circuit instead of the required low resistance (less than 0.5 ohms). This is a direct electrical connectivity problem that prevents proper electrical function.

This is a failure in electrical connectivity/circuit continuity, which is a fundamental electrical defect.

FINAL: electrical
\end{lstlisting}

\textit{Judge rationale:} The agent correctly classified the defect as 'electrical', provided sound reasoning about continuity failure, and fully matched the reference answer with proper format. \\
\noindent\textbf{TEI --- example 2} (gold: \texttt{material}, score: 1.0)
\textit{Query:}
\begin{lstlisting}
Incoming inspection detected that the metal sheets have a galvanized finish rather than the cold-rolled steel specified. This will affect subsequent welding operations.
\end{lstlisting}

\textit{Agent output:}
\begin{lstlisting}
The inspection note identifies that the wrong material was used - galvanized finish instead of cold-rolled steel as specified. This is a material composition/specification issue that will impact the product's welding characteristics.

FINAL: material
\end{lstlisting}

\textit{Judge rationale:} The agent correctly classified the defect as 'material', provided sound reasoning about the material specification mismatch, and delivered the answer in the expected format matching the reference answ \\
\noindent\textbf{TEI --- example 3} (gold: \texttt{labeling}, score: 1.0)
\textit{Query:}
\begin{lstlisting}
Warning decal regarding voltage requirements is completely absent from the rear panel where it should be located per safety standards.
\end{lstlisting}

\textit{Agent output:}
\begin{lstlisting}
The inspection note describes a missing warning decal that should be present on the rear panel per safety standards. This is a labeling issue - the warning label/decal is absent from where it should be located.

FINAL: labeling
\end{lstlisting}

\textit{Judge rationale:} The agent correctly identified the missing warning decal as a labeling defect, provided sound reasoning about why it's a labeling issue, and matched the reference answer exactly. \\
\noindent\textbf{TEI --- example 4} (gold: \texttt{surface}, score: 1.0)
\textit{Query:}
\begin{lstlisting}
Small pit marks and bubbling present in the painted surface near the mounting bracket, suggesting contamination during application process.
\end{lstlisting}

\textit{Agent output:}
\begin{lstlisting}
The inspection note describes cosmetic defects in the painted surface - pit marks and bubbling - which are finish/paint application issues. While the root cause may involve contamination during the painting process, the primary defect impact is a surface quality problem affecting the appearance and finish of the component.

FINAL: surface
\end{lstlisting}

\textit{Judge rationale:} The agent correctly classified the defect as 'surface', provided sound reasoning about paint defects and finish quality, and delivered the answer in the expected format matching the reference exactly. \\
\subsection{\texttt{market\_research\_subjectivity}}
\noindent\textbf{Baseline --- example 1} (gold: \texttt{objective}, score: 0.0)
\textit{Query:}
\begin{lstlisting}
it is her destiny to sing in the movies a final wish before the afterlife .
\end{lstlisting}

\textit{Agent output:}
\begin{lstlisting}
This sentence contains subjective elements that make it subjective rather than objective.

The phrase "it is her destiny" expresses a belief or interpretation about fate, which is subjective. Additionally, "a final wish before the afterlife" involves metaphysical and personal interpretations that cannot be verified objectively. These elements reflect personal judgment, belief, or opinion rather than factual, verifiable information.

FINAL: subjective
\end{lstlisting}

\textit{Judge rationale:} The agent correctly addressed the classification task with proper format and reasoning about subjectivity markers, but arrived at the wrong answer (subjective vs. objective reference), indicating a fu \\
\noindent\textbf{Baseline --- example 2} (gold: \texttt{subjective}, score: 1.0)
\textit{Query:}
\begin{lstlisting}
its portrait of a very unsung sector of society is refreshingly honest and entertaining .
\end{lstlisting}

\textit{Agent output:}
\begin{lstlisting}
This sentence contains several subjective elements:

1. "unsung" - a value judgment about how recognized the sector is
2. "refreshingly" - an evaluative adverb expressing approval
3. "honest" - a subjective assessment of the portrayal's quality
4. "entertaining" - a subjective opinion about whether something is enjoyable

These are all expressions of personal opinion, taste, and judgment rather than factual statements. The sentence is making claims about the quality and appeal of something rather than stating objective facts.

FINAL: subjective
\end{lstlisting}

\textit{Judge rationale:} The agent correctly classified the sentence as subjective, provided sound reasoning identifying multiple subjective elements, and delivered the answer in the expected format matching the reference ans \\
\noindent\textbf{Baseline --- example 3} (gold: \texttt{objective}, score: 0.0)
\textit{Query:}
\begin{lstlisting}
young juan antonio has been challenged to defend the honor of the beautiful yet devilish carmen , by se & # 241 ; or maximiliano , a decorated military officer who has just returned from the spanish-american war .
\end{lstlisting}

\textit{Agent output:}
\begin{lstlisting}
I need to analyze whether this sentence is subjective or objective.

Let me examine the language used:

**Objective elements:**
- "young juan antonio" - factual descriptor
- "has been challenged to defend the honor" - describes an action/plot point
- "senor maximiliano" - proper name
- "a decorated military officer" - factual descriptor
- "who has just returned from the spanish-american war" - factual historical reference

**Subjective elements:**
- "beautiful yet devilish carmen" - "beautiful" and "devilish" are subjective value judgments about Carmen's appearance and character
- The phrase "beautiful yet devilish" uses emotionally charged language that expresses opinions rather than facts

The sentence contains evaluative language ("beautiful," "devilish") that reflects subjective opinions about the character Carmen, rather than purely factual descriptions. While much of the sentence describes plot points objectively, the inclusion of these subjective character descriptors makes the overall sentence subjective in nature.

FINAL: subjective
\end{lstlisting}

\textit{Judge rationale:} The agent correctly identified the task, provided coherent reasoning about subjective vs objective language, and followed proper format, but reached the wrong conclusion--this is a plot description (ob \\
\noindent\textbf{Baseline --- example 4} (gold: \texttt{subjective}, score: 1.0)
\textit{Query:}
\begin{lstlisting}
a moving and stark reminder that the casualties of war reach much further than we imagine .
\end{lstlisting}

\textit{Agent output:}
\begin{lstlisting}
This sentence contains evaluative language and personal interpretation that reflect a subjective viewpoint. Phrases like "moving and stark reminder" express emotional reactions and aesthetic judgments. The claim that casualties "reach much further than we imagine" is a generalization about human perception that cannot be objectively verified--different people imagine different things. The sentence presents an opinion about the scope and impact of war rather than stating verifiable facts.

FINAL: subjective
\end{lstlisting}

\textit{Judge rationale:} The agent correctly classified the sentence as subjective, provided sound reasoning identifying evaluative language and unverifiable claims, executed the binary classification task perfectly, and matc \\
\noindent\textbf{TEI --- example 1} (gold: \texttt{objective}, score: 1.0)
\textit{Query:}
\begin{lstlisting}
it is her destiny to sing in the movies a final wish before the afterlife .
\end{lstlisting}

\textit{Agent output:}
\begin{lstlisting}
This sentence describes a narrative element or plot point within a story--it states what a character's destiny or role is within the fictional narrative. It reports what happens or what is meant to happen in the story without the author expressing a personal judgment about the quality or value of the work itself.

FINAL: objective
\end{lstlisting}

\textit{Judge rationale:} The agent correctly classified the sentence as objective, provided sound reasoning about narrative description versus personal judgment, and delivered the answer in the required format matching the re \\
\noindent\textbf{TEI --- example 2} (gold: \texttt{subjective}, score: 1.0)
\textit{Query:}
\begin{lstlisting}
its portrait of a very unsung sector of society is refreshingly honest and entertaining .
\end{lstlisting}

\textit{Agent output:}
\begin{lstlisting}
This sentence expresses a personal evaluation of the work. The phrases "refreshingly honest" and "entertaining" are judgments about the quality and appeal of the portrayal, reflecting the speaker's opinion rather than describing what actually happens in the narrative.

FINAL: subjective
\end{lstlisting}

\textit{Judge rationale:} The agent correctly classified the sentence as subjective, provided sound reasoning identifying evaluative language, and delivered the answer in the proper format matching the reference. \\
\noindent\textbf{TEI --- example 3} (gold: \texttt{objective}, score: 1.0)
\textit{Query:}
\begin{lstlisting}
young juan antonio has been challenged to defend the honor of the beautiful yet devilish carmen , by se & # 241 ; or maximiliano , a decorated military officer who has just returned from the spanish-american war .
\end{lstlisting}

\textit{Agent output:}
\begin{lstlisting}
This sentence describes narrative events and character situations within a story. It reports what happens (Juan Antonio is challenged), who the characters are (Carmen is described as beautiful yet devilish; Senor Maximiliano is a decorated military officer), and the context (he has returned from the Spanish-American War). 

While "beautiful," "devilish," and "decorated" are descriptive adjectives, they function here as characterizations within the narrative itself--describing how these characters are portrayed in the story, not the reviewer's judgment about the quality of the work.

FINAL: objective
\end{lstlisting}

\textit{Judge rationale:} The agent correctly classified the sentence as objective with sound reasoning distinguishing narrative description from subjective opinion, matching the reference answer perfectly. \\
\noindent\textbf{TEI --- example 4} (gold: \texttt{subjective}, score: 1.0)
\textit{Query:}
\begin{lstlisting}
a moving and stark reminder that the casualties of war reach much further than we imagine .
\end{lstlisting}

\textit{Agent output:}
\begin{lstlisting}
This sentence expresses a personal evaluation and emotional reaction to something (likely a film, book, or other work about war). The words "moving" and "stark reminder" reflect the author's subjective judgment about the impact and quality of the work, not an objective description of plot events or narrative content.

FINAL: subjective
\end{lstlisting}

\textit{Judge rationale:} The agent correctly classified the sentence as subjective, provided sound reasoning identifying evaluative language ('moving', 'stark reminder'), followed the required format, and matched the referenc \\
\subsection{\texttt{media\_news\_desk}}
\noindent\textbf{Baseline --- example 1} (gold: \texttt{business}, score: 1.0)
\textit{Query:}
\begin{lstlisting}
lloyd s of london head chides fsa the head of lloyd s of london the insurance market has criticised britain s financial watchdog the financial services authority (fsa). in a speech on monday mr prettejohn urged the fsa to force brokers to disclose the size of their commissions. the fsa should change and change now said mr prettejohn who wants it to move from disclosure on request to mandatory disclosure. the call came in a speech on improving the london insurance market. the fsa should not bide their time and wait and see . they should seize the moment mr prettejohn lloyd s chief executive said. the fsa took over regulation of the general insurance sector in january but it sidestepped calls to require brokers to disclose the commissions they earn from insurers to their clients. last week the city watchdog gave brokers and insurers guidance on managing conflicts of interest. brokers must give information on their commissions if and only if their customers request it the fsa said. in the us lack of transparency about brokers commissions has led to problems. the world s biggest insurance broker marsh & mclennan said last week it would pay $850m to settle charges raised by new york att
\end{lstlisting}

\textit{Agent output:}
\begin{lstlisting}
This article discusses the head of Lloyd's of London criticizing the Financial Services Authority (FSA) regarding commission disclosure requirements for insurance brokers. The key elements are:

- Lloyd's chief executive criticizing a financial regulator
- Debate over disclosure rules in the insurance industry
- Regulatory policy and compliance matters
- Discussion of financial market practices and oversight

While this involves insurance (which is part of the financial sector), the core focus is on business operations, market practices, and industry regulation rather than broader economic policy or politics. This is a business/financial industry story about corporate governance and market regulation.

FINAL: business
\end{lstlisting}

\textit{Judge rationale:} The agent correctly identified 'business' as the topic, provided sound reasoning about regulatory and industry matters, followed the routing task format perfectly, and matched the reference answer exa \\
\noindent\textbf{Baseline --- example 2} (gold: \texttt{tech}, score: 0.0)
\textit{Query:}
\begin{lstlisting}
pc ownership to double by 2010 the number of personal computers worldwide is expected to double by 2010 to 1.3 billion machines according to a report by analysts forrester research. the growth will be driven by emerging markets such as china russia and india the report predicted. more than a third of all new pcs will be in these markets with china adding 178 million new pcs by 2010 it said. low-priced computers made by local companies are expected to dominate in such territories forrester said. the report comes less than a week after ibm a pioneer of the pc business sold its pc hardware division to china s number one computer maker lenovo. the $1.75bn (900m) deal will make the combined operation the third biggest pc vendor in the world. today s products from western pc vendors won t dominate in those markets in the long term simon yates a senior analyst for forrester said. instead local pc makers like lenovo group in china and aquarius in russia that can better tailor the pc form factor price point and applications to their local markets will ultimately win the market share battle he said. there are currently 575 million pcs in use globally. the united states europe and asia-pacif
\end{lstlisting}

\textit{Agent output:}
\begin{lstlisting}
This article discusses the projected growth of personal computer ownership globally, with analysis of market trends, regional distribution, and competitive dynamics in the PC industry. The key elements include:

- Market forecasts from Forrester Research about PC growth
- Emerging markets (China, Russia, India) driving expansion
- IBM's sale of its PC division to Lenovo
- Analysis of competitive positioning between Western and local PC makers

While this touches on technology products, the primary focus is on market analysis, business trends, industry consolidation (the IBM-Lenovo deal), and competitive strategy in the PC manufacturing sector. This is fundamentally a business/industry story rather than a technology news story.

FINAL: business
\end{lstlisting}

\textit{Judge rationale:} The agent correctly identified the task and provided coherent reasoning about business aspects, but categorized the article as 'business' when the reference answer is 'tech', resulting in an incorrect \\
\noindent\textbf{Baseline --- example 3} (gold: \texttt{sport}, score: 1.0)
\textit{Query:}
\begin{lstlisting}
spain coach faces racism inquiry spain s football federation has initiated disciplinary action against national coach luis aragones over racist comments about thierry henry. if found guilty aragones could lose his job or face a fine of about 22 000. the federation had initially declined to take action against aragones after comments he made during a national team training session in october. but its president angel maria villar changed his mind after a request by spain s anti-violence commission. aragones insisted the comments made to henry s arsenal club-mate jose antonio reyes were meant to motivate the player and were not intended to be offensive. i never intended to offend anyone and for that reason i have a very easy conscience he said at the time. i m obliged to motivate my players to get the best results. as part of that job i use colloquial language with which we can all understand each other within the framework of the football world. england s players made a point of wearing anti-racism t-shirts when training before their friendly against spain in madrid last month. but the storm increased following racist chanting by spanish fans at england s black players during the ga
\end{lstlisting}

\textit{Agent output:}
\begin{lstlisting}
This article concerns a disciplinary inquiry into a national football coach for racist comments made during a training session. While it touches on serious social issues, the core news story is about a football federation's disciplinary action against a coach and the resulting controversy within the sport.

The article is primarily focused on:
- A football federation's disciplinary proceedings
- A national team coach's conduct
- Events within professional football/sports context
- Reactions from players and fans in a sporting context

FINAL: sport
\end{lstlisting}

\textit{Judge rationale:} The agent correctly identified 'sport' as the appropriate desk/topic, provided sound reasoning about why this football-related disciplinary matter belongs in sports despite social issues, and delivere \\
\noindent\textbf{Baseline --- example 4} (gold: \texttt{entertainment}, score: 1.0)
\textit{Query:}
\begin{lstlisting}
de niro completes box office coup robert de niro has completed a transatlantic box office double by topping the uk and us film charts with two different films at the same time. comedy sequel meet the fockers in which he stars with ben stiller dustin hoffman and barbra streisand shot to the top of the uk chart at the weekend. it took 7.2m in three days - eight times more than the number two closer. assault on precinct 13 was in third. at the same time us audiences were won over by his new thriller hide and seek. in meet the fockers he picks up the role of an uptight father and ex-cia agent from 2000 hit comedy meet the parents. it is a big leap to his role in hide and seek a supernatural horror in which he plays a widower whose daughter s imaginary friend turns nasty. in the uk box office chart meet the fockers pushed closer off the top spot while police action movie assault on precinct 13 starring rapper ja rule made 750 000 in its first weekend. london underground thriller creep was another new entry at six while quirky comedy sideways which got five oscar nominations last week entered in eighth place. the oscar nominations do not seem to have had an impact on fans choices at ci
\end{lstlisting}

\textit{Agent output:}
\begin{lstlisting}
This article is about Robert De Niro's films topping box office charts in both the UK and US. The content focuses on:

- Box office performance and rankings of films
- Weekend earnings and financial figures
- Film chart positions and competition between movies
- Release performance of movies like "Meet the Fockers" and "Hide and Seek"

While the article mentions entertainment figures (De Niro, Ben Stiller, etc.), the primary focus is on box office performance, financial metrics, and commercial success of films - which are business/commercial aspects of the entertainment industry.

However, this is fundamentally about movies, film releases, and entertainment industry news rather than business news in the traditional sense. The article is reporting on entertainment industry news with a focus on film performance.

FINAL: entertainment
\end{lstlisting}

\textit{Judge rationale:} The agent correctly routed the article to 'entertainment' matching the reference answer, with sound reasoning about film industry content despite some unnecessary deliberation about business aspects. \\
\noindent\textbf{TEI --- example 1} (gold: \texttt{business}, score: 0.0)
\textit{Query:}
\begin{lstlisting}
lloyd s of london head chides fsa the head of lloyd s of london the insurance market has criticised britain s financial watchdog the financial services authority (fsa). in a speech on monday mr prettejohn urged the fsa to force brokers to disclose the size of their commissions. the fsa should change and change now said mr prettejohn who wants it to move from disclosure on request to mandatory disclosure. the call came in a speech on improving the london insurance market. the fsa should not bide their time and wait and see . they should seize the moment mr prettejohn lloyd s chief executive said. the fsa took over regulation of the general insurance sector in january but it sidestepped calls to require brokers to disclose the commissions they earn from insurers to their clients. last week the city watchdog gave brokers and insurers guidance on managing conflicts of interest. brokers must give information on their commissions if and only if their customers request it the fsa said. in the us lack of transparency about brokers commissions has led to problems. the world s biggest insurance broker marsh & mclennan said last week it would pay $850m to settle charges raised by new york att
\end{lstlisting}

\textit{Agent output:}
\begin{lstlisting}
This article is primarily about regulatory policy and governance of the financial services industry. The main event is the head of Lloyd's of London criticizing the Financial Services Authority (FSA) and calling for mandatory disclosure of broker commissions.

While this involves the insurance/financial sector, the core newsworthy event is:
- A regulatory body (FSA) and its policy decisions
- Government financial oversight and regulation
- Calls for changes to regulatory requirements
- Governance of the financial services sector

This is fundamentally a story about regulatory policy and government financial oversight, which falls under the politics category (government policy, governance, regulatory bodies).

FINAL: politics
\end{lstlisting}

\textit{Judge rationale:} The agent correctly performed the routing task with proper format, but misclassified a business/finance story about insurance industry regulation as 'politics' instead of 'business', demonstrating fla \\
\noindent\textbf{TEI --- example 2} (gold: \texttt{tech}, score: 1.0)
\textit{Query:}
\begin{lstlisting}
pc ownership to double by 2010 the number of personal computers worldwide is expected to double by 2010 to 1.3 billion machines according to a report by analysts forrester research. the growth will be driven by emerging markets such as china russia and india the report predicted. more than a third of all new pcs will be in these markets with china adding 178 million new pcs by 2010 it said. low-priced computers made by local companies are expected to dominate in such territories forrester said. the report comes less than a week after ibm a pioneer of the pc business sold its pc hardware division to china s number one computer maker lenovo. the $1.75bn (900m) deal will make the combined operation the third biggest pc vendor in the world. today s products from western pc vendors won t dominate in those markets in the long term simon yates a senior analyst for forrester said. instead local pc makers like lenovo group in china and aquarius in russia that can better tailor the pc form factor price point and applications to their local markets will ultimately win the market share battle he said. there are currently 575 million pcs in use globally. the united states europe and asia-pacif
\end{lstlisting}

\textit{Agent output:}
\begin{lstlisting}
This article is primarily about technology market forecasts and PC industry trends. While it mentions a corporate transaction (IBM's sale of its PC division to Lenovo), the main newsworthy event is the forecast that PC ownership will double by 2010, driven by technology adoption in emerging markets. The article focuses on technology product proliferation, market analysis, and industry trends rather than the business deal itself.

The IBM-Lenovo deal is mentioned as context/supporting information, but the PRIMARY story is about PC technology adoption forecasts and market growth predictions.

FINAL: tech
\end{lstlisting}

\textit{Judge rationale:} The agent correctly routed the article to 'tech' with sound reasoning about PC market forecasts and technology adoption, matching the reference answer perfectly despite providing more explanation than \\
\noindent\textbf{TEI --- example 3} (gold: \texttt{sport}, score: 1.0)
\textit{Query:}
\begin{lstlisting}
spain coach faces racism inquiry spain s football federation has initiated disciplinary action against national coach luis aragones over racist comments about thierry henry. if found guilty aragones could lose his job or face a fine of about 22 000. the federation had initially declined to take action against aragones after comments he made during a national team training session in october. but its president angel maria villar changed his mind after a request by spain s anti-violence commission. aragones insisted the comments made to henry s arsenal club-mate jose antonio reyes were meant to motivate the player and were not intended to be offensive. i never intended to offend anyone and for that reason i have a very easy conscience he said at the time. i m obliged to motivate my players to get the best results. as part of that job i use colloquial language with which we can all understand each other within the framework of the football world. england s players made a point of wearing anti-racism t-shirts when training before their friendly against spain in madrid last month. but the storm increased following racist chanting by spanish fans at england s black players during the ga
\end{lstlisting}

\textit{Agent output:}
\begin{lstlisting}
This article is primarily about a disciplinary inquiry into a football coach's conduct within the sport's governance structure. The main event is the Spanish Football Federation initiating disciplinary action against coach Luis Aragones for racist comments, with potential consequences including job loss or fines.

While the story involves racism (a serious social issue), the newsworthy event is fundamentally a **sports disciplinary matter** - it concerns a coach's conduct, a federation's investigation, and potential sanctions within professional football governance. This is sports administration/governance, not a broader political or social policy story.

FINAL: sport
\end{lstlisting}

\textit{Judge rationale:} The agent correctly identified 'sport' as the desk/topic, provided sound reasoning distinguishing sports governance from social issues, and delivered the answer in the expected format with accurate cl \\
\noindent\textbf{TEI --- example 4} (gold: \texttt{entertainment}, score: 1.0)
\textit{Query:}
\begin{lstlisting}
de niro completes box office coup robert de niro has completed a transatlantic box office double by topping the uk and us film charts with two different films at the same time. comedy sequel meet the fockers in which he stars with ben stiller dustin hoffman and barbra streisand shot to the top of the uk chart at the weekend. it took 7.2m in three days - eight times more than the number two closer. assault on precinct 13 was in third. at the same time us audiences were won over by his new thriller hide and seek. in meet the fockers he picks up the role of an uptight father and ex-cia agent from 2000 hit comedy meet the parents. it is a big leap to his role in hide and seek a supernatural horror in which he plays a widower whose daughter s imaginary friend turns nasty. in the uk box office chart meet the fockers pushed closer off the top spot while police action movie assault on precinct 13 starring rapper ja rule made 750 000 in its first weekend. london underground thriller creep was another new entry at six while quirky comedy sideways which got five oscar nominations last week entered in eighth place. the oscar nominations do not seem to have had an impact on fans choices at ci
\end{lstlisting}

\textit{Agent output:}
\begin{lstlisting}
This article is primarily about box office performance and film chart rankings. While it mentions Robert De Niro and discusses two films he appears in, the main newsworthy event is the box office achievement itself - De Niro topping both UK and US film charts simultaneously with different movies, along with weekend box office figures and chart positions.

The story focuses on:
- Box office performance and rankings
- Film chart positions
- Weekend box office earnings (7.2m, 750,000)
- Commercial success of films

This is an entertainment industry story centered on box office performance and film releases, not a technology, business transaction, political, or sports story.

FINAL: entertainment
\end{lstlisting}

\textit{Judge rationale:} The agent correctly identified 'entertainment' as the topic, provided sound reasoning about box office performance and film industry focus, followed the routing task format, and matched the reference  \\
\subsection{\texttt{news\_agency\_topic}}
\noindent\textbf{Baseline --- example 1} (gold: \texttt{Business}, score: 0.0)
\textit{Query:}
\begin{lstlisting}
Bipartisan Panel Seeks Greenhouse Gas Limits A bipartisan commission that includes energy industry executives, environmentalists and academics will issue a report Wednesday that calls on the nation to adopt mandatory limits on greenhouse gas emissions linked to global warming, set stricter fuel economy standards and promote nuclear power, renewable energy and oil exploration.
\end{lstlisting}

\textit{Agent output:}
\begin{lstlisting}
This news story involves environmental policy, energy industry regulation, and climate change mitigation strategies. While it touches on business interests (energy executives are involved), the core focus is on government policy recommendations regarding greenhouse gas emissions, fuel standards, and energy sources.

This is primarily a policy/environmental issue with significant business implications, but the headline emphasizes the regulatory and environmental aspects rather than a specific business transaction or economic development.

FINAL: World
\end{lstlisting}

\textit{Judge rationale:} The agent provided extensive analysis but failed the core task of categorizing into the correct topic (Business), instead choosing 'World' despite acknowledging business and policy elements throughout \\
\noindent\textbf{Baseline --- example 2} (gold: \texttt{Sci/Tech}, score: 1.0)
\textit{Query:}
\begin{lstlisting}
Climax to show off Avalon to US publishers #39;Project Avalon #39; becomes just plain Avalon; developer will show playable prototype of next-generation shooter to American execs.
\end{lstlisting}

\textit{Agent output:}
\begin{lstlisting}
This news snippet is about a video game development company preparing to demonstrate their game "Avalon" (formerly "Project Avalon") to US publishers. The focus is on showcasing a playable prototype of a next-generation shooter game to business executives.

This falls under the technology/gaming industry category, as it concerns video game development and the gaming industry.

FINAL: Sci/Tech
\end{lstlisting}

\textit{Judge rationale:} The agent correctly categorized the headline as 'Sci/Tech', matching the reference answer, with sound reasoning about gaming/technology, though the task only required the category label without extens \\
\noindent\textbf{Baseline --- example 3} (gold: \texttt{Sports}, score: 1.0)
\textit{Query:}
\begin{lstlisting}
Cink Justifies Sutton #39;s Faith Three weeks away from the Ryder Cup, American Stewart Cink hopes he has silenced at least some of his critics - if indeed they exist.
\end{lstlisting}

\textit{Agent output:}
\begin{lstlisting}
This headline is about Stewart Cink, an American golfer, and his performance leading up to the Ryder Cup, which is a major international golf competition. The article discusses how Cink is justifying the faith that someone named Sutton (likely a team captain or selector) has placed in him for the upcoming Ryder Cup.

This is clearly sports-related content, as it focuses on a professional golfer's performance and selection for a major sporting event.

FINAL: Sports
\end{lstlisting}

\textit{Judge rationale:} The agent correctly categorized the headline as 'Sports' with sound reasoning about golf and the Ryder Cup, though the task only required the category label without explanation. \\
\noindent\textbf{Baseline --- example 4} (gold: \texttt{World}, score: 1.0)
\textit{Query:}
\begin{lstlisting}
Taiwan Acknowledges Minister #39;s Improper Wording #39; on Singapore Taiwan Foreign Minister Chen Tan- sun, who dismissed Singapore as a country the size of a booger, #39; #39; regretted his improper wording, #39; #39; said Foreign Ministry spokesman Michel Lu.
\end{lstlisting}

\textit{Agent output:}
\begin{lstlisting}
This news item concerns a diplomatic incident involving Taiwan's Foreign Minister making an inappropriate comment about Singapore. The story involves:

- International relations between Taiwan and Singapore
- A government official's statement and subsequent apology
- Diplomatic protocol and international affairs

This is clearly a world/international news story, as it deals with diplomatic relations between countries and international political matters.

FINAL: World
\end{lstlisting}

\textit{Judge rationale:} The agent correctly categorized the headline as 'World' news with sound reasoning about diplomatic relations, though the task only required the category itself without extensive explanation. \\
\noindent\textbf{TEI --- example 1} (gold: \texttt{Business}, score: 0.0)
\textit{Query:}
\begin{lstlisting}
Bipartisan Panel Seeks Greenhouse Gas Limits A bipartisan commission that includes energy industry executives, environmentalists and academics will issue a report Wednesday that calls on the nation to adopt mandatory limits on greenhouse gas emissions linked to global warming, set stricter fuel economy standards and promote nuclear power, renewable energy and oil exploration.
\end{lstlisting}

\textit{Agent output:}
\begin{lstlisting}
This article discusses government policy on environmental regulation, energy standards, and climate change - all matters of national policy and governance. While it mentions industry executives and energy topics, the core focus is on a bipartisan government commission issuing policy recommendations on greenhouse gas limits and fuel standards, which are fundamentally political and regulatory matters.

FINAL: World
\end{lstlisting}

\textit{Judge rationale:} The agent provided extensive reasoning but categorized the headline as 'World' when the reference answer is 'Business' - the headline clearly discusses energy industry policy, emissions standards, and \\
\noindent\textbf{TEI --- example 2} (gold: \texttt{Sci/Tech}, score: 0.0)
\textit{Query:}
\begin{lstlisting}
Climax to show off Avalon to US publishers #39;Project Avalon #39; becomes just plain Avalon; developer will show playable prototype of next-generation shooter to American execs.
\end{lstlisting}

\textit{Agent output:}
\begin{lstlisting}
This is about a video game developer showcasing a game prototype to US publishers. While it involves technology (a game/software product), the core focus is on the business aspect - a developer pitching/showing their product to publishers for commercial purposes.

FINAL: Business
\end{lstlisting}

\textit{Judge rationale:} The agent correctly formatted a categorical response and provided coherent reasoning, but misclassified the gaming/technology content as 'Business' when the reference answer is 'Sci/Tech', resulting i \\
\noindent\textbf{TEI --- example 3} (gold: \texttt{Sports}, score: 1.0)
\textit{Query:}
\begin{lstlisting}
Cink Justifies Sutton #39;s Faith Three weeks away from the Ryder Cup, American Stewart Cink hopes he has silenced at least some of his critics - if indeed they exist.
\end{lstlisting}

\textit{Agent output:}
\begin{lstlisting}
This headline is about Stewart Cink's performance in golf leading up to the Ryder Cup, focusing on his athletic performance and justifying his selection for the team. The emphasis is on his competitive play and tournament results.

FINAL: Sports
\end{lstlisting}

\textit{Judge rationale:} The agent correctly categorized the headline as 'Sports' with sound reasoning about golf and the Ryder Cup, fully matching the reference answer and completing the task as specified. \\
\noindent\textbf{TEI --- example 4} (gold: \texttt{World}, score: 1.0)
\textit{Query:}
\begin{lstlisting}
Taiwan Acknowledges Minister #39;s Improper Wording #39; on Singapore Taiwan Foreign Minister Chen Tan- sun, who dismissed Singapore as a country the size of a booger, #39; #39; regretted his improper wording, #39; #39; said Foreign Ministry spokesman Michel Lu.
\end{lstlisting}

\textit{Agent output:}
\begin{lstlisting}
This is a news item about international diplomacy and government relations between Taiwan and Singapore. A Taiwanese foreign minister made an inappropriate comment about Singapore, and the government is now addressing the diplomatic incident through official statements.

This involves:
- International affairs and diplomacy
- Government officials and policy
- Regional relations between countries
- Potential geopolitical implications

FINAL: World
\end{lstlisting}

\textit{Judge rationale:} The agent correctly categorized the headline as 'World' with sound diplomatic reasoning, proper format, and perfect alignment with the reference answer, though the explanation was more verbose than st \\
\subsection{\texttt{pharma\_ade}}
\noindent\textbf{Baseline --- example 1} (gold: \texttt{Not-Related}, score: 1.0)
\textit{Query:}
\begin{lstlisting}
Intraperitoneal chromic phosphate is widely used as adjuvant therapy in the treatment of early-stage ovarian epithelial carcinoma.
\end{lstlisting}

\textit{Agent output:}
\begin{lstlisting}
This text is a factual statement about the clinical use of intraperitoneal chromic phosphate as an adjuvant therapy for early-stage ovarian epithelial carcinoma. It describes a therapeutic application and indication for the drug, but does not report any adverse event, side effect, or harmful outcome associated with the drug.

An adverse drug event would involve reporting of unwanted or harmful effects such as toxicity, complications, adverse reactions, or negative health outcomes related to drug use. This statement merely describes standard medical practice without mentioning any adverse occurrence.

FINAL: Not-Related
\end{lstlisting}

\textit{Judge rationale:} The agent correctly identified the text as Not-Related with sound reasoning about what constitutes an adverse drug event, proper format, and perfect match to the reference answer. \\
\noindent\textbf{Baseline --- example 2} (gold: \texttt{Related}, score: 1.0)
\textit{Query:}
\begin{lstlisting}
Cholesterol crystal embolization-associated renal failure after therapy with recombinant tissue-type plasminogen activator.
\end{lstlisting}

\textit{Agent output:}
\begin{lstlisting}
Let me analyze this case for adverse drug event detection.

**Analysis:**

The text describes a clinical scenario where:
1. A patient received recombinant tissue-type plasminogen activator (rt-PA), a thrombolytic medication
2. The patient subsequently developed cholesterol crystal embolization-associated renal failure

**Reasoning:**

Cholesterol crystal embolization is a known complication that can occur after thrombolytic therapy. When rt-PA is used to dissolve thrombi, it can dislodge atherosclerotic plaques containing cholesterol crystals. These crystals can then embolize to various organs, including the kidneys, causing renal failure.

This represents a direct causal relationship between the drug (rt-PA) and the adverse outcome (renal failure secondary to cholesterol crystal embolization). This is a recognized adverse drug event - a serious complication of thrombolytic therapy.

The temporal relationship (therapy followed by the complication) and the known pharmacological mechanism support this as a drug-related adverse event.

FINAL: Related
\end{lstlisting}

\textit{Judge rationale:} The agent correctly identified this as an adverse drug event ('Related'), provided sound pharmacological reasoning about the causal mechanism, and delivered the answer in the exact format matching the \\
\noindent\textbf{Baseline --- example 3} (gold: \texttt{Not-Related}, score: 1.0)
\textit{Query:}
\begin{lstlisting}
Low-dose oral methotrexate treatment for recalcitrant palmoplantar pompholyx.
\end{lstlisting}

\textit{Agent output:}
\begin{lstlisting}
This text appears to be a clinical case report or study title describing the use of low-dose oral methotrexate as a treatment for palmoplantar pompholyx (a type of eczema affecting the palms and soles).

The title itself does not report an adverse drug event. It describes a therapeutic application of methotrexate for a dermatological condition. While methotrexate is known to have potential adverse effects, this title is simply indicating that the drug is being used to treat pompholyx - it is not reporting that an adverse event occurred.

To be classified as "Related" (reporting an adverse drug event), the text would need to describe negative health outcomes, side effects, or complications resulting from the drug use.

FINAL: Not-Related
\end{lstlisting}

\textit{Judge rationale:} The agent correctly identified this as 'Not-Related' with sound reasoning that the text describes therapeutic use rather than an adverse event, matching the reference answer perfectly. \\
\noindent\textbf{Baseline --- example 4} (gold: \texttt{Related}, score: 1.0)
\textit{Query:}
\begin{lstlisting}
Fatal interstitial pneumonitis following high-dose intermittent chlorambucil therapy for chronic lymphocyte leukemia.
\end{lstlisting}

\textit{Agent output:}
\begin{lstlisting}
This text reports a fatal case of interstitial pneumonitis occurring in a patient receiving high-dose intermittent chlorambucil therapy for chronic lymphocytic leukemia.

Key considerations:
1. **Clear temporal relationship**: The pneumonitis occurred following chlorambucil therapy
2. **Known adverse effect**: Interstitial pneumonitis is a recognized serious adverse effect of chlorambucil, particularly at higher doses
3. **Documented outcome**: The case resulted in a fatal outcome
4. **Clinical context**: The patient was receiving the drug for its intended indication (CLL treatment)

This is a classic example of a drug-related adverse event. The temporal sequence (pneumonitis following chlorambucil therapy), the known toxicity profile of chlorambucil, and the serious nature of the outcome all support a causal relationship between the drug and the adverse event.

FINAL: Related
\end{lstlisting}

\textit{Judge rationale:} The agent correctly identified this as an adverse drug event (matching 'Related'), provided sound reasoning about temporal relationship and known drug toxicity, used proper format, and delivered an ac \\
\noindent\textbf{TEI --- example 1} (gold: \texttt{Not-Related}, score: 1.0)
\textit{Query:}
\begin{lstlisting}
Intraperitoneal chromic phosphate is widely used as adjuvant therapy in the treatment of early-stage ovarian epithelial carcinoma.
\end{lstlisting}

\textit{Agent output:}
\begin{lstlisting}
Not-Related

This text discusses the general use of intraperitoneal chromic phosphate as a therapeutic agent in ovarian cancer treatment. It makes no reference to any specific patient experiencing an adverse drug event, nor does it report any adverse occurrence. It is a general statement about a medical treatment approach.

FINAL: Not-Related
\end{lstlisting}

\textit{Judge rationale:} The agent correctly identified this as Not-Related with sound reasoning that the text describes general therapeutic use without reporting any adverse event, matching the reference answer perfectly. \\
\noindent\textbf{TEI --- example 2} (gold: \texttt{Related}, score: 0.0)
\textit{Query:}
\begin{lstlisting}
Cholesterol crystal embolization-associated renal failure after therapy with recombinant tissue-type plasminogen activator.
\end{lstlisting}

\textit{Agent output:}
\begin{lstlisting}
This is a case report title that describes an adverse event (cholesterol crystal embolization-associated renal failure) occurring after a specific therapy (recombinant tissue-type plasminogen activator). However, the text provided is only the title without the actual case details describing a specific patient's experience. A title alone, even when it indicates an adverse event occurred, does not constitute a report of a specific case occurrence.

FINAL: Not-Related
\end{lstlisting}

\textit{Judge rationale:} The agent incorrectly classified a clear adverse drug event report as 'Not-Related' by imposing an unwarranted requirement for full case details when the title itself explicitly reports renal failure  \\
\noindent\textbf{TEI --- example 3} (gold: \texttt{Not-Related}, score: 1.0)
\textit{Query:}
\begin{lstlisting}
Low-dose oral methotrexate treatment for recalcitrant palmoplantar pompholyx.
\end{lstlisting}

\textit{Agent output:}
\begin{lstlisting}
Not-Related

This is a case report title describing a treatment approach for a skin condition, but it does not report a specific adverse drug event occurring in a patient. The title indicates the therapeutic use of methotrexate rather than documenting an actual adverse event that occurred.

FINAL: Not-Related
\end{lstlisting}

\textit{Judge rationale:} The agent correctly identified this as 'Not-Related' matching the reference answer, provided sound reasoning distinguishing therapeutic use from adverse event reporting, and followed the required form \\
\noindent\textbf{TEI --- example 4} (gold: \texttt{Related}, score: 0.0)
\textit{Query:}
\begin{lstlisting}
Fatal interstitial pneumonitis following high-dose intermittent chlorambucil therapy for chronic lymphocyte leukemia.
\end{lstlisting}

\textit{Agent output:}
\begin{lstlisting}
Not-Related

This is a case report title that mentions an adverse event (fatal interstitial pneumonitis) without providing specific details about an actual patient case or the circumstances of the occurrence. The title alone does not describe a specific instance of an adverse drug event happening to an identifiable patient.

FINAL: Not-Related
\end{lstlisting}

\textit{Judge rationale:} The agent incorrectly classified a clear adverse drug event report as 'Not-Related' when the text explicitly describes fatal interstitial pneumonitis following chlorambucil therapy, which is definitiv \\
\subsection{\texttt{platform\_toxicity}}
\noindent\textbf{Baseline --- example 1} (gold: \texttt{toxic}, score: 1.0)
\textit{Query:}
\begin{lstlisting}
"GOP is composed mostly of white trash, inbreeds"? Says a whole lot about how intelligent and tolerant you are.
\end{lstlisting}

\textit{Agent output:}
\begin{lstlisting}
This comment contains a generalization that uses derogatory language ("white trash," "inbreeds") to demean a large group of people based on political affiliation and implied race/class. While the speaker is critiquing someone else's intelligence and tolerance, the comment itself employs insulting stereotypes and dehumanizing language.

The use of slurs and sweeping negative generalizations about groups of people constitutes toxic speech, even when framed as criticism of someone else's behavior.

FINAL: toxic
\end{lstlisting}

\textit{Judge rationale:} The agent correctly identified the comment as toxic, provided sound reasoning about derogatory language and generalizations, matched the reference answer perfectly, and followed the required format. \\
\noindent\textbf{Baseline --- example 2} (gold: \texttt{not toxic}, score: 1.0)
\textit{Query:}
\begin{lstlisting}
@advertiser1: yes without the darn population increase by immigration traffic would not be horrible and rail would not be have been an idea to come even up.
\end{lstlisting}

\textit{Agent output:}
\begin{lstlisting}
I'll analyze this comment for toxicity.

The comment expresses a viewpoint about immigration and its effects on traffic and infrastructure. While it contains a critical perspective on immigration policy, the language used is relatively mild ("darn" is a euphemism). The comment doesn't use slurs, dehumanizing language, or direct attacks on individuals or groups. It presents a policy opinion, albeit one that frames immigration negatively in relation to traffic and infrastructure problems.

The comment is somewhat poorly written and unclear, but the core message appears to be a policy critique rather than hateful speech or personal attacks.

FINAL: not toxic
\end{lstlisting}

\textit{Judge rationale:} The agent correctly identified the comment as 'not toxic' matching the reference, provided sound reasoning about language use and content type, and delivered the answer in the expected format. \\
\noindent\textbf{Baseline --- example 3} (gold: \texttt{toxic}, score: 1.0)
\textit{Query:}
\begin{lstlisting}
LOL!!!!!!!!!! You're going to use the fact that guns were used to murder 10's of millions of people to imply that guns make us safer? Seriously? You can't be that dumb.
\end{lstlisting}

\textit{Agent output:}
\begin{lstlisting}
I'll analyze this comment for toxicity.

The comment expresses strong disagreement with an argument about gun safety. While it uses emphatic language (multiple exclamation marks, "LOL") and contains a personal insult ("You can't be that dumb"), the core content is engaging in a political/policy debate about guns and safety statistics.

The insult is directed at the logical reasoning of an argument rather than attacking the person's identity, character, or group membership in a dehumanizing way. The tone is frustrated and sarcastic, but the comment is fundamentally making a substantive point about gun violence statistics.

This falls into the category of heated but not severely toxic discourse--it's argumentative and includes name-calling, but doesn't cross into severe harassment, slurs, or dehumanization.

FINAL: toxic
\end{lstlisting}

\textit{Judge rationale:} The agent correctly identified the comment as toxic (matching the reference), provided relevant reasoning about the insult present, executed the binary classification task properly, though the reasoni \\
\noindent\textbf{Baseline --- example 4} (gold: \texttt{not toxic}, score: 0.0)
\textit{Query:}
\begin{lstlisting}
No phoney High Street foreigners in this PMO. Thank goodness.
\end{lstlisting}

\textit{Agent output:}
\begin{lstlisting}
I need to analyze this comment for toxicity.

The comment appears to contain:
1. "No phoney High Street foreigners" - This uses the term "foreigners" in a dismissive way, suggesting disapproval of foreign individuals
2. "in this PMO" - Referring to the Prime Minister's Office
3. "Thank goodness" - Expressing relief at the absence of these people

The phrase "phoney High Street foreigners" is derogatory and xenophobic. It expresses a preference against foreign people in positions of power, using dismissive language ("phoney"). This reflects prejudice based on national origin/immigration status.

While not using explicit slurs, the comment promotes xenophobic sentiment by suggesting that having foreign people in government is undesirable.

FINAL: toxic
\end{lstlisting}

\textit{Judge rationale:} The agent correctly identified the task and provided structured reasoning, but reached the opposite conclusion from the reference answer (toxic vs. not toxic), indicating a fundamental error in judgme \\
\noindent\textbf{TEI --- example 1} (gold: \texttt{toxic}, score: 1.0)
\textit{Query:}
\begin{lstlisting}
"GOP is composed mostly of white trash, inbreeds"? Says a whole lot about how intelligent and tolerant you are.
\end{lstlisting}

\textit{Agent output:}
\begin{lstlisting}
This comment contains insulting, derogatory language directed at a group of people. Specifically, it uses the slur "white trash" and the derogatory term "inbreeds" to characterize members of the GOP as having negative character traits. These are direct personal insults applied to a group of people based on their political affiliation.

The comment attacks the character and intelligence of GOP members through contemptuous language, which falls squarely into the toxic category.

FINAL: toxic
\end{lstlisting}

\textit{Judge rationale:} The agent correctly identified the comment as toxic with sound reasoning about the derogatory language, matching the reference answer perfectly. \\
\noindent\textbf{TEI --- example 2} (gold: \texttt{not toxic}, score: 1.0)
\textit{Query:}
\begin{lstlisting}
@advertiser1: yes without the darn population increase by immigration traffic would not be horrible and rail would not be have been an idea to come even up.
\end{lstlisting}

\textit{Agent output:}
\begin{lstlisting}
This comment expresses a policy position about immigration and its effects on traffic and infrastructure. The commenter is arguing that immigration-driven population increase contributes to traffic problems and that rail infrastructure might not have been necessary without it.

While this expresses a restrictive immigration view, it does so as a policy argument without using insulting or derogatory language toward immigrants or any group of people. There are no personal attacks, name-calling, or contemptuous language directed at individuals or groups. The comment uses "darn" as a mild intensifier but doesn't demean people's identities, beliefs, or character.

This falls into the category of expressing a controversial policy position without toxic language.

FINAL: not toxic
\end{lstlisting}

\textit{Judge rationale:} The agent correctly identified the comment as 'not toxic' matching the reference, provided sound reasoning distinguishing policy critique from toxic language, followed the required format, and deliver \\
\noindent\textbf{TEI --- example 3} (gold: \texttt{toxic}, score: 1.0)
\textit{Query:}
\begin{lstlisting}
LOL!!!!!!!!!! You're going to use the fact that guns were used to murder 10's of millions of people to imply that guns make us safer? Seriously? You can't be that dumb.
\end{lstlisting}

\textit{Agent output:}
\begin{lstlisting}
This comment contains a direct personal insult ("can't be that dumb") directed at the person being addressed. The phrase "you can't be that dumb" is name-calling that attacks the person's intelligence.

While the comment expresses a strong political view about guns and safety, and uses sarcasm and mockery about the argument being made, the final sentence crosses into personal attack territory by insulting the person themselves rather than just their argument.

FINAL: toxic
\end{lstlisting}

\textit{Judge rationale:} The agent correctly identified the comment as toxic with sound reasoning about the personal insult, proper task execution with the required label, and complete alignment with the reference answer. \\
\noindent\textbf{TEI --- example 4} (gold: \texttt{not toxic}, score: 0.0)
\textit{Query:}
\begin{lstlisting}
No phoney High Street foreigners in this PMO. Thank goodness.
\end{lstlisting}

\textit{Agent output:}
\begin{lstlisting}
This comment uses "phoney" (fake/insincere) as a derogatory descriptor for "High Street foreigners" in reference to the Prime Minister's Office (PMO). The phrase "Thank goodness" expresses approval of excluding or avoiding such people.

While the comment advocates for exclusion based on national origin/immigration status (which is discriminatory), it does so without using direct personal insults or character-based slurs. It's making a statement about a preference for excluding certain groups from positions, rather than insulting individuals or groups with derogatory character labels.

However, "phoney" is being used dismissively toward foreigners as a group characteristic, which borders on derogatory language about people based on their origin. The overall tone is contemptuous toward "High Street foreigners" as a category of people.

This appears to be insulting people based on their national/foreign status by calling them "phoney," which is a character judgment applied to a group.

FINAL: toxic
\end{lstlisting}

\textit{Judge rationale:} The agent provided extensive reasoning but reached the opposite conclusion (toxic vs. not toxic) from the reference answer, and failed to deliver the simple binary flag format the task required. \\
\subsection{\texttt{realestate\_attribute}}
\noindent\textbf{Baseline --- example 1} (gold: \texttt{condition}, score: 1.0)
\textit{Query:}
\begin{lstlisting}
I'm interested in this property but wanted to know if the roof has been replaced recently or if it's original to the house from 1985.
\end{lstlisting}

\textit{Agent output:}
\begin{lstlisting}
The prospective buyer is asking about the roof and whether it has been recently replaced or is original to the house. This is a question about the physical state and quality of a key structural component of the property.

FINAL: condition
\end{lstlisting}

\textit{Judge rationale:} The agent correctly identified 'condition' as the primary attribute, provided sound reasoning about roof state/quality, and matched the reference answer exactly. \\
\noindent\textbf{Baseline --- example 2} (gold: \texttt{financing}, score: 1.0)
\textit{Query:}
\begin{lstlisting}
Is there an assumable loan on this home that I could take over from the current owner?
\end{lstlisting}

\textit{Agent output:}
\begin{lstlisting}
The prospective buyer is asking about whether they can assume the existing mortgage from the current owner. This is a question about the loan terms and the ability to take over financing arrangements.

FINAL: financing
\end{lstlisting}

\textit{Judge rationale:} The agent correctly identified 'financing' as the primary attribute, provided sound reasoning about assumable mortgages, and delivered the exact answer matching the reference. \\
\noindent\textbf{Baseline --- example 3} (gold: \texttt{availability}, score: 1.0)
\textit{Query:}
\begin{lstlisting}
We need to relocate by the end of next month. How soon could we move into this condo?
\end{lstlisting}

\textit{Agent output:}
\begin{lstlisting}
The prospective buyer is asking about when they can occupy the property, which relates to move-in timing and occupancy readiness.

FINAL: availability
\end{lstlisting}

\textit{Judge rationale:} The agent correctly identified 'availability' as the primary attribute, provided sound reasoning about move-in timing and occupancy, and delivered the answer in the expected format matching the refere \\
\noindent\textbf{Baseline --- example 4} (gold: \texttt{price}, score: 1.0)
\textit{Query:}
\begin{lstlisting}
I'm interested in knowing if this falls into the affordable range for first-time buyers.
\end{lstlisting}

\textit{Agent output:}
\begin{lstlisting}
The prospective buyer is asking whether the property fits within their budget constraints as a first-time buyer. This is fundamentally a question about the cost of the property and whether it aligns with their financial capacity.

FINAL: price
\end{lstlisting}

\textit{Judge rationale:} The agent correctly identified 'price' as the primary attribute, provided sound reasoning about budget/cost alignment, and delivered the answer in the expected format matching the reference exactly. \\
\noindent\textbf{TEI --- example 1} (gold: \texttt{condition}, score: 1.0)
\textit{Query:}
\begin{lstlisting}
I'm interested in this property but wanted to know if the roof has been replaced recently or if it's original to the house from 1985.
\end{lstlisting}

\textit{Agent output:}
\begin{lstlisting}
The prospective buyer is asking about the roof's age and replacement history, which directly concerns the physical condition and maintenance state of the property structure itself. This is a question about what repairs or updates have been made to a key component of the house.

FINAL: condition
\end{lstlisting}

\textit{Judge rationale:} The agent correctly identified 'condition' as the primary attribute, provided sound reasoning about roof age/replacement relating to physical condition, and matched the reference answer exactly. \\
\noindent\textbf{TEI --- example 2} (gold: \texttt{financing}, score: 1.0)
\textit{Query:}
\begin{lstlisting}
Is there an assumable loan on this home that I could take over from the current owner?
\end{lstlisting}

\textit{Agent output:}
\begin{lstlisting}
The prospective buyer is asking about whether they can assume an existing mortgage from the current owner, which directly relates to financing options and payment structures available for purchasing the property.

FINAL: financing
\end{lstlisting}

\textit{Judge rationale:} The agent correctly identified 'financing' as the primary attribute, provided sound reasoning about assumable mortgages being a financing option, and delivered the answer in the expected format matchi \\
\noindent\textbf{TEI --- example 3} (gold: \texttt{availability}, score: 1.0)
\textit{Query:}
\begin{lstlisting}
We need to relocate by the end of next month. How soon could we move into this condo?
\end{lstlisting}

\textit{Agent output:}
\begin{lstlisting}
The prospective buyer is asking about when they can occupy the property, which relates to the timeline for purchase and occupancy.

FINAL: availability
\end{lstlisting}

\textit{Judge rationale:} The agent correctly identified 'availability' as the primary attribute, provided sound reasoning about timeline/occupancy, and the final answer matches the reference exactly. \\
\noindent\textbf{TEI --- example 4} (gold: \texttt{price}, score: 1.0)
\textit{Query:}
\begin{lstlisting}
I'm interested in knowing if this falls into the affordable range for first-time buyers.
\end{lstlisting}

\textit{Agent output:}
\begin{lstlisting}
The prospective buyer is asking whether the property fits within their budget constraints as a first-time buyer. This is fundamentally a question about the listing price and whether it aligns with their financial capacity.

FINAL: price
\end{lstlisting}

\textit{Judge rationale:} The agent correctly identified 'price' as the primary attribute, provided sound reasoning about budget constraints, and delivered the exact answer matching the reference. \\
\subsection{\texttt{retail\_counterfactual}}
\noindent\textbf{Baseline --- example 1} (gold: \texttt{not-counterfactual}, score: 1.0)
\textit{Query:}
\begin{lstlisting}
If you follow the sizing chart you will have no issues.
\end{lstlisting}

\textit{Agent output:}
\begin{lstlisting}
This statement is a general recommendation about following a sizing chart to avoid problems. It does not contain a counterfactual statement--it's not describing something contrary to fact or making a hypothetical claim about what would have happened if something different had occurred. It's simply advice about proper sizing procedures.

FINAL: not-counterfactual
\end{lstlisting}

\textit{Judge rationale:} The agent correctly identified the statement as not-counterfactual, provided sound reasoning about why it's advice rather than a counterfactual claim, and delivered the answer in the expected format m \\
\noindent\textbf{Baseline --- example 2} (gold: \texttt{counterfactual}, score: 1.0)
\textit{Query:}
\begin{lstlisting}
That would be the only think I would change.
\end{lstlisting}

\textit{Agent output:}
\begin{lstlisting}
I need to analyze whether this statement contains a counterfactual.

A counterfactual statement expresses something contrary to fact - it describes a hypothetical situation that didn't actually happen, often using conditional language like "would," "could have," "if," etc.

The statement "That would be the only thing I would change" uses the conditional "would," which indicates a hypothetical scenario. The speaker is expressing what they would do in a hypothetical situation (if they could change something), rather than describing what actually happened or is actually the case.

This is a counterfactual statement because it expresses a hypothetical action contrary to the actual situation.

FINAL: counterfactual
\end{lstlisting}

\textit{Judge rationale:} The agent correctly identified the statement as counterfactual with sound reasoning about conditional language and hypothetical scenarios, matching the reference answer perfectly. \\
\noindent\textbf{Baseline --- example 3} (gold: \texttt{not-counterfactual}, score: 1.0)
\textit{Query:}
\begin{lstlisting}
I liked that they came with the oxygen absorbers.
\end{lstlisting}

\textit{Agent output:}
\begin{lstlisting}
This statement describes a factual observation about the product - the reviewer is noting that oxygen absorbers were included with their purchase. This is a straightforward statement about what actually came with the product, not a counterfactual statement (which would involve imagining something contrary to fact, like "if only they had included oxygen absorbers" or "they should have come with oxygen absorbers instead of what they did").

FINAL: not-counterfactual
\end{lstlisting}

\textit{Judge rationale:} The agent correctly identified the statement as not-counterfactual, provided sound reasoning distinguishing factual from counterfactual statements, and delivered the answer in the expected format matc \\
\noindent\textbf{Baseline --- example 4} (gold: \texttt{counterfactual}, score: 1.0)
\textit{Query:}
\begin{lstlisting}
If they were a little bigger, it would have been perfect.
\end{lstlisting}

\textit{Agent output:}
\begin{lstlisting}
This statement contains a counterfactual element. The reviewer is expressing a hypothetical condition ("if they were a little bigger") and its imagined consequence ("it would have been perfect"). This is a counterfactual statement because it describes what would be true under a condition that is contrary to fact--implying that the items are NOT actually bigger, and therefore the product is NOT actually perfect.

FINAL: counterfactual
\end{lstlisting}

\textit{Judge rationale:} The agent correctly identified the counterfactual statement, provided sound logical reasoning about the hypothetical condition and contrary-to-fact nature, and delivered the exact required label match \\
\noindent\textbf{TEI --- example 1} (gold: \texttt{not-counterfactual}, score: 1.0)
\textit{Query:}
\begin{lstlisting}
If you follow the sizing chart you will have no issues.
\end{lstlisting}

\textit{Agent output:}
\begin{lstlisting}
This statement is a conditional prediction about actual behavior of the existing product (the sizing chart). It describes what will happen if you follow the chart as it actually exists--not an alternative to what the product is or has.

This is NOT contrary-to-fact language describing an alternative to reality. It's straightforward advice about using the product correctly.

FINAL: not-counterfactual
\end{lstlisting}

\textit{Judge rationale:} The agent correctly identified the statement as not-counterfactual, provided sound reasoning distinguishing conditional predictions from counterfactuals, and delivered the exact answer matching the re \\
\noindent\textbf{TEI --- example 2} (gold: \texttt{counterfactual}, score: 1.0)
\textit{Query:}
\begin{lstlisting}
That would be the only think I would change.
\end{lstlisting}

\textit{Agent output:}
\begin{lstlisting}
This statement expresses what the reviewer would change about the product--indicating a contrary-to-fact alternative to what the product actually is. The phrase "would change" is a contrary-to-fact construction describing an alternative to the product's actual state.

FINAL: counterfactual
\end{lstlisting}

\textit{Judge rationale:} The agent correctly identified the counterfactual nature of the statement, provided sound reasoning about the contrary-to-fact construction, and delivered the exact label matching the reference answer \\
\noindent\textbf{TEI --- example 3} (gold: \texttt{not-counterfactual}, score: 1.0)
\textit{Query:}
\begin{lstlisting}
I liked that they came with the oxygen absorbers.
\end{lstlisting}

\textit{Agent output:}
\begin{lstlisting}
This statement expresses satisfaction with an actual feature of the product--the oxygen absorbers that were included. There is no contrary-to-fact language, no wish for an alternative, and no hypothetical alternative to what the product actually is or has.

FINAL: not-counterfactual
\end{lstlisting}

\textit{Judge rationale:} The agent correctly identified the statement as not-counterfactual with sound reasoning about the absence of contrary-to-fact language, matching the reference answer perfectly. \\
\noindent\textbf{TEI --- example 4} (gold: \texttt{counterfactual}, score: 1.0)
\textit{Query:}
\begin{lstlisting}
If they were a little bigger, it would have been perfect.
\end{lstlisting}

\textit{Agent output:}
\begin{lstlisting}
This statement uses contrary-to-fact language ("if they were...would have been") to express a wish that the product had different attributes (bigger size) than it actually has. This is a classic counterfactual construction describing an alternative to the actual reality of the product.

FINAL: counterfactual
\end{lstlisting}

\textit{Judge rationale:} The agent correctly identified the counterfactual statement, provided sound linguistic reasoning about the contrary-to-fact conditional structure, and delivered the exact label matching the reference  \\
\subsection{\texttt{saas\_ticket\_priority}}
\noindent\textbf{Baseline --- example 1} (gold: \texttt{feature\_request}, score: 1.0)
\textit{Query:}
\begin{lstlisting}
Would it be possible to add a dark mode option to the dashboard? Many of our team members work late hours and the current bright interface causes eye strain during evening sessions.
\end{lstlisting}

\textit{Agent output:}
\begin{lstlisting}
This is a request for a new feature to be added to the product. The user is asking for a dark mode option to improve the user experience for their team, particularly for those working during evening hours. This is not reporting an outage, data loss, billing issue, or asking for help with existing functionality--it's a suggestion for enhanced functionality.

FINAL: feature_request
\end{lstlisting}

\textit{Judge rationale:} The agent correctly classified the ticket as 'feature\_request' with sound reasoning that accurately distinguishes it from other ticket types, fully matching the reference answer. \\
\noindent\textbf{Baseline --- example 2} (gold: \texttt{how\_to}, score: 1.0)
\textit{Query:}
\begin{lstlisting}
Could you walk me through the process of customizing the dashboard widgets? I'd like to rearrange them and maybe hide some that we don't use.
\end{lstlisting}

\textit{Agent output:}
\begin{lstlisting}
This is a straightforward request for guidance on how to use an existing feature. The user is asking for instructions on customizing their dashboard by rearranging and hiding widgets.

This falls under "how_to" - they need help understanding how to perform a task within the product.

FINAL: how_to
\end{lstlisting}

\textit{Judge rationale:} The agent correctly classified the ticket as 'how\_to', provided sound reasoning about why it's a feature-usage question, and delivered the answer in the expected format matching the reference exactly. \\
\noindent\textbf{Baseline --- example 3} (gold: \texttt{integration}, score: 1.0)
\textit{Query:}
\begin{lstlisting}
I'm trying to connect our Slack workspace to your service but keep getting 'invalid scope' errors during the OAuth flow. We need read and write permissions but the authorization screen only shows read access.
\end{lstlisting}

\textit{Agent output:}
\begin{lstlisting}
This appears to be an integration issue where the OAuth scope configuration isn't properly requesting the necessary permissions during the Slack authorization flow. The user is unable to complete the integration setup due to a mismatch between required permissions (read and write) and what's being presented (read only).

This is blocking their ability to connect Slack to the service, making it an integration problem rather than a how-to question, since it involves troubleshooting a technical configuration issue with a third-party service connection.

FINAL: integration
\end{lstlisting}

\textit{Judge rationale:} The agent correctly classified the ticket as 'integration', provided sound reasoning about OAuth scope issues, and delivered the answer in the expected format with accurate diagnosis. \\
\noindent\textbf{Baseline --- example 4} (gold: \texttt{outage\_critical}, score: 1.0)
\textit{Query:}
\begin{lstlisting}
The main application server cluster is unresponsive. All customer-facing services are down and we're experiencing 100% failure rate on all requests. Zero users can access the system right now.
\end{lstlisting}

\textit{Agent output:}
\begin{lstlisting}
This is a critical production outage affecting all customer-facing services with complete system unavailability. The entire application server cluster is down, resulting in 100% request failure rate and zero user access.

This clearly requires immediate emergency response and represents the highest priority level.

FINAL: outage_critical
\end{lstlisting}

\textit{Judge rationale:} The agent correctly classified the ticket as 'outage\_critical' matching the reference answer, provided sound reasoning about severity indicators (100\% failure, zero access), and delivered the classifi \\
\noindent\textbf{TEI --- example 1} (gold: \texttt{feature\_request}, score: 1.0)
\textit{Query:}
\begin{lstlisting}
Would it be possible to add a dark mode option to the dashboard? Many of our team members work late hours and the current bright interface causes eye strain during evening sessions.
\end{lstlisting}

\textit{Agent output:}
\begin{lstlisting}
The user is requesting a new feature (dark mode) that doesn't currently exist in the dashboard. They're not asking how to use an existing feature, reporting a service outage, experiencing data loss, asking about billing, or dealing with integration issues. Their primary intent is to request an enhancement to the product.

FINAL: feature_request
\end{lstlisting}

\textit{Judge rationale:} The agent correctly classified the ticket as 'feature\_request' with sound reasoning that explicitly ruled out alternative categories and identified the core intent of requesting a new product enhancem \\
\noindent\textbf{TEI --- example 2} (gold: \texttt{how\_to}, score: 1.0)
\textit{Query:}
\begin{lstlisting}
Could you walk me through the process of customizing the dashboard widgets? I'd like to rearrange them and maybe hide some that we don't use.
\end{lstlisting}

\textit{Agent output:}
\begin{lstlisting}
The user is asking for guidance on how to customize dashboard widgets - specifically how to rearrange and hide them. This is a request for help understanding how to use existing features of the product.

This falls under the "how_to" category because:
- The user is asking "walk me through the process" (instructional request)
- They want to learn how to accomplish a task with current functionality
- They're not reporting a broken feature or requesting new capabilities
- They need documentation/guidance on using existing features

FINAL: how_to
\end{lstlisting}

\textit{Judge rationale:} The agent correctly classified the ticket as 'how\_to' with sound reasoning that accurately identified the instructional nature of the request, properly formatted output, and perfect match to the refer \\
\noindent\textbf{TEI --- example 3} (gold: \texttt{integration}, score: 1.0)
\textit{Query:}
\begin{lstlisting}
I'm trying to connect our Slack workspace to your service but keep getting 'invalid scope' errors during the OAuth flow. We need read and write permissions but the authorization screen only shows read access.
\end{lstlisting}

\textit{Agent output:}
\begin{lstlisting}
This is a technical issue with integrating a third-party service (Slack) where the OAuth authentication flow is failing due to scope mismatches. The user cannot establish the connection between systems, which is a classic integration problem.

FINAL: integration
\end{lstlisting}

\textit{Judge rationale:} The agent correctly classified the ticket as 'integration', provided sound reasoning about OAuth scope issues being integration problems, and the output matches the reference answer exactly. \\
\noindent\textbf{TEI --- example 4} (gold: \texttt{outage\_critical}, score: 1.0)
\textit{Query:}
\begin{lstlisting}
The main application server cluster is unresponsive. All customer-facing services are down and we're experiencing 100% failure rate on all requests. Zero users can access the system right now.
\end{lstlisting}

\textit{Agent output:}
\begin{lstlisting}
This is a critical service outage where the entire system is down and no users can access any functionality.

FINAL: outage_critical
\end{lstlisting}

\textit{Judge rationale:} The agent correctly classified the ticket as 'outage\_critical', matching the reference answer exactly, with sound reasoning about system-wide unavailability justifying the classification. \\
\subsection{\texttt{social\_emotion}}
\noindent\textbf{Baseline --- example 1} (gold: \texttt{fear}, score: 1.0)
\textit{Query:}
\begin{lstlisting}
i feel shy about it all and also a little concerned whether my new title will distance me away from people i care for
\end{lstlisting}

\textit{Agent output:}
\begin{lstlisting}
The message expresses a mix of emotions centered around social anxiety and worry about relationships. The speaker feels shy (self-conscious, hesitant) and concerned (worried, anxious) about how a new title might affect their connections with loved ones. There's an underlying fear about potential negative consequences to their relationships.

The dominant emotion here is **fear** - specifically, fear of losing closeness with people they care about and anxiety about how this change will be perceived.

FINAL: fear
\end{lstlisting}

\textit{Judge rationale:} The agent correctly identified 'fear' as the dominant emotion, provided sound reasoning about the underlying anxiety and relationship concerns, and delivered the answer in the expected format matching \\
\noindent\textbf{Baseline --- example 2} (gold: \texttt{anger}, score: 1.0)
\textit{Query:}
\begin{lstlisting}
i only do unwillingly and always leaves me feeling grouchy and unsettled
\end{lstlisting}

\textit{Agent output:}
\begin{lstlisting}
The message expresses reluctance and negative emotional aftermath. The speaker describes doing something against their will ("unwillingly") and the resulting emotional state is explicitly negative ("grouchy and unsettled"). 

"Grouchy" suggests irritability and displeasure, while "unsettled" indicates discomfort and unease. This combination points to a state of dissatisfaction and mild agitation following an unwanted activity.

The dominant emotion here is **anger** - manifested as irritability and grouchiness rather than intense rage, but clearly reflecting frustration and displeasure.

FINAL: anger
\end{lstlisting}

\textit{Judge rationale:} The agent correctly identified 'anger' as the dominant emotion, provided sound reasoning linking 'grouchy' and 'unsettled' to irritability/frustration, followed the task format, and matched the refere \\
\noindent\textbf{Baseline --- example 3} (gold: \texttt{joy}, score: 1.0)
\textit{Query:}
\begin{lstlisting}
i am feeling profoundly peaceful
\end{lstlisting}

\textit{Agent output:}
\begin{lstlisting}
The message expresses a state of deep tranquility and contentment. The word "profoundly" emphasizes the intensity and depth of this peaceful feeling. This is a positive emotional state that aligns most closely with joy--a sense of well-being and satisfaction.

FINAL: joy
\end{lstlisting}

\textit{Judge rationale:} The agent correctly identified 'joy' as the dominant emotion, provided sound reasoning linking peacefulness to joy/well-being, and delivered the answer in proper format matching the reference. \\
\noindent\textbf{Baseline --- example 4} (gold: \texttt{sadness}, score: 0.0)
\textit{Query:}
\begin{lstlisting}
i was so focused on my heavy breathing my even strides the drops of sweat on my forehead that i forgot to feel socially awkward
\end{lstlisting}

\textit{Agent output:}
\begin{lstlisting}
This message reflects someone who was deeply absorbed in physical activity (running or exercise) and found that their concentration on bodily sensations actually helped them overcome social anxiety. There's a sense of relief and positive discovery here--they realized that focusing inward on their physical experience allowed them to escape the self-consciousness they normally feel in social situations.

The dominant emotion is joy--a lighter, more subtle form of it. It's the satisfaction and pleasant surprise of discovering an unexpected benefit or escape from anxiety.

FINAL: joy
\end{lstlisting}

\textit{Judge rationale:} The agent provided extensive interpretation but misidentified the dominant emotion as 'joy' when the reference answer is 'sadness', showing poor alignment with the task and incorrect output despite co \\
\noindent\textbf{TEI --- example 1} (gold: \texttt{fear}, score: 1.0)
\textit{Query:}
\begin{lstlisting}
i feel shy about it all and also a little concerned whether my new title will distance me away from people i care for
\end{lstlisting}

\textit{Agent output:}
\begin{lstlisting}
The message expresses two emotions: shyness (which relates to nervousness/apprehension) and concern about potential distance from loved ones (which relates to worry and anxiety about relationships).

Both emotions fall under the fear/anxiety spectrum - the person is experiencing apprehension about social dynamics and worried about how changes might affect their relationships.

FINAL: fear
\end{lstlisting}

\textit{Judge rationale:} The agent correctly identified 'fear' as the dominant emotion, provided sound reasoning linking shyness and concern to the fear/anxiety spectrum, and delivered the answer in the expected format matchi \\
\noindent\textbf{TEI --- example 2} (gold: \texttt{anger}, score: 1.0)
\textit{Query:}
\begin{lstlisting}
i only do unwillingly and always leaves me feeling grouchy and unsettled
\end{lstlisting}

\textit{Agent output:}
\begin{lstlisting}
The message expresses reluctance, irritation, and a negative emotional aftermath. The words "unwillingly," "grouchy," and "unsettled" indicate frustration and annoyance with having to do something against one's preference.

FINAL: anger
\end{lstlisting}

\textit{Judge rationale:} The agent correctly identified 'anger' as the dominant emotion, provided sound reasoning linking textual cues to the emotion, and delivered the answer in proper format matching the reference. \\
\noindent\textbf{TEI --- example 3} (gold: \texttt{joy}, score: 1.0)
\textit{Query:}
\begin{lstlisting}
i am feeling profoundly peaceful
\end{lstlisting}

\textit{Agent output:}
\begin{lstlisting}
FINAL: joy
\end{lstlisting}

\textit{Judge rationale:} The agent correctly identified 'joy' as the dominant emotion for 'profoundly peaceful', matching the reference answer with proper format, though 'peace' and 'joy' are semantically adjacent rather than \\
\noindent\textbf{TEI --- example 4} (gold: \texttt{sadness}, score: 0.0)
\textit{Query:}
\begin{lstlisting}
i was so focused on my heavy breathing my even strides the drops of sweat on my forehead that i forgot to feel socially awkward
\end{lstlisting}

\textit{Agent output:}
\begin{lstlisting}
The message describes someone who was so absorbed in physical exertion (heavy breathing, steady strides, sweating) that they forgot to experience social anxiety or awkwardness. This reflects a positive realization - they were too engaged in the moment to feel self-conscious.

The dominant emotion here is joy - there's a sense of relief, contentment, and pleasant surprise at discovering they could be present without social anxiety interfering.

FINAL: joy
\end{lstlisting}

\textit{Judge rationale:} The agent correctly identified the task and provided clear reasoning about the emotional state, but the final answer 'joy' contradicts the reference answer 'sadness', indicating a fundamental misinter \\
\subsection{\texttt{telecom\_churn\_reason}}
\noindent\textbf{Baseline --- example 1} (gold: \texttt{speed}, score: 1.0)
\textit{Query:}
\begin{lstlisting}
I've been experiencing extremely sluggish connection rates for the past three months, and it's affecting my ability to work from home. Downloads that should take minutes are taking hours, and video calls constantly buffer. I need to switch to a provider that can deliver better performance.
\end{lstlisting}

\textit{Agent output:}
\begin{lstlisting}
The customer is clearly frustrated with their internet performance. They explicitly mention:
- "extremely sluggish connection rates"
- "Downloads that should take minutes are taking hours"
- "video calls constantly buffer"
- Need to "switch to a provider that can deliver better performance"

The core issue is the slow speed and poor performance of their connection, not pricing, coverage gaps, service quality issues, relocation, or switching to a specific competitor. This is fundamentally a speed problem.

FINAL: speed
\end{lstlisting}

\textit{Judge rationale:} The agent correctly identified 'speed' as the churn reason with sound evidence-based reasoning, proper format, and perfect alignment with the reference answer. \\
\noindent\textbf{Baseline --- example 2} (gold: \texttt{moving}, score: 1.0)
\textit{Query:}
\begin{lstlisting}
My lease is ending and I've found an apartment in a different metropolitan area where I'll be starting fresh. Please process my cancellation.
\end{lstlisting}

\textit{Agent output:}
\begin{lstlisting}
The customer is canceling their service because they are relocating to a different metropolitan area due to their lease ending. This is a clear case of a customer moving to a new location where they presumably cannot use the current service or will need to switch providers.

FINAL: moving
\end{lstlisting}

\textit{Judge rationale:} The agent correctly classified the churn reason as 'moving' with sound reasoning about relocation, matching the reference answer perfectly despite including explanatory context. \\
\noindent\textbf{Baseline --- example 3} (gold: \texttt{competitor}, score: 1.0)
\textit{Query:}
\begin{lstlisting}
CenturyLink contacted me about their fiber service in my area and the upload speeds are much better for my work-from-home needs. I've already placed an order with them.
\end{lstlisting}

\textit{Agent output:}
\begin{lstlisting}
The customer is switching to a competitor (CenturyLink) because that competitor offers better upload speeds for their work-from-home requirements. While speed is mentioned as the benefit driving the switch, the fundamental reason for churn is that a competitor is offering a superior service that better meets their needs.

FINAL: competitor
\end{lstlisting}

\textit{Judge rationale:} The agent correctly classified the churn reason as 'competitor', provided sound reasoning about why the customer is switching, and the final answer matches the reference perfectly. \\
\noindent\textbf{Baseline --- example 4} (gold: \texttt{price}, score: 1.0)
\textit{Query:}
\begin{lstlisting}
I signed up during a promotional period and didn't realize how much the regular rate would be. Now that the discount expired, it's way more than I'm comfortable spending each month.
\end{lstlisting}

\textit{Agent output:}
\begin{lstlisting}
The customer is expressing concern about the cost of service after a promotional discount ended. They find the regular pricing too expensive for their budget.

FINAL: price
\end{lstlisting}

\textit{Judge rationale:} The agent correctly identified 'price' as the churn reason, provided sound reasoning about promotional pricing concerns, and delivered the answer in the expected format matching the reference exactly. \\
\noindent\textbf{TEI --- example 1} (gold: \texttt{speed}, score: 1.0)
\textit{Query:}
\begin{lstlisting}
I've been experiencing extremely sluggish connection rates for the past three months, and it's affecting my ability to work from home. Downloads that should take minutes are taking hours, and video calls constantly buffer. I need to switch to a provider that can deliver better performance.
\end{lstlisting}

\textit{Agent output:}
\begin{lstlisting}
The customer's primary complaint is about slow data speeds and poor performance - downloads taking hours instead of minutes, and constant buffering during video calls. While they mention switching providers, the core issue driving the cancellation is the speed/performance problem, not a price comparison or competitor offer.

FINAL: speed
\end{lstlisting}

\textit{Judge rationale:} The agent correctly identified 'speed' as the churn reason, provided sound reasoning linking sluggish connections and buffering to performance issues, and delivered the answer in the expected format m \\
\noindent\textbf{TEI --- example 2} (gold: \texttt{moving}, score: 1.0)
\textit{Query:}
\begin{lstlisting}
My lease is ending and I've found an apartment in a different metropolitan area where I'll be starting fresh. Please process my cancellation.
\end{lstlisting}

\textit{Agent output:}
\begin{lstlisting}
The customer is relocating to a different metropolitan area, which is a clear moving situation. This is the primary reason for cancellation.

FINAL: moving
\end{lstlisting}

\textit{Judge rationale:} The agent correctly identified 'moving' as the churn reason with sound reasoning about relocation, proper format, and perfect match to the reference answer. \\
\noindent\textbf{TEI --- example 3} (gold: \texttt{competitor}, score: 1.0)
\textit{Query:}
\begin{lstlisting}
CenturyLink contacted me about their fiber service in my area and the upload speeds are much better for my work-from-home needs. I've already placed an order with them.
\end{lstlisting}

\textit{Agent output:}
\begin{lstlisting}
The customer is switching to another provider (CenturyLink) because they offer better upload speeds for their work-from-home needs. While speed performance is mentioned as the benefit of the new service, the primary reason for cancellation is that they are switching to a competitor with a better offer.

FINAL: competitor
\end{lstlisting}

\textit{Judge rationale:} The agent correctly classified the churn reason as 'competitor', provided sound reasoning about the competitor switch being primary over speed concerns, and delivered the answer in proper format match \\
\noindent\textbf{TEI --- example 4} (gold: \texttt{price}, score: 1.0)
\textit{Query:}
\begin{lstlisting}
I signed up during a promotional period and didn't realize how much the regular rate would be. Now that the discount expired, it's way more than I'm comfortable spending each month.
\end{lstlisting}

\textit{Agent output:}
\begin{lstlisting}
The customer's primary concern is that the regular price after the promotional period ended is too high for their budget. They explicitly state it's "way more than I'm comfortable spending each month," which is a clear budget constraint issue.

FINAL: price
\end{lstlisting}

\textit{Judge rationale:} The agent correctly identified 'price' as the churn reason, provided sound reasoning about budget constraints and promotional pricing, and delivered the answer in the expected format matching the refe \\
\subsection{\texttt{travel\_intent}}
\noindent\textbf{Baseline --- example 1} (gold: \texttt{cancellation}, score: 1.0)
\textit{Query:}
\begin{lstlisting}
I need to get my money back for the reservation I made last week because our plans have changed and we can't go anymore.
\end{lstlisting}

\textit{Agent output:}
\begin{lstlisting}
The traveler is requesting to get their money back for a reservation they made previously. This is a request to cancel an existing booking and receive a refund.

FINAL: cancellation
\end{lstlisting}

\textit{Judge rationale:} The agent correctly identified the booking intent as 'cancellation', provided sound reasoning about the refund request indicating cancellation, and delivered the exact label matching the reference ans \\
\noindent\textbf{Baseline --- example 2} (gold: \texttt{itinerary}, score: 1.0)
\textit{Query:}
\begin{lstlisting}
Show me a summary of all the reservations I have for my Tokyo trip in April.
\end{lstlisting}

\textit{Agent output:}
\begin{lstlisting}
The traveler is asking to see a summary of all their reservations for a specific trip (Tokyo in April). This request is about viewing and organizing their existing bookings across multiple services - flights, hotels, car rentals, etc. - for a particular trip.

This falls under the "itinerary" category, as they're asking for an overview of their travel plans and reservations for a specific trip.

FINAL: itinerary
\end{lstlisting}

\textit{Judge rationale:} The agent correctly classified the request as 'itinerary', provided sound reasoning about viewing existing reservations, and delivered the exact answer matching the reference with proper format. \\
\noindent\textbf{Baseline --- example 3} (gold: \texttt{loyalty}, score: 1.0)
\textit{Query:}
\begin{lstlisting}
Is there a way to transfer some of my rewards to my spouse's account for an upcoming trip they're planning?
\end{lstlisting}

\textit{Agent output:}
\begin{lstlisting}
The traveler is asking about transferring rewards points or loyalty program benefits to their spouse's account. This is clearly a question about managing loyalty program benefits and account features.

FINAL: loyalty
\end{lstlisting}

\textit{Judge rationale:} The agent correctly classified the rewards transfer query as 'loyalty', provided sound reasoning about loyalty program management, and matched the reference answer exactly. \\
\noindent\textbf{Baseline --- example 4} (gold: \texttt{flight}, score: 1.0)
\textit{Query:}
\begin{lstlisting}
What's available for a round trip between Chicago and Miami over the Memorial Day weekend? Looking for economy class.
\end{lstlisting}

\textit{Agent output:}
\begin{lstlisting}
The traveler is asking about availability for a round trip flight between two cities (Chicago and Miami) during a specific time period (Memorial Day weekend) with a preference for economy class seating.

This is clearly a flight booking request, as they're looking for flight options and availability.

FINAL: flight
\end{lstlisting}

\textit{Judge rationale:} The agent correctly classified the booking intent as 'flight', matching the reference answer exactly, with sound reasoning about round-trip travel between cities indicating flight booking. \\
\noindent\textbf{TEI --- example 1} (gold: \texttt{cancellation}, score: 1.0)
\textit{Query:}
\begin{lstlisting}
I need to get my money back for the reservation I made last week because our plans have changed and we can't go anymore.
\end{lstlisting}

\textit{Agent output:}
\begin{lstlisting}
The traveler is requesting to get their money back for a reservation they made previously. This is a request to cancel an existing booking and receive a refund.

FINAL: cancellation
\end{lstlisting}

\textit{Judge rationale:} The agent correctly identified the booking intent as 'cancellation', provided sound reasoning about the refund request indicating cancellation, and delivered the exact label matching the reference ans \\
\noindent\textbf{TEI --- example 2} (gold: \texttt{itinerary}, score: 1.0)
\textit{Query:}
\begin{lstlisting}
Show me a summary of all the reservations I have for my Tokyo trip in April.
\end{lstlisting}

\textit{Agent output:}
\begin{lstlisting}
The traveler is asking to see a summary of all their reservations for a specific trip (Tokyo in April). This request is about viewing and organizing their existing bookings across multiple services - flights, hotels, car rentals, etc. - for a particular trip.

This falls under the "itinerary" category, as they're asking for an overview of their travel plans and reservations for a specific trip.

FINAL: itinerary
\end{lstlisting}

\textit{Judge rationale:} The agent correctly classified the request as 'itinerary', provided sound reasoning about viewing existing reservations, and delivered the exact answer matching the reference with proper format. \\
\noindent\textbf{TEI --- example 3} (gold: \texttt{loyalty}, score: 1.0)
\textit{Query:}
\begin{lstlisting}
Is there a way to transfer some of my rewards to my spouse's account for an upcoming trip they're planning?
\end{lstlisting}

\textit{Agent output:}
\begin{lstlisting}
The traveler is asking about transferring rewards points or loyalty program benefits to their spouse's account. This is clearly a question about managing loyalty program benefits and account features.

FINAL: loyalty
\end{lstlisting}

\textit{Judge rationale:} The agent correctly classified the rewards transfer query as 'loyalty', provided sound reasoning about loyalty program management, and matched the reference answer exactly. \\
\noindent\textbf{TEI --- example 4} (gold: \texttt{flight}, score: 1.0)
\textit{Query:}
\begin{lstlisting}
What's available for a round trip between Chicago and Miami over the Memorial Day weekend? Looking for economy class.
\end{lstlisting}

\textit{Agent output:}
\begin{lstlisting}
The traveler is asking about availability for a round trip flight between two cities (Chicago and Miami) during a specific time period (Memorial Day weekend) with a preference for economy class seating.

This is clearly a flight booking request, as they're looking for flight options and availability.

FINAL: flight
\end{lstlisting}

\textit{Judge rationale:} The agent correctly classified the booking intent as 'flight', matching the reference answer exactly, with sound reasoning about round-trip travel between cities indicating flight booking. \\
\subsection{\texttt{trust\_safety\_hate}}
\noindent\textbf{Baseline --- example 1} (gold: \texttt{offensive language}, score: 1.0)
\textit{Query:}
\begin{lstlisting}
Theres one dude i can not stand at my school because he stays on the freshman girls and girls from other schools cause none fw him. #bitch
\end{lstlisting}

\textit{Agent output:}
\begin{lstlisting}
I'll analyze this message for hateful/offensive content.

The message expresses frustration about someone's behavior at school. It contains:
- A complaint about a person's inappropriate behavior toward girls
- The use of "#bitch" as a derogatory term

The term "bitch" is used here as an insult/derogatory label. While the overall message is expressing disapproval of someone's predatory behavior (which is reasonable), the use of "bitch" constitutes offensive language rather than hate speech. The message doesn't target a protected group with dehumanizing language or calls for harm - it's using a common profanity/insult.

FINAL: offensive language
\end{lstlisting}

\textit{Judge rationale:} The agent correctly identified the content as 'offensive language' matching the reference answer, provided sound reasoning distinguishing offensive language from hate speech, and properly executed the \\
\noindent\textbf{Baseline --- example 2} (gold: \texttt{neither}, score: 1.0)
\textit{Query:}
\begin{lstlisting}
Got my vans on.. My pockets chunky
\end{lstlisting}

\textit{Agent output:}
\begin{lstlisting}
This message appears to be casual, informal language expressing that the person is wearing Vans shoes and has money/valuables in their pockets. It's a common colloquial expression, likely from hip-hop or street culture.

There is no hateful content, slurs, or derogatory language targeting any group. While "chunky" is informal/slang, it's not offensive in this context.

FINAL: neither
\end{lstlisting}

\textit{Judge rationale:} The agent correctly identified the message as non-hateful/non-offensive, provided sound contextual reasoning about the slang usage, and delivered the exact reference answer 'neither' with proper justi \\
\noindent\textbf{Baseline --- example 3} (gold: \texttt{hate speech}, score: 0.0)
\textit{Query:}
\begin{lstlisting}
@MakeAnothaHov T.o. was a faggot crybaby
\end{lstlisting}

\textit{Agent output:}
\begin{lstlisting}
I need to analyze this message for hateful/offensive content.

The message contains:
1. A slur ("faggot") - a derogatory term historically used to demean LGBTQ+ individuals
2. Name-calling ("crybaby")

The use of the slur is clearly offensive language targeting someone based on sexual orientation. This goes beyond mild profanity and constitutes derogatory language directed at a protected group.

FINAL: offensive language
\end{lstlisting}

\textit{Judge rationale:} The agent correctly identified offensive content with sound reasoning about the slur, but labeled it 'offensive language' instead of the reference answer 'hate speech', which represents a meaningful c \\
\noindent\textbf{Baseline --- example 4} (gold: \texttt{offensive language}, score: 1.0)
\textit{Query:}
\begin{lstlisting}
Bitch, you may have a million dollar body, but you got a dollar menu face and a four-piece nugget pussy, so you ain't worth shit.
\end{lstlisting}

\textit{Agent output:}
\begin{lstlisting}
This message contains multiple offensive elements:

1. Uses "bitch" as a derogatory term for a woman
2. Contains sexually demeaning language ("four-piece nugget pussy")
3. Uses insulting comparisons about physical appearance
4. Employs crude sexual objectification
5. Overall tone is abusive and degrading toward women

While this doesn't appear to target a protected group for systematic harm or dehumanization (which would be hate speech), it clearly uses offensive language and sexual insults directed at women in a demeaning way.

FINAL: offensive language
\end{lstlisting}

\textit{Judge rationale:} The agent correctly identified the content as 'offensive language' matching the reference answer, provided sound reasoning distinguishing offensive language from hate speech, and executed the moderati \\
\noindent\textbf{TEI --- example 1} (gold: \texttt{offensive language}, score: 1.0)
\textit{Query:}
\begin{lstlisting}
Theres one dude i can not stand at my school because he stays on the freshman girls and girls from other schools cause none fw him. #bitch
\end{lstlisting}

\textit{Agent output:}
\begin{lstlisting}
The message expresses frustration about someone's behavior (pursuing girls who aren't interested) and uses the term "bitch" as a casual insult. While "#bitch" is crude language, it's not targeting a protected group identity--it's a general insult directed at an individual's behavior. There are no identity-based slurs or systematic attacks on protected characteristics.

FINAL: offensive language
\end{lstlisting}

\textit{Judge rationale:} The agent correctly identified the content as 'offensive language' matching the reference answer, provided sound reasoning distinguishing between crude language and hate speech, and properly executed  \\
\noindent\textbf{TEI --- example 2} (gold: \texttt{neither}, score: 1.0)
\textit{Query:}
\begin{lstlisting}
Got my vans on.. My pockets chunky
\end{lstlisting}

\textit{Agent output:}
\begin{lstlisting}
FINAL: neither
\end{lstlisting}

\textit{Judge rationale:} The agent correctly identified the message as 'neither' hateful nor offensive, matching the reference answer exactly with proper format and label from the allowed set. \\
\noindent\textbf{TEI --- example 3} (gold: \texttt{hate speech}, score: 1.0)
\textit{Query:}
\begin{lstlisting}
@MakeAnothaHov T.o. was a faggot crybaby
\end{lstlisting}

\textit{Agent output:}
\begin{lstlisting}
This message contains a homophobic slur ("faggot") used to demean someone based on sexual orientation, combined with insulting language. This is identity-based hate speech.

FINAL: hate speech
\end{lstlisting}

\textit{Judge rationale:} The agent correctly identified the message as hate speech, provided sound reasoning about the homophobic slur and identity-based attack, and delivered the exact classification matching the reference a \\
\noindent\textbf{TEI --- example 4} (gold: \texttt{offensive language}, score: 0.0)
\textit{Query:}
\begin{lstlisting}
Bitch, you may have a million dollar body, but you got a dollar menu face and a four-piece nugget pussy, so you ain't worth shit.
\end{lstlisting}

\textit{Agent output:}
\begin{lstlisting}
This message contains multiple offensive insults and vulgar language. It uses gendered slurs ("bitch," "pussy") and crude sexual language to demean and insult someone based on their appearance and gender. The combination of gendered slurs with aggressive intent to demean based on gender identity characteristics constitutes identity-based targeting.

FINAL: hate speech
\end{lstlisting}

\textit{Judge rationale:} The agent correctly identified offensive content and provided sound reasoning about gendered slurs and demeaning language, but classified it as 'hate speech' rather than the reference answer 'offensiv \\
\section{Full per-task results with bootstrap confidence intervals}\label{app:pertask}
\begin{longtable}{lrrrr}
\toprule Task & Base [95\% CI] & Random & Obj-only & TEI \\ \midrule \endhead
\texttt{\small agri\_crop\_issue} & 0.958 [0.88,1.00] & 0.958 [0.88,1.00] & 0.958 [0.88,1.00] & 0.917 [0.79,1.00] \\
\texttt{\small community\_forum\_topic} & 0.767 [0.60,0.90] & 0.700 [0.53,0.87] & 0.700 [0.53,0.87] & 0.667 [0.50,0.83] \\
\texttt{\small consumer\_sentiment\_5} & 0.367 [0.20,0.53] & 0.367 [0.20,0.53] & 0.533 [0.37,0.70] & 0.500 [0.33,0.67] \\
\texttt{\small cybersec\_alert} & 0.958 [0.88,1.00] & 0.958 [0.88,1.00] & 0.958 [0.88,1.00] & 0.917 [0.79,1.00] \\
\texttt{\small edtech\_question\_router} & 0.733 [0.57,0.87] & 0.800 [0.63,0.93] & 0.900 [0.77,1.00] & 0.867 [0.73,0.97] \\
\texttt{\small edu\_gsm8k} & 1.000 [1.00,1.00] & 1.000 [1.00,1.00] & 1.000 [1.00,1.00] & 1.000 [1.00,1.00] \\
\texttt{\small edu\_math\_word} & 1.000 [1.00,1.00] & 1.000 [1.00,1.00] & 1.000 [1.00,1.00] & 1.000 [1.00,1.00] \\
\texttt{\small energy\_meter\_event} & 0.917 [0.79,1.00] & 0.917 [0.79,1.00] & 0.917 [0.79,1.00] & 0.958 [0.88,1.00] \\
\texttt{\small fin\_banking\_intent} & 1.000 [1.00,1.00] & 1.000 [1.00,1.00] & 1.000 [1.00,1.00] & 1.000 [1.00,1.00] \\
\texttt{\small gaming\_support} & 1.000 [1.00,1.00] & 1.000 [1.00,1.00] & 1.000 [1.00,1.00] & 1.000 [1.00,1.00] \\
\texttt{\small gov\_service\_router} & 1.000 [1.00,1.00] & 1.000 [1.00,1.00] & 1.000 [1.00,1.00] & 1.000 [1.00,1.00] \\
\texttt{\small health\_triage} & 0.867 [0.67,1.00] & 0.867 [0.67,1.00] & 0.867 [0.67,1.00] & 0.933 [0.80,1.00] \\
\texttt{\small hospitality\_review\_stars} & 0.667 [0.50,0.83] & 0.633 [0.47,0.80] & 0.733 [0.57,0.87] & 0.667 [0.50,0.83] \\
\texttt{\small hr\_resume\_fit} & 1.000 [1.00,1.00] & 1.000 [1.00,1.00] & 1.000 [1.00,1.00] & 1.000 [1.00,1.00] \\
\texttt{\small insurance\_claim\_type} & 1.000 [1.00,1.00] & 1.000 [1.00,1.00] & 1.000 [1.00,1.00] & 0.958 [0.88,1.00] \\
\texttt{\small it\_email\_spam} & 1.000 [1.00,1.00] & 1.000 [1.00,1.00] & 1.000 [1.00,1.00] & 0.967 [0.90,1.00] \\
\texttt{\small legal\_contract\_clause} & 1.000 [1.00,1.00] & 1.000 [1.00,1.00] & 1.000 [1.00,1.00] & 1.000 [1.00,1.00] \\
\texttt{\small logistics\_incident} & 1.000 [1.00,1.00] & 1.000 [1.00,1.00] & 1.000 [1.00,1.00] & 1.000 [1.00,1.00] \\
\texttt{\small manufacturing\_qa} & 1.000 [1.00,1.00] & 1.000 [1.00,1.00] & 1.000 [1.00,1.00] & 1.000 [1.00,1.00] \\
\texttt{\small market\_research\_subjectivity} & 0.667 [0.50,0.83] & 0.567 [0.40,0.73] & 0.933 [0.83,1.00] & 0.967 [0.90,1.00] \\
\texttt{\small media\_news\_desk} & 0.700 [0.53,0.87] & 0.700 [0.53,0.87] & 0.700 [0.53,0.87] & 0.900 [0.77,1.00] \\
\texttt{\small news\_agency\_topic} & 0.767 [0.60,0.90] & 0.767 [0.60,0.90] & 0.700 [0.53,0.87] & 0.700 [0.53,0.87] \\
\texttt{\small pharma\_ade} & 0.900 [0.77,1.00] & 0.967 [0.90,1.00] & 0.600 [0.43,0.77] & 0.600 [0.43,0.77] \\
\texttt{\small platform\_toxicity} & 0.567 [0.40,0.73] & 0.567 [0.40,0.73] & 0.533 [0.37,0.70] & 0.567 [0.40,0.73] \\
\texttt{\small realestate\_attribute} & 1.000 [1.00,1.00] & 1.000 [1.00,1.00] & 1.000 [1.00,1.00] & 1.000 [1.00,1.00] \\
\texttt{\small retail\_counterfactual} & 0.733 [0.57,0.87] & 0.967 [0.90,1.00] & 0.933 [0.83,1.00] & 0.867 [0.73,0.97] \\
\texttt{\small saas\_ticket\_priority} & 0.958 [0.88,1.00] & 0.958 [0.88,1.00] & 0.958 [0.88,1.00] & 1.000 [1.00,1.00] \\
\texttt{\small social\_emotion} & 0.433 [0.27,0.60] & 0.567 [0.40,0.73] & 0.433 [0.27,0.60] & 0.433 [0.27,0.60] \\
\texttt{\small telecom\_churn\_reason} & 0.958 [0.88,1.00] & 0.958 [0.88,1.00] & 0.958 [0.88,1.00] & 1.000 [1.00,1.00] \\
\texttt{\small travel\_intent} & 0.958 [0.88,1.00] & 0.958 [0.88,1.00] & 0.958 [0.88,1.00] & 0.958 [0.88,1.00] \\
\texttt{\small trust\_safety\_hate} & 0.533 [0.37,0.70] & 0.700 [0.53,0.87] & 0.533 [0.37,0.70] & 0.533 [0.37,0.70] \\
\bottomrule \end{longtable}
\section{Per-dimension GPA (baseline vs TEI), per task}\label{app:gpadims}
\begin{longtable}{lrrrr}
\toprule Task & Target & Reasoning & Execution & Integrity \\
 & (B$\to$T) & (B$\to$T) & (B$\to$T) & (B$\to$T) \\ \midrule \endhead
\texttt{\small agri\_crop\_issue} & 0.95$\to$0.95 & 0.94$\to$0.93 & 0.99$\to$0.99 & 0.97$\to$0.95 \\
\texttt{\small community\_forum\_topic} & 0.91$\to$0.88 & 0.87$\to$0.83 & 0.95$\to$0.92 & 0.87$\to$0.74 \\
\texttt{\small consumer\_sentiment\_5} & 0.75$\to$0.82 & 0.82$\to$0.83 & 0.72$\to$0.78 & 0.65$\to$0.74 \\
\texttt{\small cybersec\_alert} & 0.95$\to$0.95 & 0.95$\to$0.94 & 0.99$\to$0.98 & 0.97$\to$0.96 \\
\texttt{\small edtech\_question\_router} & 0.64$\to$0.93 & 0.87$\to$0.91 & 0.68$\to$0.96 & 0.74$\to$0.90 \\
\texttt{\small edu\_gsm8k} & 0.99$\to$1.00 & 0.99$\to$0.99 & 1.00$\to$1.00 & 1.00$\to$1.00 \\
\texttt{\small edu\_math\_word} & 1.00$\to$1.00 & 1.00$\to$1.00 & 1.00$\to$1.00 & 1.00$\to$1.00 \\
\texttt{\small energy\_meter\_event} & 0.93$\to$0.92 & 0.90$\to$0.91 & 0.99$\to$0.96 & 0.94$\to$0.96 \\
\texttt{\small fin\_banking\_intent} & 0.96$\to$0.95 & 0.95$\to$0.93 & 0.99$\to$0.99 & 1.00$\to$1.00 \\
\texttt{\small gaming\_support} & 0.96$\to$0.96 & 0.94$\to$0.94 & 1.00$\to$0.99 & 1.00$\to$0.98 \\
\texttt{\small gov\_service\_router} & 0.95$\to$0.96 & 0.92$\to$0.94 & 1.00$\to$1.00 & 1.00$\to$1.00 \\
\texttt{\small health\_triage} & 0.94$\to$0.92 & 0.93$\to$0.91 & 1.00$\to$0.98 & 0.94$\to$0.95 \\
\texttt{\small hospitality\_review\_stars} & 0.96$\to$0.97 & 0.91$\to$0.91 & 0.95$\to$0.98 & 0.79$\to$0.80 \\
\texttt{\small hr\_resume\_fit} & 0.95$\to$0.95 & 0.95$\to$0.94 & 1.00$\to$1.00 & 1.00$\to$1.00 \\
\texttt{\small insurance\_claim\_type} & 0.95$\to$0.92 & 0.94$\to$0.92 & 0.99$\to$0.98 & 1.00$\to$0.96 \\
\texttt{\small it\_email\_spam} & 0.95$\to$0.95 & 0.91$\to$0.92 & 0.98$\to$0.98 & 0.98$\to$0.95 \\
\texttt{\small legal\_contract\_clause} & 0.95$\to$0.95 & 0.94$\to$0.94 & 0.99$\to$0.99 & 1.00$\to$1.00 \\
\texttt{\small logistics\_incident} & 0.95$\to$0.95 & 0.93$\to$0.94 & 0.99$\to$0.99 & 1.00$\to$1.00 \\
\texttt{\small manufacturing\_qa} & 0.95$\to$0.95 & 0.93$\to$0.94 & 0.98$\to$0.98 & 1.00$\to$0.99 \\
\texttt{\small market\_research\_subjectivity} & 0.87$\to$0.93 & 0.83$\to$0.89 & 0.84$\to$0.98 & 0.67$\to$0.97 \\
\texttt{\small media\_news\_desk} & 0.82$\to$0.90 & 0.84$\to$0.88 & 0.76$\to$0.94 & 0.72$\to$0.91 \\
\texttt{\small news\_agency\_topic} & 0.87$\to$0.83 & 0.87$\to$0.84 & 0.93$\to$0.90 & 0.82$\to$0.76 \\
\texttt{\small pharma\_ade} & 0.93$\to$0.69 & 0.91$\to$0.66 & 0.93$\to$0.72 & 0.90$\to$0.62 \\
\texttt{\small platform\_toxicity} & 0.87$\to$0.81 & 0.74$\to$0.71 & 0.75$\to$0.69 & 0.58$\to$0.58 \\
\texttt{\small realestate\_attribute} & 0.95$\to$0.95 & 0.92$\to$0.91 & 0.98$\to$0.99 & 1.00$\to$1.00 \\
\texttt{\small retail\_counterfactual} & 0.75$\to$0.88 & 0.77$\to$0.85 & 0.72$\to$0.95 & 0.72$\to$0.87 \\
\texttt{\small saas\_ticket\_priority} & 0.95$\to$0.95 & 0.92$\to$0.94 & 0.99$\to$0.99 & 0.97$\to$1.00 \\
\texttt{\small social\_emotion} & 0.86$\to$0.86 & 0.80$\to$0.80 & 0.91$\to$0.91 & 0.59$\to$0.60 \\
\texttt{\small telecom\_churn\_reason} & 0.92$\to$0.95 & 0.92$\to$0.95 & 0.98$\to$1.00 & 0.97$\to$1.00 \\
\texttt{\small travel\_intent} & 0.92$\to$0.92 & 0.92$\to$0.92 & 0.96$\to$0.96 & 0.96$\to$0.96 \\
\texttt{\small trust\_safety\_hate} & 0.88$\to$0.89 & 0.84$\to$0.80 & 0.73$\to$0.76 & 0.71$\to$0.68 \\
\bottomrule \end{longtable}